\documentclass[documento]{gpt}

% Livelli di titolo: i sottocasi d'uso usano \subsubsubsection
\newcommand{\subsubsubsection}[1]{\paragraph{#1}}
\setcounter{secnumdepth}{4}
\setcounter{tocdepth}{3}
\usepackage{titlesec}
\titleformat{\paragraph}{\normalfont\normalsize\bfseries\sffamily}{\theparagraph}{1em}{}
\titlespacing*{\paragraph}{0pt}{3.25ex plus 1ex minus .2ex}{1.2ex plus .2ex}

\usepackage{adjustbox}
\usetikzlibrary{arrows.meta, backgrounds, calc, fit, positioning,
                shapes.geometric, shapes.misc}

% ---------------------------------------------------------------------
% Stili dei diagrammi UML dei casi d'uso
% ---------------------------------------------------------------------
\tikzset{
  uc/.style={ellipse, draw=gptinchiostro, line width=0.6pt, fill=gptaccento!6,
    align=center, font=\sffamily\footnotesize, inner sep=0.5pt,
    minimum width=3.3cm, minimum height=1.25cm, text width=2.7cm,
    execute at begin node={\hyphenpenalty=10000\exhyphenpenalty=10000}},
  ucfocus/.style={uc, line width=1.1pt, fill=gptaccento!16, text width=3.0cm},
  dipendenza/.style={draw=gptinchiostro, dashed, line width=0.6pt,
    -{Stealth[open, length=2.6mm, width=2.2mm]}},
  generalizzazione/.style={draw=gptinchiostro, line width=0.6pt,
    -{Triangle[open, length=3.2mm, width=3.2mm]}},
  albero/.style={draw=gptinchiostro, line width=0.6pt},
  associazione/.style={draw=gptinchiostro, line width=0.6pt},
  etich/.style={font=\sffamily\scriptsize, fill=white, inner sep=1.2pt},
  nota/.style={draw=gptgrigio, fill=white, line width=0.5pt, align=left,
    font=\sffamily\scriptsize, text width=2.75cm, inner sep=3.5pt,
    chamfered rectangle, chamfered rectangle corners=north east,
    chamfered rectangle xsep=2.2mm,
    execute at begin node={\hyphenpenalty=10000\exhyphenpenalty=10000}},
  ancora/.style={draw=gptgrigio, densely dotted, line width=0.6pt},
  confine/.style={draw=gptinchiostro, line width=0.8pt, rounded corners=2pt,
    inner sep=0.45cm},
  nomesistema/.style={anchor=north west, font=\sffamily\footnotesize\bfseries,
    inner sep=4pt},
  attorebox/.style={minimum width=0.9cm, minimum height=1.9cm, inner sep=0pt},
}
% Attore UML: \attore{nome del nodo}{(x,y)}{etichetta}
\newcommand{\attore}[3]{%
  \node[attorebox] (#1) at #2 {};
  \begin{scope}[shift={(#1.center)}, line width=0.7pt, draw=gptinchiostro]
    \draw[fill=white] (0,0.55) circle (0.22);
    \draw (0,0.33) -- (0,-0.35);
    \draw (-0.38,0.12) -- (0.38,0.12);
    \draw (0,-0.35) -- (-0.3,-0.85) (0,-0.35) -- (0.3,-0.85);
  \end{scope}
  \node[font=\sffamily\footnotesize, align=center, anchor=north, inner sep=2pt]
    (#1-l) at (#1.south) {#3};}

% Scheda di un caso d'uso
\newenvironment{scheda}
  {\par\noindent\small
   \tabularx{\linewidth}{@{}>{\raggedright\arraybackslash\sffamily\bfseries\footnotesize}p{3.0cm}
                           >{\raggedright\arraybackslash}X@{}}\toprule}
  {\bottomrule\endtabularx\par\bigskip}
\setlist[itemize]{label={\color{gptaccento}\textbullet}}

% ---------------------------------------------------------------------
\titolo{Analisi dei Requisiti}
\uso{esterno}
\destinatari{Prof.~Tullio Vardanega\newline Prof.~Riccardo Cardin\newline Eggon S.r.l.}

% Registro delle modifiche: le righe sono nel file
%     analisi-dei-requisiti.registro.tex
% accanto a questo, e le aggiunge il bot dopo il merge in base
% all'etichetta della pull request (GUIDA.md, sezione 5).
% Qui non si scrive nulla.
%
% Questo file è generato da genera.py (cartella del generatore):
% modificate dati.py, requisiti.py o testi.py e rigenerate.

\begin{document}

% =====================================================================
\section{Introduzione}
% =====================================================================

\subsection{Scopo del documento}
Il documento descrive i \glo[Requisito]{requisiti} di \textit{NEXUM Concierge},
il prodotto richiesto da Eggon S.r.l.\ con il \glo{Capitolato} C2. A partire
dal capitolato e dagli incontri con il \glo{Proponente}, individua chi usa il
sistema e per quali scopi, ne descrive il comportamento atteso con i casi
d'uso e ne ricava requisiti verificabili, ciascuno collegato alla propria
fonte.

Il documento è il riferimento comune per tre attività successive:
\begin{itemize}
  \item la \textbf{progettazione}, che deve soddisfare ogni requisito
        obbligatorio;
  \item la \textbf{verifica}, che costruisce i \glo[Test di accettazione]{test di accettazione} a partire
        dai \glo[Caso d'uso]{casi d'uso} e dai criteri di accettazione del capitolato;
  \item il \textbf{confronto con il proponente}, che valida il perimetro e le
        priorità.
\end{itemize}
Non descrive scelte tecnologiche o architetturali di dettaglio, che sono
oggetto della \textit{Specifica Tecnica}; dove un vincolo tecnologico è stato
espresso dal proponente, è riportato come requisito di vincolo con la priorità
concordata.

\subsection{Scopo del prodotto}
NEXUM è la piattaforma di Eggon per la comunicazione aziendale e la gestione dei
servizi e dei processi \glo{HR}: un \glo{Backoffice} amministrativo e una \glo{PWA} per il
personale. \textit{NEXUM Concierge} introduce nella PWA un assistente
conversazionale che permette di trovare informazioni, capire le procedure e
usare i servizi aziendali scrivendo in linguaggio naturale. Il \glo{MVP} applica
questa visione all'area \textbf{ferie e permessi}; la sezione~\ref{sec:prodotto}
ne descrive obiettivi, capacità e perimetro.

\subsection{Glossario}
I termini tecnici, gli acronimi e le parole che richiedono una definizione
precisa sono raccolti nel \textit{Glossario}. Nel testo sono marcati con il
pedice \textsubscript{G} alla loro prima occorrenza significativa.

\subsection{Sviluppo del documento}
Il documento è redatto in modo incrementale: casi d'uso e requisiti vengono
aggiornati dopo ogni incontro con il proponente e a ogni revisione dei
documenti, e le modifiche sono riportate nel registro iniziale. I codici dei
\glo[Requisito]{requisiti} sono stabili: un requisito eliminato non cede il
proprio codice e i requisiti nuovi prendono i numeri successivi all'ultimo
assegnato, così che le osservazioni del proponente restino riferite allo stesso
requisito. I codici dei casi d'uso identificano invece un caso d'uso nella
versione corrente del documento. Le tabelle di tracciamento della
sezione~\ref{sec:tracciamento} permettono di controllare a ogni versione che ogni
fonte sia coperta e che ogni requisito abbia un'origine.

Questa versione recepisce i chiarimenti dell'incontro con Eggon del 2026-11-18
e il riscontro scritto del proponente sulla bozza dei requisiti. I cambiamenti
principali riguardano l'applicativo autonomo con un proprio accesso, le
conversazioni senza storico per il dipendente e la \glo[Tracciabilità]{tracciabilità} essenziale,
ora obbligatoria. Riguardano anche i bucket, consultabili da tutti i dipendenti
dell'azienda, la cancellazione delle richieste e il versionamento delle fonti.
Le questioni ancora da concordare sono raccolte nella
sezione~\ref{sec:aperte}.

\subsection{Riferimenti}

\subsubsection{Riferimenti normativi}
\begin{itemize}
  \item \textbf{Capitolato C2} -- \textit{[NEXUM] BRD preliminare -- Concierge
        per il dipendente}, Eggon S.r.l., versione 5 del 2026-09-18;
  \item \textbf{Regolamento del progetto didattico}, corso di Ingegneria del
        Software, a.a.~2026/2027;
  \item \textbf{Norme di Progetto}, Git Push Team, versione corrente.
\end{itemize}

\subsubsection{Riferimenti informativi}
\begin{itemize}
  \item \textbf{Verbale esterno del 2026-10-26} (VE-2026-10-26), incontro con
        Eggon su perimetro, architettura, dati e organizzazione;
  \item \textbf{Verbale esterno del 2026-11-18} (VE-2026-11-18), incontro con
        Eggon sulla prima lista dei requisiti;
  \item \textbf{Riscontro di Eggon sulla bozza dei requisiti} (novembre 2026),
        osservazioni scritte del proponente;
  \item \textbf{Glossario}, Git Push Team, versione corrente;
  \item slide del corso di Ingegneria del Software: \textit{Analisi dei
        requisiti} (T.~Vardanega) e \textit{Diagrammi dei casi d'uso}
        (R.~Cardin), a.a.~2026/2027;
  \item OMG, \textit{Unified Modeling Language}, versione 2.5.1,
        \url{https://www.omg.org/spec/UML/2.5.1};
  \item ISO/IEC/IEEE 29148:2018, \textit{Systems and software engineering --
        Life cycle processes -- Requirements engineering}.
\end{itemize}


% =====================================================================
\section{Descrizione del prodotto}\label{sec:prodotto}
% =====================================================================

\subsection{Contesto}
Una persona che vuole prendere ferie si pone domande di natura diversa.
«Come funzionano le ferie in azienda?» richiede di consultare regolamenti e
procedure; «Dove posso segnare le ferie?» richiede le istruzioni di NEXUM;
«Quante ferie ho fatto?» richiede dati personali aggiornati; «Registra le mie
ferie» richiede un'operazione applicativa. Oggi queste risposte stanno in
luoghi diversi: documenti aziendali, manuali, schermate della PWA, uffici HR.

Il Concierge riunisce questi percorsi in una conversazione: riconosce di quale
capacità ha bisogno ogni richiesta, usa la fonte o il servizio appropriato e
mantiene il contesto, così che la stessa conversazione possa passare da una
domanda sulle regole alla consultazione dei dati e all'invio di una richiesta.

\subsection{Obiettivi del prodotto}
Gli obiettivi di business indicati dal capitolato sono:
\begin{itemize}
  \item \textbf{semplificare l'accesso ai servizi}: la persona esprime la propria
        esigenza senza conoscere prima schermate e procedure;
  \item \textbf{rendere consultabili le regole aziendali}, con risposte riferite
        ai documenti dell'organizzazione;
  \item \textbf{facilitare l'uso di NEXUM} con istruzioni contestuali sulle sue
        funzioni;
  \item \textbf{migliorare l'accesso ai dati personali}, aggiornati e
        comprensibili, attraverso i servizi applicativi;
  \item \textbf{accompagnare le operazioni} con raccolta guidata dei dati e
        conferma, riducendo omissioni ed errori;
  \item \textbf{ridurre le richieste ripetitive} all'ufficio HR;
  \item \textbf{verificare una base estensibile}: quattro capacità riutilizzabili
        in altri ambiti di NEXUM.
\end{itemize}

\subsection{Le quattro capacità del Concierge}
Le capacità condividono interfaccia, contesto e controlli: non sono progetti
separati, ma modi diversi di rispondere a un messaggio
(Tabella~\ref{tab:capacita}).

\begin{table}[H]
  \centering
  \small
  \caption{Le quattro capacità del Concierge\label{tab:capacita}}
  \begin{tabularx}{\linewidth}{@{}>{\raggedright\arraybackslash}p{2.9cm}>{\raggedright\arraybackslash}X>{\raggedright\arraybackslash}p{3.6cm}l@{}}
    \toprule
    \intestazione{Capacità} & \intestazione{Esempio} & \intestazione{Fonte o strumento} & \intestazione{Casi d'uso} \\
    \midrule
    Conoscenza aziendale & «Come funziona la richiesta di ferie nella nostra azienda?» & regolamenti e procedure caricati dal backoffice & \hyperref[uc:regole]{UC11} \\ \addlinespace
    Conoscenza del prodotto & «Dove posso segnare le ferie?» & manuali e istruzioni NEXUM & \hyperref[uc:prodotto]{UC14} \\ \addlinespace
    Consultazione dei dati & «Quante ferie ho fatto quest'anno?» & \glo[Abilità]{abilità} di lettura dei dati personali & \hyperref[uc:dati]{UC15}, \hyperref[uc:stato_rich]{UC19} \\ \addlinespace
    Esecuzione di operazioni & «Registra che sarò in ferie dal 10 al 15 ottobre.» & abilità di registrazione della richiesta & \hyperref[uc:registra]{UC20} \\
    \bottomrule
  \end{tabularx}
\end{table}

\subsection{Architettura di riferimento}
La Figura~\ref{fig:architettura} riassume, secondo il capitolato, il percorso
di una richiesta. Il modello linguistico (\glo{LLM}) aiuta a interpretare la
richiesta e a formulare la risposta; la gestione del dialogo mantiene il
contesto e coordina le capacità; i \textit{gate} accedono alla conoscenza o ai
servizi applicativi, che applicano autorizzazioni e validazioni. I servizi
NEXUM sono raggiunti tramite le \glo{API} di Eggon, con chiamate autenticate
con \glo{OAuth}. La lettura dei dati non richiede conferma, l'invio di una richiesta
sì. Lo schema descrive responsabilità funzionali e non prescrive tecnologie.

\begin{figure}[H]
  \centering
  \begin{adjustbox}{max width=\linewidth}
  \begin{tikzpicture}[
      passo/.style={draw=gptaccento, fill=gptaccento!7, rounded corners=3pt,
        font=\sffamily\footnotesize, align=center, minimum height=0.8cm,
        text width=3.7cm, inner sep=3pt,
        execute at begin node={\hyphenpenalty=10000\exhyphenpenalty=10000}},
      evid/.style={passo, fill=gptaccento!18},
      conf/.style={passo, draw=orange!70!black, fill=orange!12},
      freccia/.style={-{Stealth[length=2mm]}, draw=gptaccento, line width=0.7pt},
      gruppo/.style={draw=gptlinea, rounded corners=4pt, inner sep=6pt},
      tit/.style={font=\sffamily\scriptsize\bfseries, text=gptgrigio, align=left,
        anchor=south west, inner sep=2pt}]
    \node[passo] (rich) at (0,0) {Richiesta nella chat};
    \node[passo, below=0.45cm of rich] (dial) {Gestione del dialogo e del contesto};
    \node[evid, right=1.2cm of dial] (llm) {Modello linguistico};
    \node[passo, below=0.45cm of dial] (sel) {Selezione della capacità};
    % conoscenza
    \node[evid, below left=1.7cm and 0.9cm of sel] (gk) {Gate conoscenza};
    \node[passo, below=0.4cm of gk] (ric) {Ricerca nelle fonti dell'azienda};
    \node[passo, below=0.4cm of ric] (evi) {Evidenze e fonti};
    % lettura
    \node[evid, below=1.7cm of sel] (gl) {Gate di lettura};
    \node[passo, below=0.4cm of gl] (snl) {API Eggon (OAuth):\\autorizzazione e lettura};
    \node[passo, below=0.4cm of snl] (dp) {Dati personali aggiornati};
    % invio
    \node[evid, below right=1.7cm and 0.9cm of sel] (go) {Gate operativo};
    \node[passo, below=0.4cm of go] (rr) {Raccolta dati e riepilogo};
    \node[conf, below=0.4cm of rr] (cp) {Conferma della persona};
    \node[passo, below=0.4cm of cp] (sno) {API Eggon (OAuth):\\autorizzazione e validazione};
    \node[passo, below=0.4cm of sno] (es) {Esito effettivo};
    \node[passo, below=4.1cm of dp] (comp) {Composizione della risposta};
    \node[passo, below=0.45cm of comp] (risp) {Risposta, fonti o stato nella chat};
    \begin{scope}[on background layer]
      \node[gruppo, fit=(gk)(ric)(evi), label={[tit]north west:CONOSCENZA}] {};
      \node[gruppo, fit=(gl)(snl)(dp), label={[tit]north west:LETTURA\\DEI DATI}] {};
      \node[gruppo, fit=(go)(rr)(cp)(sno)(es), label={[tit]north west:INVIO DI\\UNA RICHIESTA}] {};
    \end{scope}
    \draw[freccia] (rich) -- (dial);
    \draw[freccia, {Stealth[length=2mm]}-{Stealth[length=2mm]}] (dial) -- (llm);
    \draw[freccia] (dial) -- (sel);
    \draw[freccia] (sel.south) -- ++(0,-0.3) -| (gk.north);
    \draw[freccia] (sel) -- (gl);
    \draw[freccia] (sel.south) -- ++(0,-0.3) -| (go.north);
    \foreach \a/\b in {gk/ric, ric/evi, gl/snl, snl/dp, go/rr, rr/cp, cp/sno, sno/es}
      \draw[freccia] (\a) -- (\b);
    \draw[freccia] (evi.south) |- ($(comp.north)+(0,0.5)$) -- (comp.north);
    \draw[freccia] (dp.south) -- (comp.north);
    \draw[freccia] (es.south) |- ($(comp.north)+(0,0.5)$);
    \draw[freccia] (comp) -- (risp);
  \end{tikzpicture}
  \end{adjustbox}
  \caption{Percorso di una richiesta nel Concierge (rielaborato dal capitolato)}
  \label{fig:architettura}
\end{figure}

\subsection{Perimetro dell'MVP}
L'MVP deve essere completo nell'area ferie e permessi, senza realizzare l'intera
visione evolutiva. Comprende:
\begin{itemize}
  \item un applicativo autonomo, con un proprio accesso per gli utenti di prova,
        che comunica con NEXUM tramite le API di Eggon;
  \item la chat testuale con gestione del contesto della conversazione
        corrente;
  \item il caricamento e la gestione delle fonti aziendali nel backoffice,
        organizzate in \glo[Bucket]{bucket};
  \item l'uso dei manuali NEXUM come fonte distinta dai regolamenti aziendali,
        caricati dagli amministratori in bucket di conoscenza del prodotto;
  \item la costruzione e la consultazione della conoscenza, con citazione delle
        fonti e astensione quando non bastano;
  \item le abilità di lettura dei dati personali, di consultazione dello stato,
        di registrazione e di cancellazione delle richieste di ferie e permessi;
  \item autorizzazioni, conferme e gestione degli errori;
  \item la tracciabilità essenziale delle interazioni: persona, azienda, data e
        ora, fonti con la versione, capacità o abilità invocata, conferma ed
        esito quando applicabili.
\end{itemize}
Non sono richiesti:
\begin{itemize}
  \item onboarding, sportello interno e altri domini operativi;
  \item voce e allegati nella chat;
  \item memoria personale tra conversazioni diverse e uno storico di chat che
        il dipendente possa riaprire;
  \item la segmentazione delle fonti per gruppi, dipartimenti o sedi;
  \item la riproduzione del sistema di accesso di NEXUM;
  \item un sistema generale di workflow, un marketplace di plugin o un editor di
        abilità;
  \item l'estrazione automatica di regole strutturate e i grafi di conoscenza;
  \item l'integrazione nell'ambiente di produzione di NEXUM.
\end{itemize}
L'impianto deve però
permettere di aggiungere in futuro fonti e abilità senza riprogettare il
Concierge.

Nell'incontro del 2026-10-26 il proponente ha confermato il dominio ferie e
permessi, il PDF come prima tipologia di documento e l'italiano come unica
lingua. Nell'incontro del 2026-11-18 ha chiarito che non esistono ritiro o
modifica di una richiesta: la cancellazione è obbligatoria e la correzione
avviene cancellando la richiesta e inserendone una nuova. Ha inoltre chiarito
che il Concierge è un applicativo autonomo, che le API di Eggon saranno
autenticate con OAuth e che le loro specifiche saranno proposte dal gruppo.

\subsection{Utenti e caratteristiche}
\begin{itemize}
  \item \textbf{Personale delle aziende clienti}: usa la PWA, spesso da
        dispositivo mobile, senza competenze tecniche particolari; vuole
        risposte brevi, affidabili e riferite ai documenti della propria
        azienda.
  \item \textbf{HR e amministrazione}: configura nel backoffice le conoscenze
        disponibili, compresi i manuali NEXUM, verifica e aggiorna i documenti e
        ricostruisce le operazioni dalla traccia. Gli amministratori possono
        essere più di uno e sono tutti allo stesso livello (il capitolato parla
        anche di \textit{responsabili delle fonti}: nell'MVP il ruolo coincide
        con l'amministrazione).
\end{itemize}

\subsection{Assunzioni e dipendenze}
\begin{itemize}
  \item Eggon fornisce regolamenti di esempio, manuali NEXUM pertinenti e dati
        di prova fittizi su ferie e permessi.
  \item Eggon realizza, a partire da una proposta del gruppo, API autenticate con
        OAuth. Le API coprono lettura dei dati, consultazione dello stato,
        inserimento e cancellazione delle richieste e verifica dell'esito dopo
        un'interruzione. In loro assenza o per lo sviluppo si usano simulatori
        conformi al contratto.
  \item Eggon associa alcuni utenti di prova del Concierge a utenti NEXUM: il
        Concierge passa alle API l'identificativo dell'utente per cui opera.
  \item Il contratto definisce il significato dei dati e degli stati usati nei
        conteggi, la distinzione tra ferie fruite e richieste approvate, le
        unità di misura e il trattamento degli estremi dei periodi.
  \item Eggon mette a disposizione l'accesso ai modelli tramite la propria
        infrastruttura \glo{AWS}, con limiti di budget; finché non serve, si lavora con
        servizi simulati.
  \item Sono individuati referenti di Eggon per validare risposte e operazioni.
\end{itemize}

\section{Casi d'uso}\label{sec:uc}

\subsection{Metodo e notazione}
I casi d'uso descrivono il sistema dal punto di vista di chi lo usa: chi avvia
un'interazione, con quale obiettivo, cosa deve essere vero prima e dopo, e cosa
accade nei casi normali e in quelli eccezionali. Sono stati ricavati dagli
esempi e dai requisiti di business del capitolato, dai criteri di
accettazione e dall'incontro con il proponente.

\paragraph{Granularità} Un caso d'uso corrisponde a un obiettivo dell'\glo[Attore]{attore}. I
sottocasi (\textit{UCn.m}) sono introdotti solo quando una parte dello scenario
è un'azione o un gruppo di informazioni distinto, che dà origine a un proprio
requisito; non si scompone la visualizzazione campo per campo. Un comportamento
condiviso da più casi (per esempio il riepilogo e la conferma di
un'operazione) è un caso d'uso a sé, incluso da tutti quelli che lo usano.

\paragraph{Scheda} Ogni caso d'uso è descritto da una scheda con:
\begin{itemize}
  \item \textbf{descrizione}: l'obiettivo, quando non è evidente dal titolo;
  \item \textbf{attore principale} e \textbf{attori secondari}: chi avvia il caso
        d'uso e i sistemi esterni che vi partecipano;
  \item \textbf{precondizioni} e \textbf{postcondizioni}: lo stato del sistema
        prima dell'inizio e dopo la conclusione con successo;
  \item \textbf{scenario principale}: la sequenza di passi nel caso normale;
  \item \textbf{scenari alternativi}: le varianti che non meritano un caso d'uso
        a sé, numerate in base al passo da cui partono (per esempio 3a);
  \item \textbf{relazioni}: inclusioni, estensioni con la condizione che le
        attiva, generalizzazioni;
  \item \textbf{requisiti}: i requisiti funzionali che ne derivano.
\end{itemize}

\paragraph{Diagrammi} Ogni caso d'uso ha un diagramma \glo{UML}~2 che lo mostra nel suo
contesto: attori, casi inclusi, casi che lo includono o che estende, caso
generale da cui eredita. Se ha estensioni, un secondo diagramma le mostra con le
loro condizioni. Ogni sezione si apre con un diagramma di panoramica. La
notazione è riassunta nella Figura~\ref{fig:legenda}.

\begin{figure}[H]
  \centering
  \begin{adjustbox}{max width=\linewidth}
  \begin{tikzpicture}
    \attore{la}{(0,0)}{Attore}
    \node[uc] (l1) at (3.6,0) {\textbf{UCn}\\Caso d'uso};
    \draw[associazione] (la) -- (l1);
    \node[uc] (l2) at (8.4,0.9) {\textbf{UCn.1}\\Caso incluso};
    \draw[dipendenza] (l1) -- node[etich, sloped] {\guillemotleft include\guillemotright} (l2);
    \node[uc] (l3) at (8.4,-1.6) {\textbf{UCm}\\Estensione};
    \draw[dipendenza] (l3) -- node[etich, sloped, pos=0.35] {\guillemotleft extend\guillemotright} coordinate[pos=0.6] (lm) (l1);
    \node[nota] (ln) at (4.4,-2.4) {\textbf{Condizione:}\\\{quando si attiva\}};
    \draw[ancora] (ln) -- (lm);
    \node[uc] (l4) at (3.6,2.6) {\textbf{UCk}\\Caso generale};
    \draw[generalizzazione] (l1) -- (l4);
    \node[font=\sffamily\scriptsize, text=gptgrigio, align=left, anchor=west] at (4.5,1.55) {generalizzazione};
    \begin{scope}[on background layer]
      \node[confine, fit=(l1)(l2)(l3)(ln)(l4)] (sistema) {};
    \end{scope}
    \node[nomesistema] at (sistema.north west) {NEXUM Concierge};
  \end{tikzpicture}
  \end{adjustbox}
  \caption{Notazione dei diagrammi dei casi d'uso}\label{fig:legenda}
\end{figure}

Nella Figura~\ref{fig:legenda} il riquadro è il confine del sistema: gli
attori ne restano fuori. La freccia \guillemotleft include\guillemotright{} va
dal caso che include a quello incluso, che viene sempre eseguito; la freccia
\guillemotleft extend\guillemotright{} va dall'estensione al caso esteso, e
l'estensione avviene solo se vale la condizione riportata nella nota. La
generalizzazione, con il triangolo vuoto, va dal caso o dall'attore
specializzato a quello generale, di cui eredita relazioni e comportamento.

\paragraph{Codici} I casi d'uso sono identificati da \textit{UCn}, dove
\textit{n} è progressivo nel documento, e i sottocasi da \textit{UCn.m}.


\subsection{Attori}\label{sec:attori}
Gli attori sono i ruoli che interagiscono con il sistema. Gli attori principali
sono persone; gli attori secondari sono sistemi esterni che il Concierge usa
per completare un caso d'uso. Le relazioni di generalizzazione tra attori sono
riportate nella Figura~\ref{fig:attori}: un attore specializzato può svolgere
tutti i casi d'uso dell'attore generale.

\begin{figure}[H]
  \centering
  \begin{adjustbox}{max width=\linewidth}
  \begin{tikzpicture}
    \attore{una}{(-6.2,0)}{Utente non\\autenticato}
    \attore{ua}{(0,0)}{Utente\\autenticato}
    \attore{dip}{(-3.0,-3.4)}{Dipendente}
    \attore{amm}{(3.0,-3.4)}{Amministratore}
    \draw[generalizzazione] (dip.north) -- (ua-l.south west);
    \draw[generalizzazione] (amm.north) -- (ua-l.south east);
    \draw[gptlinea, dashed] (6.0,1.6) -- (6.0,-5.4);
    \node[font=\sffamily\scriptsize\bfseries, text=gptgrigio] at (8.6,1.4) {ATTORI SECONDARI};
    \attore{s2}{(8.6,-0.4)}{Servizio ferie e\\permessi NEXUM}
    \attore{s3}{(8.6,-3.6)}{Fornitore di\\modelli AI}
  \end{tikzpicture}
  \end{adjustbox}
  \caption{Attori del sistema e relative generalizzazioni}\label{fig:attori}
\end{figure}

\subsubsection{Attori principali}
\begin{itemize}
  \item \textbf{Utente non autenticato}: persona che non ha ancora effettuato
        l'accesso. Può solo autenticarsi.
  \item \textbf{Utente autenticato}: persona che ha effettuato l'accesso al
        Concierge; può uscire dal sistema. È specializzato dagli attori seguenti.
        L'accesso è gestito dall'applicativo, autonomo dal sistema di accesso di
        NEXUM: ogni utente ha un ruolo, un'azienda e un utente NEXUM associato.
  \item \textbf{Dipendente}: persona del personale di un'azienda cliente. Usa la
        chat per consultare le conoscenze della propria azienda e per operare
        sui propri dati di ferie e permessi.
  \item \textbf{Amministratore}: persona dell'HR o dell'amministrazione di
        un'azienda cliente. Nel backoffice gestisce i bucket e le fonti, compresi
        i manuali NEXUM, e consulta la traccia delle interazioni della propria
        azienda. Gli amministratori possono essere più di uno e sono tutti allo
        stesso livello.
\end{itemize}

\subsubsection{Attori secondari}
\begin{itemize}
  \item \textbf{Servizio ferie e permessi NEXUM}: fornisce i dati di ferie e
        permessi dell'utente NEXUM associato, registra, cancella e verifica le
        richieste e ne applica autorizzazioni e validazioni. È raggiunto tramite
        le API di Eggon, con chiamate autenticate con OAuth, e può essere reale o
        simulato.
  \item \textbf{Fornitore di modelli AI}: servizio esterno che espone i modelli
        linguistici e di \glo{Embedding} usati per interpretare i messaggi,
        indicizzare le fonti e formulare le risposte. Il sistema non dipende da
        un fornitore specifico.
\end{itemize}

\subsection{Elenco dei casi d'uso}
La Tabella~\ref{tab:elenco-uc} elenca tutti i casi d'uso, divisi nelle quattro
aree descritte di seguito.

{\small
\begin{longtable}{@{}p{1.9cm}>{\raggedright\arraybackslash}p{8.7cm}>{\raggedright\arraybackslash}p{4.2cm}@{}}
\caption{Elenco dei casi d'uso}\label{tab:elenco-uc}\\
\toprule \intestazione{Codice} & \intestazione{Titolo} & \intestazione{Attore principale} \\ \midrule\endfirsthead
\toprule \intestazione{Codice} & \intestazione{Titolo} & \intestazione{Attore principale} \\ \midrule\endhead
\midrule\multicolumn{3}{r@{}}{\footnotesize\itshape continua nella pagina successiva}\\ \endfoot \bottomrule\endlastfoot
\hyperref[uc:login]{UC1} & Autenticazione & Utente non autenticato \\
\hspace{1em}\hyperref[uc:login_id]{UC1.1} & Inserimento dell'identificativo & Utente non autenticato \\
\hspace{1em}\hyperref[uc:login_pw]{UC1.2} & Inserimento della password & Utente non autenticato \\
\hyperref[uc:cred_err]{UC2} & Credenziali non valide & Utente non autenticato \\
\hyperref[uc:logout]{UC3} & Uscita dal sistema & Utente autenticato \\
\hyperref[uc:nuova_conv]{UC4} & Avvio di una nuova conversazione & Dipendente \\
\hyperref[uc:vis_conv]{UC5} & Visualizzazione della conversazione & Dipendente \\
\hyperref[uc:msg]{UC6} & Invio di un messaggio al Concierge & Dipendente \\
\hyperref[uc:chiarimento]{UC7} & Richiesta di chiarimento & Dipendente \\
\hyperref[uc:fuori_ambito]{UC8} & Richiesta fuori dal perimetro & Dipendente \\
\hyperref[uc:msg_lungo]{UC9} & Messaggio troppo lungo & Dipendente \\
\hyperref[uc:errore_tecnico]{UC10} & Comunicazione di un errore tecnico & Dipendente \\
\hyperref[uc:regole]{UC11} & Domanda sulle regole aziendali & Dipendente \\
\hyperref[uc:fonti]{UC12} & Visualizzazione delle fonti citate & Dipendente \\
\hyperref[uc:astensione]{UC13} & Astensione dalla risposta & Dipendente \\
\hyperref[uc:prodotto]{UC14} & Domanda sull'utilizzo di NEXUM & Dipendente \\
\hyperref[uc:dati]{UC15} & Consultazione dei propri dati di ferie e permessi & Dipendente \\
\hspace{1em}\hyperref[uc:dati_riep]{UC15.1} & Visualizzazione delle richieste per periodo e stato & Dipendente \\
\hspace{1em}\hyperref[uc:dati_criteri]{UC15.2} & Visualizzazione di periodo, unità di misura e criterio di conteggio & Dipendente \\
\hyperref[uc:saldo]{UC16} & Visualizzazione del saldo residuo & Dipendente \\
\hyperref[uc:nessun_dato]{UC17} & Nessun dato corrispondente & Dipendente \\
\hyperref[uc:non_aut]{UC18} & Rifiuto di una richiesta non autorizzata & Dipendente \\
\hyperref[uc:stato_rich]{UC19} & Consultazione dello stato di una richiesta & Dipendente \\
\hyperref[uc:registra]{UC20} & Registrazione di una richiesta di ferie o permesso & Dipendente \\
\hspace{1em}\hyperref[uc:reg_tipo]{UC20.1} & Determinazione del tipo di richiesta & Dipendente \\
\hyperref[uc:periodo]{UC21} & Determinazione del periodo & Dipendente \\
\hspace{1em}\hyperref[uc:reg_campi]{UC20.2} & Raccolta degli altri dati richiesti & Dipendente \\
\hyperref[uc:op_riep]{UC22} & Visualizzazione del riepilogo dell'operazione & Dipendente \\
\hyperref[uc:op_conf]{UC23} & Conferma esplicita dell'operazione & Dipendente \\
\hyperref[uc:op_esito]{UC24} & Visualizzazione dell'esito dell'operazione & Dipendente \\
\hyperref[uc:op_modifica]{UC25} & Modifica dei dati prima della conferma & Dipendente \\
\hyperref[uc:op_rinuncia]{UC26} & Rinuncia all'operazione prima della conferma & Dipendente \\
\hyperref[uc:op_rifiuto]{UC27} & Operazione respinta dal servizio & Dipendente \\
\hyperref[uc:op_incerto]{UC28} & Esito dell'operazione incerto & Dipendente \\
\hyperref[uc:annulla_inv]{UC29} & Cancellazione di una richiesta & Dipendente \\
\hyperref[uc:sel_rich]{UC30} & Individuazione della richiesta interessata & Dipendente \\
\hyperref[uc:bk_elenco]{UC31} & Visualizzazione dell'elenco dei bucket & Amministratore \\
\hspace{1em}\hyperref[uc:bk_elem]{UC31.1} & Visualizzazione di un bucket nell'elenco & Amministratore \\
\hyperref[uc:bk_crea]{UC32} & Creazione di un bucket & Amministratore \\
\hspace{1em}\hyperref[uc:bk_nome]{UC32.1} & Inserimento del nome del bucket & Amministratore \\
\hspace{1em}\hyperref[uc:bk_desc]{UC32.2} & Inserimento della descrizione del bucket & Amministratore \\
\hyperref[uc:bk_dest]{UC33} & Limitazione della consultazione a gruppi & Amministratore \\
\hyperref[uc:bk_nome_dup]{UC34} & Nome del bucket già in uso & Amministratore \\
\hyperref[uc:bk_incompleto]{UC35} & Dati del bucket mancanti & Amministratore \\
\hyperref[uc:bk_modifica]{UC36} & Modifica di un bucket & Amministratore \\
\hyperref[uc:bk_elimina]{UC37} & Eliminazione di un bucket & Amministratore \\
\hyperref[uc:fn_elenco]{UC38} & Visualizzazione delle fonti di un bucket & Amministratore \\
\hspace{1em}\hyperref[uc:fn_elem]{UC38.1} & Visualizzazione di una fonte nell'elenco & Amministratore \\
\hyperref[uc:fn_carica]{UC39} & Caricamento di una fonte & Amministratore \\
\hspace{1em}\hyperref[uc:fn_file]{UC39.1} & Selezione del file & Amministratore \\
\hspace{1em}\hyperref[uc:fn_meta]{UC39.2} & Inserimento dei metadati della fonte & Amministratore \\
\hyperref[uc:fn_formato]{UC40} & Formato del file non supportato & Amministratore \\
\hyperref[uc:fn_dim]{UC41} & Dimensione massima superata & Amministratore \\
\hyperref[uc:fn_dup]{UC42} & Fonte già presente & Amministratore \\
\hyperref[uc:fn_dettaglio]{UC43} & Visualizzazione del dettaglio di una fonte & Amministratore \\
\hspace{1em}\hyperref[uc:fn_meta_v]{UC43.1} & Visualizzazione dei metadati & Amministratore \\
\hspace{1em}\hyperref[uc:fn_stato]{UC43.2} & Visualizzazione dello stato di elaborazione & Amministratore \\
\hyperref[uc:fn_errore]{UC44} & Visualizzazione dell'errore di elaborazione & Amministratore \\
\hyperref[uc:fn_sostituisci]{UC45} & Sostituzione di una fonte & Amministratore \\
\hyperref[uc:fn_ritira]{UC46} & Ritiro di una fonte & Amministratore \\
\hyperref[uc:fn_riprova]{UC47} & Nuova elaborazione di una fonte & Amministratore \\
\hyperref[uc:tr_elenco]{UC48} & Consultazione del registro delle interazioni & Amministratore \\
\hspace{1em}\hyperref[uc:tr_voce]{UC48.1} & Visualizzazione di una voce del registro & Amministratore \\
\hyperref[uc:tr_filtro]{UC49} & Filtraggio del registro & Amministratore \\
\hspace{1em}\hyperref[uc:tr_f_persona]{UC49.1} & Filtraggio per persona & Amministratore \\
\hspace{1em}\hyperref[uc:tr_f_periodo]{UC49.2} & Filtraggio per periodo & Amministratore \\
\hspace{1em}\hyperref[uc:tr_f_capacita]{UC49.3} & Filtraggio per capacità ed esito & Amministratore \\
\hyperref[uc:tr_dettaglio]{UC50} & Visualizzazione del dettaglio di un'interazione & Amministratore \\
\end{longtable}}

\subsection{Accesso al Concierge}\label{area:accesso}
Casi d'uso con cui una persona accede al Concierge e ne esce. L'accesso è gestito dall'applicativo, autonomo dal sistema di accesso di NEXUM: ogni utente di prova ha un ruolo (dipendente o amministratore) ed è associato a un'azienda e a un utente NEXUM, usato nelle chiamate alle API Eggon.

\begin{figure}[H]
  \centering
  \begin{adjustbox}{max width=\linewidth, max totalheight=0.62\textheight}
  \begin{tikzpicture}
    \node[uc] (p0) at (0.00,0.82) {\textbf{UC1}\\Autenticazione};
    \node[uc] (p1) at (0.00,-0.82) {\textbf{UC3}\\Uscita dal sistema};
    \begin{scope}[on background layer]
      \node[confine, fit=(p0) (p1)] (sistema) {};
    \end{scope}
    \node[nomesistema] at (sistema.north west) {NEXUM Concierge};
    \attore{ap0}{(-5.60,1.70)}{Utente\\non autenticato}
    \draw[associazione] (ap0) -- (p0);
    \attore{ap1}{(-5.60,-1.70)}{Utente\\autenticato}
    \draw[associazione] (ap1) -- (p1);
  \end{tikzpicture}
  \end{adjustbox}
  \caption{Panoramica: accesso al concierge}\label{fig:pano:accesso:0}
\end{figure}

\subsubsection{UC1 -- Autenticazione}\label{uc:login}
\begin{figure}[H]
  \centering
  \begin{adjustbox}{max width=\linewidth, max totalheight=0.62\textheight}
  \begin{tikzpicture}
    \node[ucfocus] (b) at (0.00,0.00) {\textbf{UC1}\\Autenticazione};
    \node[uc] (c0) at (7.15,1.02) {\textbf{UC1.1}\\Inserimento dell'identificativo};
    \draw[dipendenza] (b) -- node[etich, sloped] {\guillemotleft include\guillemotright} (c0);
    \node[uc] (c1) at (7.15,-1.02) {\textbf{UC1.2}\\Inserimento della password};
    \draw[dipendenza] (b) -- node[etich, sloped] {\guillemotleft include\guillemotright} (c1);
    \begin{scope}[on background layer]
      \node[confine, fit=(b) (c0) (c1)] (sistema) {};
    \end{scope}
    \node[nomesistema] at (sistema.north west) {NEXUM Concierge};
    \attore{ap0}{(-4.83,0.00)}{Utente\\non autenticato}
    \draw[associazione] (ap0) -- (b);
  \end{tikzpicture}
  \end{adjustbox}
  \caption{Diagramma di UC1 -- Autenticazione}\label{fig:uc:login}
\end{figure}

\begin{figure}[H]
  \centering
  \begin{adjustbox}{max width=\linewidth, max totalheight=0.62\textheight}
  \begin{tikzpicture}
    \node[ucfocus] (b) at (0.00,0.00) {\textbf{UC1}\\Autenticazione};
    \node[uc] (e0) at (6.60,0.00) {\textbf{UC2}\\Credenziali non valide};
    \draw[dipendenza] (e0) -- (b) node[etich, pos=0.3, sloped] {\guillemotleft extend\guillemotright} coordinate[pos=0.5] (m0);
    \node[nota] (n0) at (3.3,-1.9) {\textbf{Condizione:}\\\{credenziali non riconosciute\}};
    \draw[ancora] (n0) -- (m0);
    \begin{scope}[on background layer]
      \node[confine, fit=(b) (e0) (n0)] (sistema) {};
    \end{scope}
    \node[nomesistema] at (sistema.north west) {NEXUM Concierge};
    \attore{ap0}{(-4.23,0.00)}{Utente\\non autenticato}
    \draw[associazione] (ap0) -- (b);
  \end{tikzpicture}
  \end{adjustbox}
  \caption{Estensioni di UC1 -- Autenticazione}\label{fig:uc:login:ext}
\end{figure}

\begin{scheda}
Descrizione & L'utente si identifica per accedere alla chat del Concierge o al backoffice. L'accesso è gestito dall'applicativo, autonomo dal sistema di accesso di NEXUM: nell'MVP gli utenti sono quelli di prova concordati con Eggon, ciascuno con un ruolo e associato a un'azienda e a un utente NEXUM. \\ \addlinespace
Attore principale & Utente non autenticato \\ \addlinespace
Precondizioni & \begin{itemize}[nosep, leftmargin=*, before=\vspace{-\baselineskip}\vspace{2pt}, after=\vspace{-\baselineskip}]\item L'utente non è autenticato. \item È visualizzata la pagina di accesso. \end{itemize} \\ \addlinespace
Postcondizioni & \begin{itemize}[nosep, leftmargin=*, before=\vspace{-\baselineskip}\vspace{2pt}, after=\vspace{-\baselineskip}]\item L'utente è autenticato con il suo ruolo: il dipendente accede alla chat, l'amministratore al backoffice. \item Il sistema conosce l'azienda dell'utente e l'utente NEXUM associato, da usare nelle chiamate alle API Eggon. \end{itemize} \\ \addlinespace
Scenario principale & \begin{enumerate}[nosep, leftmargin=*, before=\vspace{-\baselineskip}\vspace{2pt}, after=\vspace{-\baselineskip}]\item L'utente inserisce il proprio identificativo (\hyperref[uc:login_id]{UC1.1}). \item L'utente inserisce la password (\hyperref[uc:login_pw]{UC1.2}). \item L'utente conferma l'accesso. \item Il sistema verifica le credenziali. \item Il sistema riconosce ruolo e azienda dell'utente, ricava l'utente NEXUM associato secondo la mappatura e avvia la sessione. \end{enumerate} \\ \addlinespace
Inclusioni & \hyperref[uc:login_id]{UC1.1}, \hyperref[uc:login_pw]{UC1.2} \\ \addlinespace
Estensioni & \hyperref[uc:cred_err]{UC2} \{credenziali non riconosciute\} \\ \addlinespace
Requisiti & \hyperlink{req:RF-01}{RF-01}, \hyperlink{req:RF-02}{RF-02}, \hyperlink{req:RF-07}{RF-07} \\
\end{scheda}

\subsubsubsection{UC1.1 -- Inserimento dell'identificativo}\label{uc:login_id}
\begin{figure}[H]
  \centering
  \begin{adjustbox}{max width=\linewidth, max totalheight=0.62\textheight}
  \begin{tikzpicture}
    \node[ucfocus] (b) at (0.00,0.00) {\textbf{UC1.1}\\Inserimento dell'identificativo};
    \node[uc] (c0) at (7.15,0.00) {\textbf{UC1}\\Autenticazione};
    \draw[dipendenza] (c0) -- node[etich, sloped] {\guillemotleft include\guillemotright} (b);
    \begin{scope}[on background layer]
      \node[confine, fit=(b) (c0)] (sistema) {};
    \end{scope}
    \node[nomesistema] at (sistema.north west) {NEXUM Concierge};
    \attore{ap0}{(-4.83,0.00)}{Utente\\non autenticato}
    \draw[associazione] (ap0) -- (b);
  \end{tikzpicture}
  \end{adjustbox}
  \caption{Diagramma di UC1.1 -- Inserimento dell'identificativo}\label{fig:uc:login_id}
\end{figure}

\begin{scheda}
Attore principale & Utente non autenticato \\ \addlinespace
Precondizioni & È visualizzata la pagina di accesso. \\ \addlinespace
Postcondizioni & L'identificativo è inserito. \\ \addlinespace
Scenario principale & L'utente inserisce il proprio identificativo (quello assegnato all'utente di prova) nel campo dedicato. \\ \addlinespace
Incluso da & \hyperref[uc:login]{UC1} \\ \addlinespace
Requisiti & \hyperlink{req:RF-03}{RF-03} \\
\end{scheda}

\subsubsubsection{UC1.2 -- Inserimento della password}\label{uc:login_pw}
\begin{figure}[H]
  \centering
  \begin{adjustbox}{max width=\linewidth, max totalheight=0.62\textheight}
  \begin{tikzpicture}
    \node[ucfocus] (b) at (0.00,0.00) {\textbf{UC1.2}\\Inserimento della password};
    \node[uc] (c0) at (7.15,0.00) {\textbf{UC1}\\Autenticazione};
    \draw[dipendenza] (c0) -- node[etich, sloped] {\guillemotleft include\guillemotright} (b);
    \begin{scope}[on background layer]
      \node[confine, fit=(b) (c0)] (sistema) {};
    \end{scope}
    \node[nomesistema] at (sistema.north west) {NEXUM Concierge};
    \attore{ap0}{(-4.83,0.00)}{Utente\\non autenticato}
    \draw[associazione] (ap0) -- (b);
  \end{tikzpicture}
  \end{adjustbox}
  \caption{Diagramma di UC1.2 -- Inserimento della password}\label{fig:uc:login_pw}
\end{figure}

\begin{scheda}
Attore principale & Utente non autenticato \\ \addlinespace
Precondizioni & È visualizzata la pagina di accesso. \\ \addlinespace
Postcondizioni & La password è inserita e non è mostrata in chiaro. \\ \addlinespace
Scenario principale & L'utente inserisce la password nel campo dedicato. \\ \addlinespace
Incluso da & \hyperref[uc:login]{UC1} \\ \addlinespace
Requisiti & \hyperlink{req:RF-04}{RF-04} \\
\end{scheda}

\subsubsection{UC2 -- Credenziali non valide}\label{uc:cred_err}
\begin{figure}[H]
  \centering
  \begin{adjustbox}{max width=\linewidth, max totalheight=0.62\textheight}
  \begin{tikzpicture}
    \node[ucfocus] (b) at (0.00,0.00) {\textbf{UC2}\\Credenziali non valide};
    \node[uc] (c0) at (7.15,0.00) {\textbf{UC1}\\Autenticazione};
    \draw[dipendenza] (b) -- node[etich, sloped] {\guillemotleft extend\guillemotright} (c0);
    \begin{scope}[on background layer]
      \node[confine, fit=(b) (c0)] (sistema) {};
    \end{scope}
    \node[nomesistema] at (sistema.north west) {NEXUM Concierge};
    \attore{ap0}{(-4.83,0.00)}{Utente\\non autenticato}
    \draw[associazione] (ap0) -- (b);
  \end{tikzpicture}
  \end{adjustbox}
  \caption{Diagramma di UC2 -- Credenziali non valide}\label{fig:uc:cred_err}
\end{figure}

\begin{scheda}
Attore principale & Utente non autenticato \\ \addlinespace
Precondizioni & L'utente ha confermato l'accesso con credenziali che il sistema non riconosce. \\ \addlinespace
Postcondizioni & L'utente non è autenticato e può ripetere l'accesso. \\ \addlinespace
Scenario principale & \begin{enumerate}[nosep, leftmargin=*, before=\vspace{-\baselineskip}\vspace{2pt}, after=\vspace{-\baselineskip}]\item Il sistema mostra un messaggio d'errore che non indica quale dei due dati è errato. \item Il sistema svuota il campo della password. \end{enumerate} \\ \addlinespace
Estende & \hyperref[uc:login]{UC1} \{credenziali non riconosciute\} \\ \addlinespace
Requisiti & \hyperlink{req:RF-05}{RF-05} \\
\end{scheda}

\subsubsection{UC3 -- Uscita dal sistema}\label{uc:logout}
\begin{figure}[H]
  \centering
  \begin{adjustbox}{max width=\linewidth, max totalheight=0.62\textheight}
  \begin{tikzpicture}
    \node[ucfocus] (b) at (0.00,0.00) {\textbf{UC3}\\Uscita dal sistema};
    \begin{scope}[on background layer]
      \node[confine, fit=(b)] (sistema) {};
    \end{scope}
    \node[nomesistema] at (sistema.north west) {NEXUM Concierge};
    \attore{ap0}{(-4.83,0.00)}{Utente\\autenticato}
    \draw[associazione] (ap0) -- (b);
  \end{tikzpicture}
  \end{adjustbox}
  \caption{Diagramma di UC3 -- Uscita dal sistema}\label{fig:uc:logout}
\end{figure}

\begin{scheda}
Descrizione & L'utente termina la sessione. \\ \addlinespace
Attore principale & Utente autenticato \\ \addlinespace
Precondizioni & L'utente è autenticato. \\ \addlinespace
Postcondizioni & La sessione è chiusa e il contesto della conversazione è eliminato. \\ \addlinespace
Scenario principale & \begin{enumerate}[nosep, leftmargin=*, before=\vspace{-\baselineskip}\vspace{2pt}, after=\vspace{-\baselineskip}]\item L'utente sceglie di uscire. \item Il sistema chiude la sessione e mostra la pagina di accesso. \end{enumerate} \\ \addlinespace
Requisiti & \hyperlink{req:RF-08}{RF-08} \\
\end{scheda}

\subsection{Conversazione con il Concierge}\label{area:chat}
Casi d'uso della chat testuale della PWA. Il caso d'uso generale \hyperref[uc:msg]{UC6} descrive ciò che accomuna ogni messaggio (contesto, scelta della capacità, chiarimenti, errori); i casi d'uso che lo specializzano descrivono le quattro capacità richieste dal capitolato e le operazioni aggiuntive emerse dall'incontro con il proponente.

\begin{figure}[H]
  \centering
  \begin{adjustbox}{max width=\linewidth, max totalheight=0.62\textheight}
  \begin{tikzpicture}
    \node[uc] (p0) at (0.00,2.12) {\textbf{UC4}\\Avvio di una nuova conversazione};
    \node[uc] (p1) at (0.00,0.12) {\textbf{UC5}\\Visualizzazione della conversazione};
    \node[uc] (p2) at (0.00,-2.12) {\textbf{UC6}\\Invio di un messaggio al Concierge};
    \begin{scope}[on background layer]
      \node[confine, fit=(p0) (p1) (p2)] (sistema) {};
    \end{scope}
    \node[nomesistema] at (sistema.north west) {NEXUM Concierge};
    \attore{ap0}{(-5.60,0.00)}{Dipendente}
    \draw[associazione] (ap0) -- (p0);
    \draw[associazione] (ap0) -- (p1);
    \draw[associazione] (ap0) -- (p2);
    \attore{as0}{(5.60,0.00)}{Fornitore\\di modelli AI}
    \draw[associazione] (as0) -- (p2);
  \end{tikzpicture}
  \end{adjustbox}
  \caption{Panoramica: conversazione con il concierge}\label{fig:pano:chat:0}
\end{figure}

\subsubsection{UC4 -- Avvio di una nuova conversazione}\label{uc:nuova_conv}
\begin{figure}[H]
  \centering
  \begin{adjustbox}{max width=\linewidth, max totalheight=0.62\textheight}
  \begin{tikzpicture}
    \node[ucfocus] (b) at (0.00,0.00) {\textbf{UC4}\\Avvio di una nuova conversazione};
    \begin{scope}[on background layer]
      \node[confine, fit=(b)] (sistema) {};
    \end{scope}
    \node[nomesistema] at (sistema.north west) {NEXUM Concierge};
    \attore{ap0}{(-4.83,0.00)}{Dipendente}
    \draw[associazione] (ap0) -- (b);
  \end{tikzpicture}
  \end{adjustbox}
  \caption{Diagramma di UC4 -- Avvio di una nuova conversazione}\label{fig:uc:nuova_conv}
\end{figure}

\begin{scheda}
Descrizione & Ogni apertura della chat avvia una nuova conversazione, e il dipendente può ricominciare da capo in qualsiasi momento. Il Concierge non conserva memoria tra conversazioni diverse e non offre uno storico di chat da riaprire. \\ \addlinespace
Attore principale & Dipendente \\ \addlinespace
Precondizioni & Il dipendente è autenticato. \\ \addlinespace
Postcondizioni & \begin{itemize}[nosep, leftmargin=*, before=\vspace{-\baselineskip}\vspace{2pt}, after=\vspace{-\baselineskip}]\item La chat è vuota e il contesto della conversazione precedente non è più utilizzato. \item È mostrato l'avviso su cosa può fare il Concierge. \end{itemize} \\ \addlinespace
Scenario principale & \begin{enumerate}[nosep, leftmargin=*, before=\vspace{-\baselineskip}\vspace{2pt}, after=\vspace{-\baselineskip}]\item Il dipendente apre la chat o sceglie di avviare una nuova conversazione. \item Il sistema svuota la chat e azzera il contesto. \item Il sistema mostra un avviso che spiega cosa può fare il Concierge e invita a non inserire informazioni riservate. \end{enumerate} \\ \addlinespace
Scenari alternativi & 1a. Una richiesta di operazione è in preparazione e non ancora confermata: il sistema avvisa che verrà abbandonata e chiede conferma. \\ \addlinespace
Requisiti & \hyperlink{req:RF-09}{RF-09}, \hyperlink{req:RF-98}{RF-98} \\
\end{scheda}

\subsubsection{UC5 -- Visualizzazione della conversazione}\label{uc:vis_conv}
\begin{figure}[H]
  \centering
  \begin{adjustbox}{max width=\linewidth, max totalheight=0.62\textheight}
  \begin{tikzpicture}
    \node[ucfocus] (b) at (0.00,0.00) {\textbf{UC5}\\Visualizzazione della conversazione};
    \begin{scope}[on background layer]
      \node[confine, fit=(b)] (sistema) {};
    \end{scope}
    \node[nomesistema] at (sistema.north west) {NEXUM Concierge};
    \attore{ap0}{(-4.83,0.00)}{Dipendente}
    \draw[associazione] (ap0) -- (b);
  \end{tikzpicture}
  \end{adjustbox}
  \caption{Diagramma di UC5 -- Visualizzazione della conversazione}\label{fig:uc:vis_conv}
\end{figure}

\begin{scheda}
Descrizione & Il dipendente rilegge i messaggi della conversazione in corso. \\ \addlinespace
Attore principale & Dipendente \\ \addlinespace
Precondizioni & Il dipendente ha una conversazione aperta. \\ \addlinespace
Postcondizioni & La conversazione è visualizzata. \\ \addlinespace
Scenario principale & \begin{enumerate}[nosep, leftmargin=*, before=\vspace{-\baselineskip}\vspace{2pt}, after=\vspace{-\baselineskip}]\item Il sistema mostra i messaggi in ordine cronologico, distinguendo quelli del dipendente da quelli del Concierge e, se previsto, l'orario di ciascuno. \item Per ogni risposta il sistema mostra le fonti citate o l'esito dell'operazione, quando presenti. \end{enumerate} \\ \addlinespace
Requisiti & \hyperlink{req:RF-10}{RF-10}, \hyperlink{req:RF-97}{RF-97} \\
\end{scheda}

\subsubsection{UC6 -- Invio di un messaggio al Concierge}\label{uc:msg}
\begin{figure}[H]
  \centering
  \begin{adjustbox}{max width=\linewidth, max totalheight=0.62\textheight}
  \begin{tikzpicture}
    \node[ucfocus] (b) at (0.00,0.00) {\textbf{UC6}\\Invio di un messaggio al Concierge};
    \draw[generalizzazione] (0,-1.70) -- (b);
    \node[uc] (f00) at (-7.05,-3.34) {\textbf{UC11}\\Domanda sulle regole aziendali};
    \draw[albero] (f00.north) -- (-7.05,-1.70);
    \node[uc] (f01) at (-2.35,-3.34) {\textbf{UC14}\\Domanda sull'utilizzo di NEXUM};
    \draw[albero] (f01.north) -- (-2.35,-1.70);
    \node[uc] (f02) at (2.35,-3.34) {\textbf{UC15}\\Consultazione dei propri dati di ferie e permessi};
    \draw[albero] (f02.north) -- (2.35,-1.70);
    \node[uc] (f03) at (7.05,-3.34) {\textbf{UC19}\\Consultazione dello stato di una richiesta};
    \draw[albero] (f03.north) -- (7.05,-1.70);
    \node[uc] (f10) at (-2.35,-6.32) {\textbf{UC20}\\Registrazione di una richiesta di ferie o permesso};
    \draw[albero] (f10.north) -- (-2.35,-1.70);
    \node[uc] (f11) at (2.35,-6.32) {\textbf{UC29}\\Cancellazione di una richiesta};
    \draw[albero] (f11.north) -- (2.35,-1.70);
    \draw[albero] (-7.05,-1.70) -- (7.05,-1.70);
    \begin{scope}[on background layer]
      \node[confine, fit=(b) (f00) (f01) (f02) (f03) (f10) (f11)] (sistema) {};
    \end{scope}
    \node[nomesistema] at (sistema.north west) {NEXUM Concierge};
    \attore{ap0}{(-11.07,0.00)}{Dipendente}
    \draw[associazione] (ap0) -- (b);
    \attore{as0}{(11.07,0.00)}{Fornitore\\di modelli AI}
    \draw[associazione] (b) -- (as0);
  \end{tikzpicture}
  \end{adjustbox}
  \caption{Generalizzazione di UC6 -- Invio di un messaggio al Concierge}\label{fig:uc:msg}
\end{figure}

\begin{figure}[H]
  \centering
  \begin{adjustbox}{max width=\linewidth, max totalheight=0.62\textheight}
  \begin{tikzpicture}
    \node[ucfocus] (b) at (0.00,0.00) {\textbf{UC6}\\Invio di un messaggio al Concierge};
    \node[uc] (e0) at (-5.00,4.00) {\textbf{UC7}\\Richiesta di chiarimento};
    \draw[dipendenza] (e0) -- (b) node[etich, pos=0.3, sloped] {\guillemotleft extend\guillemotright} coordinate[pos=0.5] (m0);
    \node[nota] (n0) at (-6.45,1.75) {\textbf{Condizione:}\\\{formulazione ambigua o informazioni mancanti\}};
    \draw[ancora] (n0) -- (m0);
    \node[uc] (e1) at (5.00,4.00) {\textbf{UC8}\\Richiesta fuori dal perimetro};
    \draw[dipendenza] (e1) -- (b) node[etich, pos=0.3, sloped] {\guillemotleft extend\guillemotright} coordinate[pos=0.5] (m1);
    \node[nota] (n1) at (6.45,1.75) {\textbf{Condizione:}\\\{richiesta estranea a ferie, permessi e uso di NEXUM\}};
    \draw[ancora] (n1) -- (m1);
    \node[uc] (e2) at (5.00,-4.00) {\textbf{UC9}\\Messaggio troppo lungo};
    \draw[dipendenza] (e2) -- (b) node[etich, pos=0.3, sloped] {\guillemotleft extend\guillemotright} coordinate[pos=0.5] (m2);
    \node[nota] (n2) at (6.45,-1.75) {\textbf{Condizione:}\\\{messaggio oltre la lunghezza massima\}};
    \draw[ancora] (n2) -- (m2);
    \node[uc] (e3) at (-5.00,-4.00) {\textbf{UC10}\\Comunicazione di un errore tecnico};
    \draw[dipendenza] (e3) -- (b) node[etich, pos=0.3, sloped] {\guillemotleft extend\guillemotright} coordinate[pos=0.5] (m3);
    \node[nota] (n3) at (-6.45,-1.75) {\textbf{Condizione:}\\\{modello o servizio non disponibile o in errore\}};
    \draw[ancora] (n3) -- (m3);
    \begin{scope}[on background layer]
      \node[confine, fit=(b) (e0) (n0) (e1) (n1) (e2) (n2) (e3) (n3)] (sistema) {};
    \end{scope}
    \node[nomesistema] at (sistema.north west) {NEXUM Concierge};
    \attore{ap0}{(-10.25,0.00)}{Dipendente}
    \draw[associazione] (ap0) -- (b);
  \end{tikzpicture}
  \end{adjustbox}
  \caption{Estensioni di UC6 -- Invio di un messaggio al Concierge}\label{fig:uc:msg:ext}
\end{figure}

\begin{scheda}
Descrizione & Il dipendente scrive in linguaggio naturale una domanda o una richiesta. Il sistema la interpreta nel contesto della conversazione, sceglie la capacità o le capacità necessarie e risponde. È il caso d'uso generale: le capacità sono descritte dai casi d'uso che lo specializzano. \\ \addlinespace
Attore principale & Dipendente \\ \addlinespace
Attori secondari & Fornitore di modelli AI \\ \addlinespace
Precondizioni & \begin{itemize}[nosep, leftmargin=*, before=\vspace{-\baselineskip}\vspace{2pt}, after=\vspace{-\baselineskip}]\item Il dipendente è autenticato. \item Una conversazione è aperta. \end{itemize} \\ \addlinespace
Postcondizioni & \begin{itemize}[nosep, leftmargin=*, before=\vspace{-\baselineskip}\vspace{2pt}, after=\vspace{-\baselineskip}]\item La risposta è mostrata nella conversazione. \item Messaggio e risposta fanno parte del contesto della conversazione. \item L'interazione è registrata nella traccia. \end{itemize} \\ \addlinespace
Scenario principale & \begin{enumerate}[nosep, leftmargin=*, before=\vspace{-\baselineskip}\vspace{2pt}, after=\vspace{-\baselineskip}]\item Il dipendente scrive il messaggio e lo invia. \item Il sistema interpreta il messaggio tenendo conto dei messaggi precedenti della conversazione. \item Il sistema riconosce la capacità richiesta: conoscenza aziendale, conoscenza del prodotto, consultazione dei dati o esecuzione di un'operazione. \item Il sistema svolge la capacità con la fonte o l'abilità corrispondente. \item Il sistema compone la risposta e la mostra. \item Il sistema registra l'interazione nella traccia, con le fonti e le versioni usate o l'abilità invocata. \end{enumerate} \\ \addlinespace
Scenari alternativi & \begin{itemize}[nosep, leftmargin=*, before=\vspace{-\baselineskip}\vspace{2pt}, after=\vspace{-\baselineskip}]\item 2a. Il messaggio rimanda a elementi già citati («e a dicembre?»): il sistema li ricava dal contesto della conversazione. \item 3a. Il messaggio richiede più capacità («Quante ferie ho fatto e come ne chiedo altre?»): il sistema le svolge in sequenza e compone un'unica risposta, in cui ogni parte indica la propria fonte. \end{itemize} \\ \addlinespace
Estensioni & \hyperref[uc:chiarimento]{UC7} \{formulazione ambigua o informazioni mancanti\}; \hyperref[uc:fuori_ambito]{UC8} \{richiesta estranea a ferie, permessi e uso di NEXUM\}; \hyperref[uc:msg_lungo]{UC9} \{messaggio oltre la lunghezza massima\}; \hyperref[uc:errore_tecnico]{UC10} \{modello o servizio non disponibile o in errore\} \\ \addlinespace
Specializzato da & \hyperref[uc:regole]{UC11}, \hyperref[uc:prodotto]{UC14}, \hyperref[uc:dati]{UC15}, \hyperref[uc:stato_rich]{UC19}, \hyperref[uc:registra]{UC20}, \hyperref[uc:annulla_inv]{UC29} \\ \addlinespace
Requisiti & \hyperlink{req:RF-11}{RF-11}, \hyperlink{req:RF-12}{RF-12}, \hyperlink{req:RF-13}{RF-13}, \hyperlink{req:RF-14}{RF-14}, \hyperlink{req:RF-15}{RF-15}, \hyperlink{req:RF-99}{RF-99} \\
\end{scheda}

\subsubsection{UC7 -- Richiesta di chiarimento}\label{uc:chiarimento}
\begin{figure}[H]
  \centering
  \begin{adjustbox}{max width=\linewidth, max totalheight=0.62\textheight}
  \begin{tikzpicture}
    \node[ucfocus] (b) at (0.00,0.00) {\textbf{UC7}\\Richiesta di chiarimento};
    \node[uc] (c0) at (7.15,1.27) {\textbf{UC6}\\Invio di un messaggio al Concierge};
    \draw[dipendenza] (b) -- node[etich, sloped] {\guillemotleft extend\guillemotright} (c0);
    \node[uc] (c1) at (7.15,-1.27) {\textbf{UC20}\\Registrazione di una richiesta di ferie o permesso};
    \draw[dipendenza] (b) -- node[etich, sloped] {\guillemotleft extend\guillemotright} (c1);
    \begin{scope}[on background layer]
      \node[confine, fit=(b) (c0) (c1)] (sistema) {};
    \end{scope}
    \node[nomesistema] at (sistema.north west) {NEXUM Concierge};
    \attore{ap0}{(-4.83,0.00)}{Dipendente}
    \draw[associazione] (ap0) -- (b);
  \end{tikzpicture}
  \end{adjustbox}
  \caption{Diagramma di UC7 -- Richiesta di chiarimento}\label{fig:uc:chiarimento}
\end{figure}

\begin{scheda}
Descrizione & Il Concierge chiede l'informazione che manca o che ammette più interpretazioni, invece di supporla. \\ \addlinespace
Attore principale & Dipendente \\ \addlinespace
Precondizioni & Il messaggio è ambiguo o incompleto rispetto al periodo, all'anno, al tipo di richiesta o alla durata. \\ \addlinespace
Postcondizioni & Il sistema dispone dell'informazione e prosegue con la richiesta originaria; l'informazione resta nel contesto. \\ \addlinespace
Scenario principale & \begin{enumerate}[nosep, leftmargin=*, before=\vspace{-\baselineskip}\vspace{2pt}, after=\vspace{-\baselineskip}]\item Il sistema indica l'elemento ambiguo o mancante e pone una domanda mirata, proponendo se possibile le alternative («Intendi dal 10 al 15 ottobre 2026?»). \item Il dipendente risponde. \item Il sistema riprende la richiesta originaria con l'informazione ottenuta. \end{enumerate} \\ \addlinespace
Scenari alternativi & \begin{itemize}[nosep, leftmargin=*, before=\vspace{-\baselineskip}\vspace{2pt}, after=\vspace{-\baselineskip}]\item 1a. Una data è priva di anno ma l'anno è già esplicitato nella conversazione: il sistema usa quell'anno senza chiedere, e lo riporta per esteso nella risposta o nel riepilogo. \item 2a. Il dipendente cambia argomento: la richiesta originaria è abbandonata e nessuna operazione viene inviata. \end{itemize} \\ \addlinespace
Estende & \hyperref[uc:msg]{UC6} \{formulazione ambigua o informazioni mancanti\}; \hyperref[uc:registra]{UC20} \{dati ambigui o mancanti\} \\ \addlinespace
Requisiti & \hyperlink{req:RF-16}{RF-16}, \hyperlink{req:RF-17}{RF-17} \\
\end{scheda}

\subsubsection{UC8 -- Richiesta fuori dal perimetro}\label{uc:fuori_ambito}
\begin{figure}[H]
  \centering
  \begin{adjustbox}{max width=\linewidth, max totalheight=0.62\textheight}
  \begin{tikzpicture}
    \node[ucfocus] (b) at (0.00,0.00) {\textbf{UC8}\\Richiesta fuori dal perimetro};
    \node[uc] (c0) at (7.15,0.00) {\textbf{UC6}\\Invio di un messaggio al Concierge};
    \draw[dipendenza] (b) -- node[etich, sloped] {\guillemotleft extend\guillemotright} (c0);
    \begin{scope}[on background layer]
      \node[confine, fit=(b) (c0)] (sistema) {};
    \end{scope}
    \node[nomesistema] at (sistema.north west) {NEXUM Concierge};
    \attore{ap0}{(-4.83,0.00)}{Dipendente}
    \draw[associazione] (ap0) -- (b);
  \end{tikzpicture}
  \end{adjustbox}
  \caption{Diagramma di UC8 -- Richiesta fuori dal perimetro}\label{fig:uc:fuori_ambito}
\end{figure}

\begin{scheda}
Attore principale & Dipendente \\ \addlinespace
Precondizioni & Il messaggio riguarda un argomento estraneo a ferie, permessi e uso di NEXUM. \\ \addlinespace
Postcondizioni & Nessuna risposta viene inventata; il dipendente riceve un messaggio che spiega il limite del servizio e sa di cosa può occuparsi il Concierge. \\ \addlinespace
Scenario principale & \begin{enumerate}[nosep, leftmargin=*, before=\vspace{-\baselineskip}\vspace{2pt}, after=\vspace{-\baselineskip}]\item Il sistema spiega che la richiesta è fuori dal perimetro del servizio. \item Il sistema ricorda gli argomenti su cui può aiutare. \end{enumerate} \\ \addlinespace
Estende & \hyperref[uc:msg]{UC6} \{richiesta estranea a ferie, permessi e uso di NEXUM\} \\ \addlinespace
Requisiti & \hyperlink{req:RF-18}{RF-18} \\
\end{scheda}

\subsubsection{UC9 -- Messaggio troppo lungo}\label{uc:msg_lungo}
\begin{figure}[H]
  \centering
  \begin{adjustbox}{max width=\linewidth, max totalheight=0.62\textheight}
  \begin{tikzpicture}
    \node[ucfocus] (b) at (0.00,0.00) {\textbf{UC9}\\Messaggio troppo lungo};
    \node[uc] (c0) at (7.15,0.00) {\textbf{UC6}\\Invio di un messaggio al Concierge};
    \draw[dipendenza] (b) -- node[etich, sloped] {\guillemotleft extend\guillemotright} (c0);
    \begin{scope}[on background layer]
      \node[confine, fit=(b) (c0)] (sistema) {};
    \end{scope}
    \node[nomesistema] at (sistema.north west) {NEXUM Concierge};
    \attore{ap0}{(-4.83,0.00)}{Dipendente}
    \draw[associazione] (ap0) -- (b);
  \end{tikzpicture}
  \end{adjustbox}
  \caption{Diagramma di UC9 -- Messaggio troppo lungo}\label{fig:uc:msg_lungo}
\end{figure}

\begin{scheda}
Attore principale & Dipendente \\ \addlinespace
Precondizioni & Il dipendente ha scritto un messaggio che supera la lunghezza massima consentita. \\ \addlinespace
Postcondizioni & Il messaggio non è inviato e il testo resta modificabile. \\ \addlinespace
Scenario principale & Il sistema impedisce l'invio e indica la lunghezza massima. \\ \addlinespace
Estende & \hyperref[uc:msg]{UC6} \{messaggio oltre la lunghezza massima\} \\ \addlinespace
Requisiti & \hyperlink{req:RF-19}{RF-19} \\
\end{scheda}

\subsubsection{UC10 -- Comunicazione di un errore tecnico}\label{uc:errore_tecnico}
\begin{figure}[H]
  \centering
  \begin{adjustbox}{max width=\linewidth, max totalheight=0.62\textheight}
  \begin{tikzpicture}
    \node[ucfocus] (b) at (0.00,0.00) {\textbf{UC10}\\Comunicazione di un errore tecnico};
    \node[uc] (c0) at (7.15,0.00) {\textbf{UC6}\\Invio di un messaggio al Concierge};
    \draw[dipendenza] (b) -- node[etich, sloped] {\guillemotleft extend\guillemotright} (c0);
    \begin{scope}[on background layer]
      \node[confine, fit=(b) (c0)] (sistema) {};
    \end{scope}
    \node[nomesistema] at (sistema.north west) {NEXUM Concierge};
    \attore{ap0}{(-4.83,0.00)}{Dipendente}
    \draw[associazione] (ap0) -- (b);
  \end{tikzpicture}
  \end{adjustbox}
  \caption{Diagramma di UC10 -- Comunicazione di un errore tecnico}\label{fig:uc:errore_tecnico}
\end{figure}

\begin{scheda}
Attore principale & Dipendente \\ \addlinespace
Precondizioni & Il modello, le API Eggon o un'altra componente non risponde o restituisce un errore. \\ \addlinespace
Postcondizioni & Il dipendente sa che la richiesta non è stata completata e se può riprovare; nessuna operazione è considerata eseguita senza un esito verificato. \\ \addlinespace
Scenario principale & \begin{enumerate}[nosep, leftmargin=*, before=\vspace{-\baselineskip}\vspace{2pt}, after=\vspace{-\baselineskip}]\item Il sistema comunica con un messaggio comprensibile che la richiesta non è stata completata per un problema tecnico. \item Il sistema indica se e quando è possibile riprovare. \item Il sistema registra l'errore nella traccia, con i dettagli utili all'amministratore. \end{enumerate} \\ \addlinespace
Estende & \hyperref[uc:msg]{UC6} \{modello o servizio non disponibile o in errore\} \\ \addlinespace
Requisiti & \hyperlink{req:RF-20}{RF-20} \\
\end{scheda}

\subsubsection{UC11 -- Domanda sulle regole aziendali}\label{uc:regole}
\begin{figure}[H]
  \centering
  \begin{adjustbox}{max width=\linewidth, max totalheight=0.62\textheight}
  \begin{tikzpicture}
    \node[ucfocus] (b) at (0.00,0.00) {\textbf{UC11}\\Domanda sulle regole aziendali};
    \node[uc] (c0) at (7.15,0.00) {\textbf{UC12}\\Visualizzazione delle fonti citate};
    \draw[dipendenza] (b) -- node[etich, sloped] {\guillemotleft include\guillemotright} (c0);
    \node[uc] (g) at (0.00,2.84) {\textbf{UC6}\\Invio di un messaggio al Concierge};
    \draw[generalizzazione] (b) -- (g);
    \begin{scope}[on background layer]
      \node[confine, fit=(b) (c0) (g)] (sistema) {};
    \end{scope}
    \node[nomesistema] at (sistema.north west) {NEXUM Concierge};
    \attore{ap0}{(-4.83,0.00)}{Dipendente}
    \draw[associazione] (ap0) -- (b);
    \attore{as0}{(0.00,-2.85)}{Fornitore\\di modelli AI}
    \draw[associazione] (b) -- (as0);
  \end{tikzpicture}
  \end{adjustbox}
  \caption{Diagramma di UC11 -- Domanda sulle regole aziendali}\label{fig:uc:regole}
\end{figure}

\begin{figure}[H]
  \centering
  \begin{adjustbox}{max width=\linewidth, max totalheight=0.62\textheight}
  \begin{tikzpicture}
    \node[ucfocus] (b) at (0.00,0.00) {\textbf{UC11}\\Domanda sulle regole aziendali};
    \node[uc] (e0) at (6.60,0.00) {\textbf{UC13}\\Astensione dalla risposta};
    \draw[dipendenza] (e0) -- (b) node[etich, pos=0.3, sloped] {\guillemotleft extend\guillemotright} coordinate[pos=0.5] (m0);
    \node[nota] (n0) at (3.3,-1.9) {\textbf{Condizione:}\\\{fonti insufficienti o in contrasto non risolvibile\}};
    \draw[ancora] (n0) -- (m0);
    \begin{scope}[on background layer]
      \node[confine, fit=(b) (e0) (n0)] (sistema) {};
    \end{scope}
    \node[nomesistema] at (sistema.north west) {NEXUM Concierge};
    \attore{ap0}{(-4.23,0.00)}{Dipendente}
    \draw[associazione] (ap0) -- (b);
  \end{tikzpicture}
  \end{adjustbox}
  \caption{Estensioni di UC11 -- Domanda sulle regole aziendali}\label{fig:uc:regole:ext}
\end{figure}

\begin{scheda}
Descrizione & Il dipendente chiede come funzionano ferie e permessi nella propria azienda («Come funziona la richiesta di ferie nella nostra azienda?»). La risposta si basa solo sui documenti caricati nel backoffice per quell'azienda. \\ \addlinespace
Attore principale & Dipendente \\ \addlinespace
Attori secondari & Fornitore di modelli AI \\ \addlinespace
Precondizioni & Esiste almeno una fonte disponibile in un bucket di conoscenza aziendale dell'azienda del dipendente. \\ \addlinespace
Postcondizioni & Il dipendente riceve una risposta supportata dalle fonti, con le fonti citate. \\ \addlinespace
Scenario principale & \begin{enumerate}[nosep, leftmargin=*, before=\vspace{-\baselineskip}\vspace{2pt}, after=\vspace{-\baselineskip}]\item Il dipendente pone una domanda su regolamenti o procedure aziendali. \item Il sistema riconosce che serve la conoscenza aziendale. \item Il sistema recupera le porzioni pertinenti, cercando solo tra le fonti attive dei bucket dell'azienda del dipendente. \item Il sistema genera la risposta basandosi esclusivamente sulle porzioni recuperate. \item Il sistema mostra la risposta con le fonti citate (\hyperref[uc:fonti]{UC12}). \end{enumerate} \\ \addlinespace
Inclusioni & \hyperref[uc:fonti]{UC12} \\ \addlinespace
Estensioni & \hyperref[uc:astensione]{UC13} \{fonti insufficienti o in contrasto non risolvibile\} \\ \addlinespace
Generalizzazione & specializza \hyperref[uc:msg]{UC6} \\ \addlinespace
Requisiti & \hyperlink{req:RF-21}{RF-21}, \hyperlink{req:RF-22}{RF-22}, \hyperlink{req:RF-23}{RF-23} \\
\end{scheda}

\subsubsection{UC12 -- Visualizzazione delle fonti citate}\label{uc:fonti}
\begin{figure}[H]
  \centering
  \begin{adjustbox}{max width=\linewidth, max totalheight=0.62\textheight}
  \begin{tikzpicture}
    \node[ucfocus] (b) at (0.00,0.00) {\textbf{UC12}\\Visualizzazione delle fonti citate};
    \node[uc] (c0) at (7.15,1.14) {\textbf{UC11}\\Domanda sulle regole aziendali};
    \draw[dipendenza] (c0) -- node[etich, sloped] {\guillemotleft include\guillemotright} (b);
    \node[uc] (c1) at (7.15,-1.14) {\textbf{UC14}\\Domanda sull'utilizzo di NEXUM};
    \draw[dipendenza] (c1) -- node[etich, sloped] {\guillemotleft include\guillemotright} (b);
    \begin{scope}[on background layer]
      \node[confine, fit=(b) (c0) (c1)] (sistema) {};
    \end{scope}
    \node[nomesistema] at (sistema.north west) {NEXUM Concierge};
    \attore{ap0}{(-4.83,0.00)}{Dipendente}
    \draw[associazione] (ap0) -- (b);
  \end{tikzpicture}
  \end{adjustbox}
  \caption{Diagramma di UC12 -- Visualizzazione delle fonti citate}\label{fig:uc:fonti}
\end{figure}

\begin{scheda}
Descrizione & Ogni risposta documentale indica da dove proviene. \\ \addlinespace
Attore principale & Dipendente \\ \addlinespace
Precondizioni & Il sistema ha generato una risposta a partire da una o più fonti. \\ \addlinespace
Postcondizioni & Il dipendente vede le fonti su cui si basa la risposta. \\ \addlinespace
Scenario principale & Il sistema mostra, per ogni fonte usata: titolo, categoria (conoscenza aziendale o manuale NEXUM) e versione. \\ \addlinespace
Scenari alternativi & 1a. È disponibile il riferimento puntuale: il sistema indica anche la pagina o la sezione utilizzata. \\ \addlinespace
Incluso da & \hyperref[uc:regole]{UC11}, \hyperref[uc:prodotto]{UC14} \\ \addlinespace
Requisiti & \hyperlink{req:RF-24}{RF-24}, \hyperlink{req:RF-100}{RF-100} \\
\end{scheda}

\subsubsection{UC13 -- Astensione dalla risposta}\label{uc:astensione}
\begin{figure}[H]
  \centering
  \begin{adjustbox}{max width=\linewidth, max totalheight=0.62\textheight}
  \begin{tikzpicture}
    \node[ucfocus] (b) at (0.00,0.00) {\textbf{UC13}\\Astensione dalla risposta};
    \node[uc] (c0) at (7.15,1.14) {\textbf{UC11}\\Domanda sulle regole aziendali};
    \draw[dipendenza] (b) -- node[etich, sloped] {\guillemotleft extend\guillemotright} (c0);
    \node[uc] (c1) at (7.15,-1.14) {\textbf{UC14}\\Domanda sull'utilizzo di NEXUM};
    \draw[dipendenza] (b) -- node[etich, sloped] {\guillemotleft extend\guillemotright} (c1);
    \begin{scope}[on background layer]
      \node[confine, fit=(b) (c0) (c1)] (sistema) {};
    \end{scope}
    \node[nomesistema] at (sistema.north west) {NEXUM Concierge};
    \attore{ap0}{(-4.83,0.00)}{Dipendente}
    \draw[associazione] (ap0) -- (b);
  \end{tikzpicture}
  \end{adjustbox}
  \caption{Diagramma di UC13 -- Astensione dalla risposta}\label{fig:uc:astensione}
\end{figure}

\begin{scheda}
Descrizione & Quando le fonti non bastano, il Concierge lo dichiara invece di rispondere comunque. \\ \addlinespace
Attore principale & Dipendente \\ \addlinespace
Precondizioni & Le fonti disponibili non contengono informazioni sufficienti, oppure contengono indicazioni in contrasto che il sistema non può risolvere. \\ \addlinespace
Postcondizioni & Nessuna risposta non supportata è fornita; il dipendente sa perché. \\ \addlinespace
Scenario principale & \begin{enumerate}[nosep, leftmargin=*, before=\vspace{-\baselineskip}\vspace{2pt}, after=\vspace{-\baselineskip}]\item Il sistema dichiara di non disporre di informazioni sufficienti per rispondere. \item Il sistema suggerisce di rivolgersi all'amministrazione dell'azienda. \end{enumerate} \\ \addlinespace
Scenari alternativi & 1a. Le fonti sono in contrasto: il sistema lo dichiara e cita le fonti in contrasto, senza scegliere tra esse. \\ \addlinespace
Estende & \hyperref[uc:regole]{UC11} \{fonti insufficienti o in contrasto non risolvibile\}; \hyperref[uc:prodotto]{UC14} \{manuali insufficienti per rispondere\} \\ \addlinespace
Requisiti & \hyperlink{req:RF-25}{RF-25} \\
\end{scheda}

\subsubsection{UC14 -- Domanda sull'utilizzo di NEXUM}\label{uc:prodotto}
\begin{figure}[H]
  \centering
  \begin{adjustbox}{max width=\linewidth, max totalheight=0.62\textheight}
  \begin{tikzpicture}
    \node[ucfocus] (b) at (0.00,0.00) {\textbf{UC14}\\Domanda sull'utilizzo di NEXUM};
    \node[uc] (c0) at (7.15,0.00) {\textbf{UC12}\\Visualizzazione delle fonti citate};
    \draw[dipendenza] (b) -- node[etich, sloped] {\guillemotleft include\guillemotright} (c0);
    \node[uc] (g) at (0.00,2.84) {\textbf{UC6}\\Invio di un messaggio al Concierge};
    \draw[generalizzazione] (b) -- (g);
    \begin{scope}[on background layer]
      \node[confine, fit=(b) (c0) (g)] (sistema) {};
    \end{scope}
    \node[nomesistema] at (sistema.north west) {NEXUM Concierge};
    \attore{ap0}{(-4.83,0.00)}{Dipendente}
    \draw[associazione] (ap0) -- (b);
    \attore{as0}{(0.00,-2.85)}{Fornitore\\di modelli AI}
    \draw[associazione] (b) -- (as0);
  \end{tikzpicture}
  \end{adjustbox}
  \caption{Diagramma di UC14 -- Domanda sull'utilizzo di NEXUM}\label{fig:uc:prodotto}
\end{figure}

\begin{figure}[H]
  \centering
  \begin{adjustbox}{max width=\linewidth, max totalheight=0.62\textheight}
  \begin{tikzpicture}
    \node[ucfocus] (b) at (0.00,0.00) {\textbf{UC14}\\Domanda sull'utilizzo di NEXUM};
    \node[uc] (e0) at (6.60,0.00) {\textbf{UC13}\\Astensione dalla risposta};
    \draw[dipendenza] (e0) -- (b) node[etich, pos=0.3, sloped] {\guillemotleft extend\guillemotright} coordinate[pos=0.5] (m0);
    \node[nota] (n0) at (3.3,-1.9) {\textbf{Condizione:}\\\{manuali insufficienti per rispondere\}};
    \draw[ancora] (n0) -- (m0);
    \begin{scope}[on background layer]
      \node[confine, fit=(b) (e0) (n0)] (sistema) {};
    \end{scope}
    \node[nomesistema] at (sistema.north west) {NEXUM Concierge};
    \attore{ap0}{(-4.23,0.00)}{Dipendente}
    \draw[associazione] (ap0) -- (b);
  \end{tikzpicture}
  \end{adjustbox}
  \caption{Estensioni di UC14 -- Domanda sull'utilizzo di NEXUM}\label{fig:uc:prodotto:ext}
\end{figure}

\begin{scheda}
Descrizione & Il dipendente chiede come si usa una funzione di NEXUM («Dove posso segnare le ferie?»). La risposta si basa sui manuali NEXUM e spiega come fare, non se una richiesta sia consentita o approvata. \\ \addlinespace
Attore principale & Dipendente \\ \addlinespace
Attori secondari & Fornitore di modelli AI \\ \addlinespace
Precondizioni & I manuali NEXUM pertinenti sono caricati in un bucket di conoscenza del prodotto. \\ \addlinespace
Postcondizioni & Il dipendente riceve le istruzioni con le fonti citate. \\ \addlinespace
Scenario principale & \begin{enumerate}[nosep, leftmargin=*, before=\vspace{-\baselineskip}\vspace{2pt}, after=\vspace{-\baselineskip}]\item Il dipendente chiede come usare una funzione, una schermata o un passaggio di NEXUM. \item Il sistema riconosce che serve la conoscenza del prodotto. \item Il sistema recupera le porzioni pertinenti dai manuali NEXUM. \item Il sistema genera la risposta basandosi su quelle porzioni. \item Il sistema mostra la risposta con le fonti citate (\hyperref[uc:fonti]{UC12}). \end{enumerate} \\ \addlinespace
Scenari alternativi & 1a. La domanda riguarda anche le regole («Posso segnare le ferie ad agosto?»): il sistema distingue nella risposta le istruzioni d'uso dalle regole aziendali, citando le due fonti separatamente. \\ \addlinespace
Inclusioni & \hyperref[uc:fonti]{UC12} \\ \addlinespace
Estensioni & \hyperref[uc:astensione]{UC13} \{manuali insufficienti per rispondere\} \\ \addlinespace
Generalizzazione & specializza \hyperref[uc:msg]{UC6} \\ \addlinespace
Requisiti & \hyperlink{req:RF-28}{RF-28}, \hyperlink{req:RF-29}{RF-29} \\
\end{scheda}

\subsubsection{UC15 -- Consultazione dei propri dati di ferie e permessi}\label{uc:dati}
\begin{figure}[H]
  \centering
  \begin{adjustbox}{max width=\linewidth, max totalheight=0.62\textheight}
  \begin{tikzpicture}
    \node[ucfocus] (b) at (0.00,0.00) {\textbf{UC15}\\Consultazione dei propri dati di ferie e permessi};
    \node[uc] (c0) at (7.15,1.39) {\textbf{UC15.1}\\Visualizzazione delle richieste per periodo e stato};
    \draw[dipendenza] (b) -- node[etich, sloped] {\guillemotleft include\guillemotright} (c0);
    \node[uc] (c1) at (7.15,-1.39) {\textbf{UC15.2}\\Visualizzazione di periodo, unità di misura e criterio di conteggio};
    \draw[dipendenza] (b) -- node[etich, sloped] {\guillemotleft include\guillemotright} (c1);
    \node[uc] (g) at (0.00,3.08) {\textbf{UC6}\\Invio di un messaggio al Concierge};
    \draw[generalizzazione] (b) -- (g);
    \begin{scope}[on background layer]
      \node[confine, fit=(b) (c0) (c1) (g)] (sistema) {};
    \end{scope}
    \node[nomesistema] at (sistema.north west) {NEXUM Concierge};
    \attore{ap0}{(-4.83,0.00)}{Dipendente}
    \draw[associazione] (ap0) -- (b);
    \attore{as0}{(0.00,-4.72)}{Servizio\\ferie e permessi NEXUM}
    \draw[associazione] (b) -- (as0);
  \end{tikzpicture}
  \end{adjustbox}
  \caption{Diagramma di UC15 -- Consultazione dei propri dati di ferie e permessi}\label{fig:uc:dati}
\end{figure}

\begin{figure}[H]
  \centering
  \begin{adjustbox}{max width=\linewidth, max totalheight=0.62\textheight}
  \begin{tikzpicture}
    \node[ucfocus] (b) at (0.00,0.00) {\textbf{UC15}\\Consultazione dei propri dati di ferie e permessi};
    \node[uc] (e0) at (-5.00,4.00) {\textbf{UC16}\\Visualizzazione del saldo residuo};
    \draw[dipendenza] (e0) -- (b) node[etich, pos=0.3, sloped] {\guillemotleft extend\guillemotright} coordinate[pos=0.5] (m0);
    \node[nota] (n0) at (-6.45,1.75) {\textbf{Condizione:}\\\{il servizio fornisce il saldo residuo\}};
    \draw[ancora] (n0) -- (m0);
    \node[uc] (e1) at (5.00,4.00) {\textbf{UC17}\\Nessun dato corrispondente};
    \draw[dipendenza] (e1) -- (b) node[etich, pos=0.3, sloped] {\guillemotleft extend\guillemotright} coordinate[pos=0.5] (m1);
    \node[nota] (n1) at (6.45,1.75) {\textbf{Condizione:}\\\{nessun dato nel periodo richiesto\}};
    \draw[ancora] (n1) -- (m1);
    \node[uc] (e2) at (5.00,-4.00) {\textbf{UC18}\\Rifiuto di una richiesta non autorizzata};
    \draw[dipendenza] (e2) -- (b) node[etich, pos=0.3, sloped] {\guillemotleft extend\guillemotright} coordinate[pos=0.5] (m2);
    \node[nota] (n2) at (6.45,-1.75) {\textbf{Condizione:}\\\{richiesta sui dati di un'altra persona\}};
    \draw[ancora] (n2) -- (m2);
    \begin{scope}[on background layer]
      \node[confine, fit=(b) (e0) (n0) (e1) (n1) (e2) (n2)] (sistema) {};
    \end{scope}
    \node[nomesistema] at (sistema.north west) {NEXUM Concierge};
    \attore{ap0}{(-10.25,0.00)}{Dipendente}
    \draw[associazione] (ap0) -- (b);
  \end{tikzpicture}
  \end{adjustbox}
  \caption{Estensioni di UC15 -- Consultazione dei propri dati di ferie e permessi}\label{fig:uc:dati:ext}
\end{figure}

\begin{scheda}
Descrizione & Il dipendente chiede informazioni sulle proprie ferie e sui propri permessi («Quante ferie ho fatto quest'anno?»). I dati arrivano dall'abilità di lettura, che interroga le API Eggon, autenticate con OAuth, per l'utente NEXUM associato al dipendente; il modello non calcola né inventa valori. \\ \addlinespace
Attore principale & Dipendente \\ \addlinespace
Attori secondari & Servizio ferie e permessi NEXUM \\ \addlinespace
Precondizioni & Le API Eggon sono raggiungibili. \\ \addlinespace
Postcondizioni & Il dipendente vede i propri dati del periodo richiesto, con periodo, unità di misura e criterio di conteggio. \\ \addlinespace
Scenario principale & \begin{enumerate}[nosep, leftmargin=*, before=\vspace{-\baselineskip}\vspace{2pt}, after=\vspace{-\baselineskip}]\item Il dipendente chiede informazioni sulle proprie ferie o sui propri permessi. \item Il sistema riconosce che serve la consultazione dei dati e ricava periodo e tipo dal messaggio o dal contesto. \item Il sistema invoca l'abilità di lettura, che ottiene i dati dalle API Eggon per l'utente NEXUM associato al dipendente. \item Il sistema mostra le richieste del periodo suddivise per stato (\hyperref[uc:dati_riep]{UC15.1}). \item Il sistema indica periodo, unità di misura e criterio di conteggio (\hyperref[uc:dati_criteri]{UC15.2}). \end{enumerate} \\ \addlinespace
Inclusioni & \hyperref[uc:dati_riep]{UC15.1}, \hyperref[uc:dati_criteri]{UC15.2} \\ \addlinespace
Estensioni & \hyperref[uc:saldo]{UC16} \{il servizio fornisce il saldo residuo\}; \hyperref[uc:nessun_dato]{UC17} \{nessun dato nel periodo richiesto\}; \hyperref[uc:non_aut]{UC18} \{richiesta sui dati di un'altra persona\} \\ \addlinespace
Generalizzazione & specializza \hyperref[uc:msg]{UC6} \\ \addlinespace
Requisiti & \hyperlink{req:RF-31}{RF-31}, \hyperlink{req:RF-32}{RF-32}, \hyperlink{req:RF-96}{RF-96} \\
\end{scheda}

\subsubsubsection{UC15.1 -- Visualizzazione delle richieste per periodo e stato}\label{uc:dati_riep}
\begin{figure}[H]
  \centering
  \begin{adjustbox}{max width=\linewidth, max totalheight=0.62\textheight}
  \begin{tikzpicture}
    \node[ucfocus] (b) at (0.00,0.00) {\textbf{UC15.1}\\Visualizzazione delle richieste per periodo e stato};
    \node[uc] (c0) at (7.15,0.00) {\textbf{UC15}\\Consultazione dei propri dati di ferie e permessi};
    \draw[dipendenza] (c0) -- node[etich, sloped] {\guillemotleft include\guillemotright} (b);
    \begin{scope}[on background layer]
      \node[confine, fit=(b) (c0)] (sistema) {};
    \end{scope}
    \node[nomesistema] at (sistema.north west) {NEXUM Concierge};
    \attore{ap0}{(-4.83,0.00)}{Dipendente}
    \draw[associazione] (ap0) -- (b);
  \end{tikzpicture}
  \end{adjustbox}
  \caption{Diagramma di UC15.1 -- Visualizzazione delle richieste per periodo e stato}\label{fig:uc:dati_riep}
\end{figure}

\begin{scheda}
Attore principale & Dipendente \\ \addlinespace
Precondizioni & Il servizio ha restituito i dati del periodo. \\ \addlinespace
Postcondizioni & Le richieste del periodo sono visualizzate per periodo e stato. \\ \addlinespace
Scenario principale & Il sistema mostra le richieste distinguendo il periodo (passate o future) dallo stato (per esempio approvate, in attesa o rifiutate), secondo il significato dei dati concordato con Eggon, con il totale di ciascun gruppo. \\ \addlinespace
Incluso da & \hyperref[uc:dati]{UC15} \\ \addlinespace
Requisiti & \hyperlink{req:RF-33}{RF-33} \\
\end{scheda}

\subsubsubsection{UC15.2 -- Visualizzazione di periodo, unità di misura e criterio di conteggio}\label{uc:dati_criteri}
\begin{figure}[H]
  \centering
  \begin{adjustbox}{max width=\linewidth, max totalheight=0.62\textheight}
  \begin{tikzpicture}
    \node[ucfocus] (b) at (0.00,0.00) {\textbf{UC15.2}\\Visualizzazione di periodo, unità di misura e criterio di conteggio};
    \node[uc] (c0) at (7.15,0.00) {\textbf{UC15}\\Consultazione dei propri dati di ferie e permessi};
    \draw[dipendenza] (c0) -- node[etich, sloped] {\guillemotleft include\guillemotright} (b);
    \begin{scope}[on background layer]
      \node[confine, fit=(b) (c0)] (sistema) {};
    \end{scope}
    \node[nomesistema] at (sistema.north west) {NEXUM Concierge};
    \attore{ap0}{(-4.83,0.00)}{Dipendente}
    \draw[associazione] (ap0) -- (b);
  \end{tikzpicture}
  \end{adjustbox}
  \caption{Diagramma di UC15.2 -- Visualizzazione di periodo, unità di misura e criterio di conteggio}\label{fig:uc:dati_criteri}
\end{figure}

\begin{scheda}
Attore principale & Dipendente \\ \addlinespace
Precondizioni & Il servizio ha restituito i dati del periodo. \\ \addlinespace
Postcondizioni & Il dipendente sa a cosa si riferiscono i valori mostrati. \\ \addlinespace
Scenario principale & Il sistema indica il periodo considerato con date complete, l'unità di misura (giorni od ore) e il significato degli stati conteggiati. \\ \addlinespace
Incluso da & \hyperref[uc:dati]{UC15} \\ \addlinespace
Requisiti & \hyperlink{req:RF-34}{RF-34} \\
\end{scheda}

\subsubsection{UC16 -- Visualizzazione del saldo residuo}\label{uc:saldo}
\begin{figure}[H]
  \centering
  \begin{adjustbox}{max width=\linewidth, max totalheight=0.62\textheight}
  \begin{tikzpicture}
    \node[ucfocus] (b) at (0.00,0.00) {\textbf{UC16}\\Visualizzazione del saldo residuo};
    \node[uc] (c0) at (7.15,0.00) {\textbf{UC15}\\Consultazione dei propri dati di ferie e permessi};
    \draw[dipendenza] (b) -- node[etich, sloped] {\guillemotleft extend\guillemotright} (c0);
    \begin{scope}[on background layer]
      \node[confine, fit=(b) (c0)] (sistema) {};
    \end{scope}
    \node[nomesistema] at (sistema.north west) {NEXUM Concierge};
    \attore{ap0}{(-4.83,0.00)}{Dipendente}
    \draw[associazione] (ap0) -- (b);
  \end{tikzpicture}
  \end{adjustbox}
  \caption{Diagramma di UC16 -- Visualizzazione del saldo residuo}\label{fig:uc:saldo}
\end{figure}

\begin{scheda}
Descrizione & Il saldo è mostrato solo se il servizio lo fornisce: nell'esperienza di Eggon viene aggiornato mensilmente dal consulente del lavoro. La disponibilità del dato va verificata nel contratto delle API. \\ \addlinespace
Attore principale & Dipendente \\ \addlinespace
Precondizioni & \begin{itemize}[nosep, leftmargin=*, before=\vspace{-\baselineskip}\vspace{2pt}, after=\vspace{-\baselineskip}]\item Il dipendente ha chiesto quante ferie o permessi gli restano. \item Il servizio NEXUM fornisce il saldo residuo. \end{itemize} \\ \addlinespace
Postcondizioni & Il dipendente vede il saldo e il riferimento al suo aggiornamento. \\ \addlinespace
Scenario principale & Il sistema mostra il saldo residuo restituito dal servizio, con unità di misura e riferimento all'aggiornamento. \\ \addlinespace
Scenari alternativi & 1a. Il servizio non fornisce il saldo: il sistema dichiara che il dato non è disponibile e non lo ricostruisce dalle richieste. \\ \addlinespace
Estende & \hyperref[uc:dati]{UC15} \{il servizio fornisce il saldo residuo\} \\ \addlinespace
Requisiti & \hyperlink{req:RF-35}{RF-35}, \hyperlink{req:RF-36}{RF-36} \\
\end{scheda}

\subsubsection{UC17 -- Nessun dato corrispondente}\label{uc:nessun_dato}
\begin{figure}[H]
  \centering
  \begin{adjustbox}{max width=\linewidth, max totalheight=0.62\textheight}
  \begin{tikzpicture}
    \node[ucfocus] (b) at (0.00,0.00) {\textbf{UC17}\\Nessun dato corrispondente};
    \node[uc] (c0) at (7.15,1.27) {\textbf{UC15}\\Consultazione dei propri dati di ferie e permessi};
    \draw[dipendenza] (b) -- node[etich, sloped] {\guillemotleft extend\guillemotright} (c0);
    \node[uc] (c1) at (7.15,-1.27) {\textbf{UC19}\\Consultazione dello stato di una richiesta};
    \draw[dipendenza] (b) -- node[etich, sloped] {\guillemotleft extend\guillemotright} (c1);
    \begin{scope}[on background layer]
      \node[confine, fit=(b) (c0) (c1)] (sistema) {};
    \end{scope}
    \node[nomesistema] at (sistema.north west) {NEXUM Concierge};
    \attore{ap0}{(-4.83,0.00)}{Dipendente}
    \draw[associazione] (ap0) -- (b);
  \end{tikzpicture}
  \end{adjustbox}
  \caption{Diagramma di UC17 -- Nessun dato corrispondente}\label{fig:uc:nessun_dato}
\end{figure}

\begin{scheda}
Attore principale & Dipendente \\ \addlinespace
Precondizioni & Il servizio non restituisce dati o richieste per il periodo o i criteri indicati. \\ \addlinespace
Postcondizioni & Il dipendente sa che non risultano dati, distinto da un errore. \\ \addlinespace
Scenario principale & Il sistema dichiara che non risultano richieste o dati per il periodo e i criteri indicati, riportandoli. \\ \addlinespace
Estende & \hyperref[uc:dati]{UC15} \{nessun dato nel periodo richiesto\}; \hyperref[uc:stato_rich]{UC19} \{nessuna richiesta corrispondente\} \\ \addlinespace
Requisiti & \hyperlink{req:RF-37}{RF-37} \\
\end{scheda}

\subsubsection{UC18 -- Rifiuto di una richiesta non autorizzata}\label{uc:non_aut}
\begin{figure}[H]
  \centering
  \begin{adjustbox}{max width=\linewidth, max totalheight=0.62\textheight}
  \begin{tikzpicture}
    \node[ucfocus] (b) at (0.00,0.00) {\textbf{UC18}\\Rifiuto di una richiesta non autorizzata};
    \node[uc] (c0) at (7.15,1.27) {\textbf{UC15}\\Consultazione dei propri dati di ferie e permessi};
    \draw[dipendenza] (b) -- node[etich, sloped] {\guillemotleft extend\guillemotright} (c0);
    \node[uc] (c1) at (7.15,-1.27) {\textbf{UC20}\\Registrazione di una richiesta di ferie o permesso};
    \draw[dipendenza] (b) -- node[etich, sloped] {\guillemotleft extend\guillemotright} (c1);
    \begin{scope}[on background layer]
      \node[confine, fit=(b) (c0) (c1)] (sistema) {};
    \end{scope}
    \node[nomesistema] at (sistema.north west) {NEXUM Concierge};
    \attore{ap0}{(-4.83,0.00)}{Dipendente}
    \draw[associazione] (ap0) -- (b);
  \end{tikzpicture}
  \end{adjustbox}
  \caption{Diagramma di UC18 -- Rifiuto di una richiesta non autorizzata}\label{fig:uc:non_aut}
\end{figure}

\begin{scheda}
Descrizione & Il testo della conversazione non attribuisce privilegi: le autorizzazioni sono verificate dai servizi applicativi e dalle API Eggon. \\ \addlinespace
Attore principale & Dipendente \\ \addlinespace
Precondizioni & La richiesta riguarda dati o operazioni di un'altra persona o di un'azienda diversa dalla propria. \\ \addlinespace
Postcondizioni & Nessun dato altrui è mostrato e nessuna operazione è inviata. \\ \addlinespace
Scenario principale & \begin{enumerate}[nosep, leftmargin=*, before=\vspace{-\baselineskip}\vspace{2pt}, after=\vspace{-\baselineskip}]\item Il sistema dichiara di poter operare solo sui dati del dipendente nella sua azienda. \item Il sistema registra il tentativo nella traccia. \end{enumerate} \\ \addlinespace
Estende & \hyperref[uc:dati]{UC15} \{richiesta sui dati di un'altra persona\}; \hyperref[uc:registra]{UC20} \{richiesta per conto di un'altra persona\} \\ \addlinespace
Requisiti & \hyperlink{req:RF-38}{RF-38} \\
\end{scheda}

\subsubsection{UC19 -- Consultazione dello stato di una richiesta}\label{uc:stato_rich}
\begin{figure}[H]
  \centering
  \begin{adjustbox}{max width=\linewidth, max totalheight=0.62\textheight}
  \begin{tikzpicture}
    \node[ucfocus] (b) at (0.00,0.00) {\textbf{UC19}\\Consultazione dello stato di una richiesta};
    \node[uc] (c0) at (7.15,0.00) {\textbf{UC30}\\Individuazione della richiesta interessata};
    \draw[dipendenza] (b) -- node[etich, sloped] {\guillemotleft include\guillemotright} (c0);
    \node[uc] (g) at (0.00,2.84) {\textbf{UC6}\\Invio di un messaggio al Concierge};
    \draw[generalizzazione] (b) -- (g);
    \begin{scope}[on background layer]
      \node[confine, fit=(b) (c0) (g)] (sistema) {};
    \end{scope}
    \node[nomesistema] at (sistema.north west) {NEXUM Concierge};
    \attore{ap0}{(-4.83,0.00)}{Dipendente}
    \draw[associazione] (ap0) -- (b);
    \attore{as0}{(0.00,-3.09)}{Servizio\\ferie e permessi NEXUM}
    \draw[associazione] (b) -- (as0);
  \end{tikzpicture}
  \end{adjustbox}
  \caption{Diagramma di UC19 -- Consultazione dello stato di una richiesta}\label{fig:uc:stato_rich}
\end{figure}

\begin{figure}[H]
  \centering
  \begin{adjustbox}{max width=\linewidth, max totalheight=0.62\textheight}
  \begin{tikzpicture}
    \node[ucfocus] (b) at (0.00,0.00) {\textbf{UC19}\\Consultazione dello stato di una richiesta};
    \node[uc] (e0) at (6.60,0.00) {\textbf{UC17}\\Nessun dato corrispondente};
    \draw[dipendenza] (e0) -- (b) node[etich, pos=0.3, sloped] {\guillemotleft extend\guillemotright} coordinate[pos=0.5] (m0);
    \node[nota] (n0) at (3.3,-1.9) {\textbf{Condizione:}\\\{nessuna richiesta corrispondente\}};
    \draw[ancora] (n0) -- (m0);
    \begin{scope}[on background layer]
      \node[confine, fit=(b) (e0) (n0)] (sistema) {};
    \end{scope}
    \node[nomesistema] at (sistema.north west) {NEXUM Concierge};
    \attore{ap0}{(-4.23,0.00)}{Dipendente}
    \draw[associazione] (ap0) -- (b);
  \end{tikzpicture}
  \end{adjustbox}
  \caption{Estensioni di UC19 -- Consultazione dello stato di una richiesta}\label{fig:uc:stato_rich:ext}
\end{figure}

\begin{scheda}
Descrizione & Il dipendente chiede a che punto è una propria richiesta («Le ferie di agosto sono state approvate?»). Lo stato è quello effettivo in NEXUM, anche a distanza di tempo dall'inserimento: con l'approvazione automatica la richiesta può risultare subito approvata, con quella manuale resta in attesa finché un operatore non decide. \\ \addlinespace
Attore principale & Dipendente \\ \addlinespace
Attori secondari & Servizio ferie e permessi NEXUM \\ \addlinespace
Precondizioni & Le API Eggon sono raggiungibili. \\ \addlinespace
Postcondizioni & Il dipendente conosce lo stato aggiornato della richiesta. \\ \addlinespace
Scenario principale & \begin{enumerate}[nosep, leftmargin=*, before=\vspace{-\baselineskip}\vspace{2pt}, after=\vspace{-\baselineskip}]\item Il dipendente chiede lo stato di una propria richiesta. \item Il sistema individua la richiesta interessata (\hyperref[uc:sel_rich]{UC30}). \item Il sistema ottiene lo stato dalle API Eggon. \item Il sistema mostra lo stato (registrata, in attesa di approvazione, approvata o rifiutata) con la data dell'ultimo cambiamento. \end{enumerate} \\ \addlinespace
Inclusioni & \hyperref[uc:sel_rich]{UC30} \\ \addlinespace
Estensioni & \hyperref[uc:nessun_dato]{UC17} \{nessuna richiesta corrispondente\} \\ \addlinespace
Generalizzazione & specializza \hyperref[uc:msg]{UC6} \\ \addlinespace
Requisiti & \hyperlink{req:RF-39}{RF-39}, \hyperlink{req:RF-96}{RF-96} \\
\end{scheda}

\subsubsection{UC20 -- Registrazione di una richiesta di ferie o permesso}\label{uc:registra}
\begin{figure}[H]
  \centering
  \begin{adjustbox}{max width=\linewidth, max totalheight=0.62\textheight}
  \begin{tikzpicture}
    \node[ucfocus] (b) at (0.00,0.00) {\textbf{UC20}\\Registrazione di una richiesta di ferie o permesso};
    \node[uc] (c0) at (7.15,5.72) {\textbf{UC20.1}\\Determinazione del tipo di richiesta};
    \draw[dipendenza] (b) -- node[etich, sloped] {\guillemotleft include\guillemotright} (c0);
    \node[uc] (c1) at (7.15,3.68) {\textbf{UC21}\\Determinazione del periodo};
    \draw[dipendenza] (b) -- node[etich, sloped] {\guillemotleft include\guillemotright} (c1);
    \node[uc] (c2) at (7.15,1.39) {\textbf{UC20.2}\\Raccolta degli altri dati richiesti};
    \draw[dipendenza] (b) -- node[etich, sloped] {\guillemotleft include\guillemotright} (c2);
    \node[uc] (c3) at (7.15,-1.14) {\textbf{UC22}\\Visualizzazione del riepilogo dell'operazione};
    \draw[dipendenza] (b) -- node[etich, sloped] {\guillemotleft include\guillemotright} (c3);
    \node[uc] (c4) at (7.15,-3.43) {\textbf{UC23}\\Conferma esplicita dell'operazione};
    \draw[dipendenza] (b) -- node[etich, sloped] {\guillemotleft include\guillemotright} (c4);
    \node[uc] (c5) at (7.15,-5.72) {\textbf{UC24}\\Visualizzazione dell'esito dell'operazione};
    \draw[dipendenza] (b) -- node[etich, sloped] {\guillemotleft include\guillemotright} (c5);
    \node[uc] (g) at (0.00,3.08) {\textbf{UC6}\\Invio di un messaggio al Concierge};
    \draw[generalizzazione] (b) -- (g);
    \begin{scope}[on background layer]
      \node[confine, fit=(b) (c0) (c1) (c2) (c3) (c4) (c5) (g)] (sistema) {};
    \end{scope}
    \node[nomesistema] at (sistema.north west) {NEXUM Concierge};
    \attore{ap0}{(-4.83,0.00)}{Dipendente}
    \draw[associazione] (ap0) -- (b);
    \attore{as0}{(0.00,-8.82)}{Servizio\\ferie e permessi NEXUM}
    \draw[associazione] (b) -- (as0);
  \end{tikzpicture}
  \end{adjustbox}
  \caption{Diagramma di UC20 -- Registrazione di una richiesta di ferie o permesso}\label{fig:uc:registra}
\end{figure}

\begin{figure}[H]
  \centering
  \begin{adjustbox}{max width=\linewidth, max totalheight=0.62\textheight}
  \begin{tikzpicture}
    \node[ucfocus] (b) at (0.00,0.00) {\textbf{UC20}\\Registrazione di una richiesta di ferie o permesso};
    \node[uc] (e0) at (-5.00,4.00) {\textbf{UC7}\\Richiesta di chiarimento};
    \draw[dipendenza] (e0) -- (b) node[etich, pos=0.3, sloped] {\guillemotleft extend\guillemotright} coordinate[pos=0.5] (m0);
    \node[nota] (n0) at (-6.45,1.75) {\textbf{Condizione:}\\\{dati ambigui o mancanti\}};
    \draw[ancora] (n0) -- (m0);
    \node[uc] (e1) at (0.00,6.40) {\textbf{UC25}\\Modifica dei dati prima della conferma};
    \draw[dipendenza] (e1) -- (b) node[etich, pos=0.3, sloped] {\guillemotleft extend\guillemotright} coordinate[pos=0.5] (m1);
    \node[nota] (n1) at (-3.75,6.5) {\textbf{Condizione:}\\\{il dipendente cambia i dati dopo il riepilogo\}};
    \draw[ancora] (n1) -- (m1);
    \node[uc] (e2) at (5.00,4.00) {\textbf{UC26}\\Rinuncia all'operazione prima della conferma};
    \draw[dipendenza] (e2) -- (b) node[etich, pos=0.3, sloped] {\guillemotleft extend\guillemotright} coordinate[pos=0.5] (m2);
    \node[nota] (n2) at (6.45,1.75) {\textbf{Condizione:}\\\{il dipendente rinuncia prima della conferma\}};
    \draw[ancora] (n2) -- (m2);
    \node[uc] (e3) at (5.00,-4.00) {\textbf{UC27}\\Operazione respinta dal servizio};
    \draw[dipendenza] (e3) -- (b) node[etich, pos=0.3, sloped] {\guillemotleft extend\guillemotright} coordinate[pos=0.5] (m3);
    \node[nota] (n3) at (6.45,-1.75) {\textbf{Condizione:}\\\{il servizio respinge la validazione\}};
    \draw[ancora] (n3) -- (m3);
    \node[uc] (e4) at (0.00,-6.40) {\textbf{UC28}\\Esito dell'operazione incerto};
    \draw[dipendenza] (e4) -- (b) node[etich, pos=0.3, sloped] {\guillemotleft extend\guillemotright} coordinate[pos=0.5] (m4);
    \node[nota] (n4) at (3.75,-6.5) {\textbf{Condizione:}\\\{invio interrotto con esito sconosciuto\}};
    \draw[ancora] (n4) -- (m4);
    \node[uc] (e5) at (-5.00,-4.00) {\textbf{UC18}\\Rifiuto di una richiesta non autorizzata};
    \draw[dipendenza] (e5) -- (b) node[etich, pos=0.3, sloped] {\guillemotleft extend\guillemotright} coordinate[pos=0.5] (m5);
    \node[nota] (n5) at (-6.45,-1.75) {\textbf{Condizione:}\\\{richiesta per conto di un'altra persona\}};
    \draw[ancora] (n5) -- (m5);
    \begin{scope}[on background layer]
      \node[confine, fit=(b) (e0) (n0) (e1) (n1) (e2) (n2) (e3) (n3) (e4) (n4) (e5) (n5)] (sistema) {};
    \end{scope}
    \node[nomesistema] at (sistema.north west) {NEXUM Concierge};
    \attore{ap0}{(-10.25,0.00)}{Dipendente}
    \draw[associazione] (ap0) -- (b);
  \end{tikzpicture}
  \end{adjustbox}
  \caption{Estensioni di UC20 -- Registrazione di una richiesta di ferie o permesso}\label{fig:uc:registra:ext}
\end{figure}

\begin{scheda}
Descrizione & Il dipendente chiede di registrare ferie o un permesso («Registra che sarò in ferie dal 10 al 15 ottobre»). Il messaggio avvia la preparazione della richiesta: l'invio avviene solo dopo un riepilogo e una conferma esplicita, e la registrazione non equivale all'approvazione. \\ \addlinespace
Attore principale & Dipendente \\ \addlinespace
Attori secondari & Servizio ferie e permessi NEXUM \\ \addlinespace
Precondizioni & Le API Eggon sono raggiungibili. \\ \addlinespace
Postcondizioni & \begin{itemize}[nosep, leftmargin=*, before=\vspace{-\baselineskip}\vspace{2pt}, after=\vspace{-\baselineskip}]\item La richiesta è registrata in NEXUM e il dipendente ne conosce l'esito effettivo. \item La traccia riporta persona, azienda, data e ora, abilità invocata, conferma ed esito. \end{itemize} \\ \addlinespace
Scenario principale & \begin{enumerate}[nosep, leftmargin=*, before=\vspace{-\baselineskip}\vspace{2pt}, after=\vspace{-\baselineskip}]\item Il dipendente chiede di registrare ferie o un permesso. \item Il sistema riconosce l'operazione e avvia la preparazione della richiesta, senza inviarla. \item Il sistema determina il tipo di richiesta (\hyperref[uc:reg_tipo]{UC20.1}). \item Il sistema determina il periodo (\hyperref[uc:periodo]{UC21}). \item Il sistema raccoglie gli altri dati richiesti dal servizio (\hyperref[uc:reg_campi]{UC20.2}). \item Il sistema mostra il riepilogo (\hyperref[uc:op_riep]{UC22}). \item Il dipendente conferma esplicitamente l'invio (\hyperref[uc:op_conf]{UC23}). \item Il sistema invia la richiesta tramite l'abilità di registrazione, con una chiamata autenticata alle API Eggon. \item Il sistema mostra l'esito restituito dal servizio (\hyperref[uc:op_esito]{UC24}). \end{enumerate} \\ \addlinespace
Inclusioni & \hyperref[uc:reg_tipo]{UC20.1}, \hyperref[uc:periodo]{UC21}, \hyperref[uc:reg_campi]{UC20.2}, \hyperref[uc:op_riep]{UC22}, \hyperref[uc:op_conf]{UC23}, \hyperref[uc:op_esito]{UC24} \\ \addlinespace
Estensioni & \hyperref[uc:chiarimento]{UC7} \{dati ambigui o mancanti\}; \hyperref[uc:op_modifica]{UC25} \{il dipendente cambia i dati dopo il riepilogo\}; \hyperref[uc:op_rinuncia]{UC26} \{il dipendente rinuncia prima della conferma\}; \hyperref[uc:op_rifiuto]{UC27} \{il servizio respinge la validazione\}; \hyperref[uc:op_incerto]{UC28} \{invio interrotto con esito sconosciuto\}; \hyperref[uc:non_aut]{UC18} \{richiesta per conto di un'altra persona\} \\ \addlinespace
Generalizzazione & specializza \hyperref[uc:msg]{UC6} \\ \addlinespace
Requisiti & \hyperlink{req:RF-40}{RF-40}, \hyperlink{req:RF-41}{RF-41}, \hyperlink{req:RF-42}{RF-42}, \hyperlink{req:RF-43}{RF-43}, \hyperlink{req:RF-96}{RF-96} \\
\end{scheda}

\subsubsubsection{UC20.1 -- Determinazione del tipo di richiesta}\label{uc:reg_tipo}
\begin{figure}[H]
  \centering
  \begin{adjustbox}{max width=\linewidth, max totalheight=0.62\textheight}
  \begin{tikzpicture}
    \node[ucfocus] (b) at (0.00,0.00) {\textbf{UC20.1}\\Determinazione del tipo di richiesta};
    \node[uc] (c0) at (7.15,0.00) {\textbf{UC20}\\Registrazione di una richiesta di ferie o permesso};
    \draw[dipendenza] (c0) -- node[etich, sloped] {\guillemotleft include\guillemotright} (b);
    \begin{scope}[on background layer]
      \node[confine, fit=(b) (c0)] (sistema) {};
    \end{scope}
    \node[nomesistema] at (sistema.north west) {NEXUM Concierge};
    \attore{ap0}{(-4.83,0.00)}{Dipendente}
    \draw[associazione] (ap0) -- (b);
  \end{tikzpicture}
  \end{adjustbox}
  \caption{Diagramma di UC20.1 -- Determinazione del tipo di richiesta}\label{fig:uc:reg_tipo}
\end{figure}

\begin{scheda}
Attore principale & Dipendente \\ \addlinespace
Precondizioni & È in preparazione una richiesta di registrazione. \\ \addlinespace
Postcondizioni & Il tipo di richiesta è determinato. \\ \addlinespace
Scenario principale & Il sistema ricava il tipo di richiesta (ferie o uno dei tipi di permesso previsti dal servizio) dal messaggio o dal contesto. \\ \addlinespace
Scenari alternativi & 1a. Il tipo non è ricavabile: il sistema propone i tipi disponibili (\hyperref[uc:chiarimento]{UC7}). \\ \addlinespace
Incluso da & \hyperref[uc:registra]{UC20} \\ \addlinespace
Requisiti & \hyperlink{req:RF-44}{RF-44} \\
\end{scheda}

\subsubsubsection{UC20.2 -- Raccolta degli altri dati richiesti}\label{uc:reg_campi}
\begin{figure}[H]
  \centering
  \begin{adjustbox}{max width=\linewidth, max totalheight=0.62\textheight}
  \begin{tikzpicture}
    \node[ucfocus] (b) at (0.00,0.00) {\textbf{UC20.2}\\Raccolta degli altri dati richiesti};
    \node[uc] (c0) at (7.15,0.00) {\textbf{UC20}\\Registrazione di una richiesta di ferie o permesso};
    \draw[dipendenza] (c0) -- node[etich, sloped] {\guillemotleft include\guillemotright} (b);
    \begin{scope}[on background layer]
      \node[confine, fit=(b) (c0)] (sistema) {};
    \end{scope}
    \node[nomesistema] at (sistema.north west) {NEXUM Concierge};
    \attore{ap0}{(-4.83,0.00)}{Dipendente}
    \draw[associazione] (ap0) -- (b);
  \end{tikzpicture}
  \end{adjustbox}
  \caption{Diagramma di UC20.2 -- Raccolta degli altri dati richiesti}\label{fig:uc:reg_campi}
\end{figure}

\begin{scheda}
Attore principale & Dipendente \\ \addlinespace
Precondizioni & Tipo e periodo della richiesta sono determinati. \\ \addlinespace
Postcondizioni & Tutti i campi obbligatori previsti dal servizio sono valorizzati. \\ \addlinespace
Scenario principale & \begin{enumerate}[nosep, leftmargin=*, before=\vspace{-\baselineskip}\vspace{2pt}, after=\vspace{-\baselineskip}]\item Il sistema chiede i campi obbligatori previsti dal servizio per quel tipo di richiesta (per esempio una nota) che non sono già noti. \item Il dipendente li fornisce. \end{enumerate} \\ \addlinespace
Incluso da & \hyperref[uc:registra]{UC20} \\ \addlinespace
Requisiti & \hyperlink{req:RF-46}{RF-46} \\
\end{scheda}

\subsubsection{UC21 -- Determinazione del periodo}\label{uc:periodo}
\begin{figure}[H]
  \centering
  \begin{adjustbox}{max width=\linewidth, max totalheight=0.62\textheight}
  \begin{tikzpicture}
    \node[ucfocus] (b) at (0.00,0.00) {\textbf{UC21}\\Determinazione del periodo};
    \node[uc] (c0) at (7.15,0.00) {\textbf{UC20}\\Registrazione di una richiesta di ferie o permesso};
    \draw[dipendenza] (c0) -- node[etich, sloped] {\guillemotleft include\guillemotright} (b);
    \begin{scope}[on background layer]
      \node[confine, fit=(b) (c0)] (sistema) {};
    \end{scope}
    \node[nomesistema] at (sistema.north west) {NEXUM Concierge};
    \attore{ap0}{(-4.83,0.00)}{Dipendente}
    \draw[associazione] (ap0) -- (b);
  \end{tikzpicture}
  \end{adjustbox}
  \caption{Diagramma di UC21 -- Determinazione del periodo}\label{fig:uc:periodo}
\end{figure}

\begin{scheda}
Descrizione & Caso d'uso incluso dalle operazioni che richiedono un periodo. \\ \addlinespace
Attore principale & Dipendente \\ \addlinespace
Precondizioni & È in preparazione una richiesta che richiede un periodo. \\ \addlinespace
Postcondizioni & Date di inizio e fine, ed eventuali orari, sono determinati e completi di anno. \\ \addlinespace
Scenario principale & \begin{enumerate}[nosep, leftmargin=*, before=\vspace{-\baselineskip}\vspace{2pt}, after=\vspace{-\baselineskip}]\item Il sistema ricava dal messaggio o dal contesto le date di inizio e fine e, per i permessi, gli orari o la durata. \item Il sistema completa le date prive di anno con l'anno già esplicitato nella conversazione. \end{enumerate} \\ \addlinespace
Scenari alternativi & \begin{itemize}[nosep, leftmargin=*, before=\vspace{-\baselineskip}\vspace{2pt}, after=\vspace{-\baselineskip}]\item 2a. L'anno non è ricavabile dal contesto: il sistema lo chiede (\hyperref[uc:chiarimento]{UC7}). \item 1a. Il periodo è incoerente (fine precedente all'inizio): il sistema lo segnala e chiede di correggerlo. \end{itemize} \\ \addlinespace
Incluso da & \hyperref[uc:registra]{UC20} \\ \addlinespace
Requisiti & \hyperlink{req:RF-45}{RF-45} \\
\end{scheda}

\subsubsection{UC22 -- Visualizzazione del riepilogo dell'operazione}\label{uc:op_riep}
\begin{figure}[H]
  \centering
  \begin{adjustbox}{max width=\linewidth, max totalheight=0.62\textheight}
  \begin{tikzpicture}
    \node[ucfocus] (b) at (0.00,0.00) {\textbf{UC22}\\Visualizzazione del riepilogo dell'operazione};
    \node[uc] (c0) at (7.15,1.14) {\textbf{UC20}\\Registrazione di una richiesta di ferie o permesso};
    \draw[dipendenza] (c0) -- node[etich, sloped] {\guillemotleft include\guillemotright} (b);
    \node[uc] (c1) at (7.15,-1.14) {\textbf{UC29}\\Cancellazione di una richiesta};
    \draw[dipendenza] (c1) -- node[etich, sloped] {\guillemotleft include\guillemotright} (b);
    \begin{scope}[on background layer]
      \node[confine, fit=(b) (c0) (c1)] (sistema) {};
    \end{scope}
    \node[nomesistema] at (sistema.north west) {NEXUM Concierge};
    \attore{ap0}{(-4.83,0.00)}{Dipendente}
    \draw[associazione] (ap0) -- (b);
  \end{tikzpicture}
  \end{adjustbox}
  \caption{Diagramma di UC22 -- Visualizzazione del riepilogo dell'operazione}\label{fig:uc:op_riep}
\end{figure}

\begin{scheda}
Descrizione & Caso d'uso incluso da ogni operazione che modifica dati in NEXUM. \\ \addlinespace
Attore principale & Dipendente \\ \addlinespace
Precondizioni & Tutti i dati dell'operazione sono noti. \\ \addlinespace
Postcondizioni & Il dipendente vede esattamente cosa verrà inviato. \\ \addlinespace
Scenario principale & \begin{enumerate}[nosep, leftmargin=*, before=\vspace{-\baselineskip}\vspace{2pt}, after=\vspace{-\baselineskip}]\item Il sistema mostra un riepilogo completo: operazione, tipo, date di inizio e fine complete di giorno, mese e anno, orari, durata e altri campi. \item Il sistema precisa che la richiesta non è ancora stata inviata e chiede la conferma. \end{enumerate} \\ \addlinespace
Incluso da & \hyperref[uc:registra]{UC20}, \hyperref[uc:annulla_inv]{UC29} \\ \addlinespace
Requisiti & \hyperlink{req:RF-47}{RF-47} \\
\end{scheda}

\subsubsection{UC23 -- Conferma esplicita dell'operazione}\label{uc:op_conf}
\begin{figure}[H]
  \centering
  \begin{adjustbox}{max width=\linewidth, max totalheight=0.62\textheight}
  \begin{tikzpicture}
    \node[ucfocus] (b) at (0.00,0.00) {\textbf{UC23}\\Conferma esplicita dell'operazione};
    \node[uc] (c0) at (7.15,1.14) {\textbf{UC20}\\Registrazione di una richiesta di ferie o permesso};
    \draw[dipendenza] (c0) -- node[etich, sloped] {\guillemotleft include\guillemotright} (b);
    \node[uc] (c1) at (7.15,-1.14) {\textbf{UC29}\\Cancellazione di una richiesta};
    \draw[dipendenza] (c1) -- node[etich, sloped] {\guillemotleft include\guillemotright} (b);
    \begin{scope}[on background layer]
      \node[confine, fit=(b) (c0) (c1)] (sistema) {};
    \end{scope}
    \node[nomesistema] at (sistema.north west) {NEXUM Concierge};
    \attore{ap0}{(-4.83,0.00)}{Dipendente}
    \draw[associazione] (ap0) -- (b);
  \end{tikzpicture}
  \end{adjustbox}
  \caption{Diagramma di UC23 -- Conferma esplicita dell'operazione}\label{fig:uc:op_conf}
\end{figure}

\begin{scheda}
Attore principale & Dipendente \\ \addlinespace
Precondizioni & Il riepilogo dell'operazione è visualizzato. \\ \addlinespace
Postcondizioni & L'operazione è confermata e può essere inviata. \\ \addlinespace
Scenario principale & Il dipendente conferma esplicitamente, con il comando di conferma o una risposta affermativa inequivocabile. \\ \addlinespace
Scenari alternativi & 1a. La risposta è ambigua: il sistema non la considera una conferma e ripete la domanda. \\ \addlinespace
Incluso da & \hyperref[uc:registra]{UC20}, \hyperref[uc:annulla_inv]{UC29} \\ \addlinespace
Requisiti & \hyperlink{req:RF-48}{RF-48} \\
\end{scheda}

\subsubsection{UC24 -- Visualizzazione dell'esito dell'operazione}\label{uc:op_esito}
\begin{figure}[H]
  \centering
  \begin{adjustbox}{max width=\linewidth, max totalheight=0.62\textheight}
  \begin{tikzpicture}
    \node[ucfocus] (b) at (0.00,0.00) {\textbf{UC24}\\Visualizzazione dell'esito dell'operazione};
    \node[uc] (c0) at (7.15,1.14) {\textbf{UC20}\\Registrazione di una richiesta di ferie o permesso};
    \draw[dipendenza] (c0) -- node[etich, sloped] {\guillemotleft include\guillemotright} (b);
    \node[uc] (c1) at (7.15,-1.14) {\textbf{UC29}\\Cancellazione di una richiesta};
    \draw[dipendenza] (c1) -- node[etich, sloped] {\guillemotleft include\guillemotright} (b);
    \begin{scope}[on background layer]
      \node[confine, fit=(b) (c0) (c1)] (sistema) {};
    \end{scope}
    \node[nomesistema] at (sistema.north west) {NEXUM Concierge};
    \attore{ap0}{(-4.83,0.00)}{Dipendente}
    \draw[associazione] (ap0) -- (b);
  \end{tikzpicture}
  \end{adjustbox}
  \caption{Diagramma di UC24 -- Visualizzazione dell'esito dell'operazione}\label{fig:uc:op_esito}
\end{figure}

\begin{scheda}
Attore principale & Dipendente \\ \addlinespace
Precondizioni & Il servizio ha risposto all'invio. \\ \addlinespace
Postcondizioni & Il dipendente conosce l'esito effettivo. \\ \addlinespace
Scenario principale & \begin{enumerate}[nosep, leftmargin=*, before=\vspace{-\baselineskip}\vspace{2pt}, after=\vspace{-\baselineskip}]\item Il sistema mostra l'esito restituito dal servizio: identificativo della richiesta e stato (registrata o in attesa di approvazione). \item Il sistema precisa che la registrazione non equivale all'approvazione, che resta nel processo NEXUM. \end{enumerate} \\ \addlinespace
Incluso da & \hyperref[uc:registra]{UC20}, \hyperref[uc:annulla_inv]{UC29} \\ \addlinespace
Requisiti & \hyperlink{req:RF-49}{RF-49} \\
\end{scheda}

\subsubsection{UC25 -- Modifica dei dati prima della conferma}\label{uc:op_modifica}
\begin{figure}[H]
  \centering
  \begin{adjustbox}{max width=\linewidth, max totalheight=0.62\textheight}
  \begin{tikzpicture}
    \node[ucfocus] (b) at (0.00,0.00) {\textbf{UC25}\\Modifica dei dati prima della conferma};
    \node[uc] (c0) at (7.15,0.00) {\textbf{UC20}\\Registrazione di una richiesta di ferie o permesso};
    \draw[dipendenza] (b) -- node[etich, sloped] {\guillemotleft extend\guillemotright} (c0);
    \begin{scope}[on background layer]
      \node[confine, fit=(b) (c0)] (sistema) {};
    \end{scope}
    \node[nomesistema] at (sistema.north west) {NEXUM Concierge};
    \attore{ap0}{(-4.83,0.00)}{Dipendente}
    \draw[associazione] (ap0) -- (b);
  \end{tikzpicture}
  \end{adjustbox}
  \caption{Diagramma di UC25 -- Modifica dei dati prima della conferma}\label{fig:uc:op_modifica}
\end{figure}

\begin{scheda}
Attore principale & Dipendente \\ \addlinespace
Precondizioni & Il riepilogo è visualizzato e il dipendente chiede di cambiare un dato («fai fino al 16»). \\ \addlinespace
Postcondizioni & È visualizzato un nuovo riepilogo; la conferma precedente non è più valida. \\ \addlinespace
Scenario principale & \begin{enumerate}[nosep, leftmargin=*, before=\vspace{-\baselineskip}\vspace{2pt}, after=\vspace{-\baselineskip}]\item Il sistema aggiorna i dati. \item Il sistema mostra un nuovo riepilogo e chiede una nuova conferma. \end{enumerate} \\ \addlinespace
Estende & \hyperref[uc:registra]{UC20} \{il dipendente cambia i dati dopo il riepilogo\} \\ \addlinespace
Requisiti & \hyperlink{req:RF-50}{RF-50} \\
\end{scheda}

\subsubsection{UC26 -- Rinuncia all'operazione prima della conferma}\label{uc:op_rinuncia}
\begin{figure}[H]
  \centering
  \begin{adjustbox}{max width=\linewidth, max totalheight=0.62\textheight}
  \begin{tikzpicture}
    \node[ucfocus] (b) at (0.00,0.00) {\textbf{UC26}\\Rinuncia all'operazione prima della conferma};
    \node[uc] (c0) at (7.15,1.14) {\textbf{UC20}\\Registrazione di una richiesta di ferie o permesso};
    \draw[dipendenza] (b) -- node[etich, sloped] {\guillemotleft extend\guillemotright} (c0);
    \node[uc] (c1) at (7.15,-1.14) {\textbf{UC29}\\Cancellazione di una richiesta};
    \draw[dipendenza] (b) -- node[etich, sloped] {\guillemotleft extend\guillemotright} (c1);
    \begin{scope}[on background layer]
      \node[confine, fit=(b) (c0) (c1)] (sistema) {};
    \end{scope}
    \node[nomesistema] at (sistema.north west) {NEXUM Concierge};
    \attore{ap0}{(-4.83,0.00)}{Dipendente}
    \draw[associazione] (ap0) -- (b);
  \end{tikzpicture}
  \end{adjustbox}
  \caption{Diagramma di UC26 -- Rinuncia all'operazione prima della conferma}\label{fig:uc:op_rinuncia}
\end{figure}

\begin{scheda}
Attore principale & Dipendente \\ \addlinespace
Precondizioni & Il riepilogo è visualizzato e il dipendente rinuncia. \\ \addlinespace
Postcondizioni & Nulla è inviato e la preparazione è chiusa. \\ \addlinespace
Scenario principale & Il sistema conferma che l'operazione non è stata inviata. \\ \addlinespace
Estende & \hyperref[uc:registra]{UC20} \{il dipendente rinuncia prima della conferma\}; \hyperref[uc:annulla_inv]{UC29} \{il dipendente rinuncia prima della conferma\} \\ \addlinespace
Requisiti & \hyperlink{req:RF-51}{RF-51} \\
\end{scheda}

\subsubsection{UC27 -- Operazione respinta dal servizio}\label{uc:op_rifiuto}
\begin{figure}[H]
  \centering
  \begin{adjustbox}{max width=\linewidth, max totalheight=0.62\textheight}
  \begin{tikzpicture}
    \node[ucfocus] (b) at (0.00,0.00) {\textbf{UC27}\\Operazione respinta dal servizio};
    \node[uc] (c0) at (7.15,1.14) {\textbf{UC20}\\Registrazione di una richiesta di ferie o permesso};
    \draw[dipendenza] (b) -- node[etich, sloped] {\guillemotleft extend\guillemotright} (c0);
    \node[uc] (c1) at (7.15,-1.14) {\textbf{UC29}\\Cancellazione di una richiesta};
    \draw[dipendenza] (b) -- node[etich, sloped] {\guillemotleft extend\guillemotright} (c1);
    \begin{scope}[on background layer]
      \node[confine, fit=(b) (c0) (c1)] (sistema) {};
    \end{scope}
    \node[nomesistema] at (sistema.north west) {NEXUM Concierge};
    \attore{ap0}{(-4.83,0.00)}{Dipendente}
    \draw[associazione] (ap0) -- (b);
  \end{tikzpicture}
  \end{adjustbox}
  \caption{Diagramma di UC27 -- Operazione respinta dal servizio}\label{fig:uc:op_rifiuto}
\end{figure}

\begin{scheda}
Attore principale & Dipendente \\ \addlinespace
Precondizioni & Il servizio NEXUM ha respinto l'operazione per una validazione negativa (per esempio un periodo sovrapposto a un'altra richiesta). \\ \addlinespace
Postcondizioni & Il dipendente conosce il motivo; nulla è registrato. \\ \addlinespace
Scenario principale & \begin{enumerate}[nosep, leftmargin=*, before=\vspace{-\baselineskip}\vspace{2pt}, after=\vspace{-\baselineskip}]\item Il sistema mostra il motivo restituito dal servizio in modo comprensibile. \item Il sistema propone di correggere i dati; la correzione porta a un nuovo riepilogo e a una nuova conferma. \end{enumerate} \\ \addlinespace
Estende & \hyperref[uc:registra]{UC20} \{il servizio respinge la validazione\}; \hyperref[uc:annulla_inv]{UC29} \{cancellazione non consentita dal servizio\} \\ \addlinespace
Requisiti & \hyperlink{req:RF-52}{RF-52} \\
\end{scheda}

\subsubsection{UC28 -- Esito dell'operazione incerto}\label{uc:op_incerto}
\begin{figure}[H]
  \centering
  \begin{adjustbox}{max width=\linewidth, max totalheight=0.62\textheight}
  \begin{tikzpicture}
    \node[ucfocus] (b) at (0.00,0.00) {\textbf{UC28}\\Esito dell'operazione incerto};
    \node[uc] (c0) at (7.15,1.14) {\textbf{UC20}\\Registrazione di una richiesta di ferie o permesso};
    \draw[dipendenza] (b) -- node[etich, sloped] {\guillemotleft extend\guillemotright} (c0);
    \node[uc] (c1) at (7.15,-1.14) {\textbf{UC29}\\Cancellazione di una richiesta};
    \draw[dipendenza] (b) -- node[etich, sloped] {\guillemotleft extend\guillemotright} (c1);
    \begin{scope}[on background layer]
      \node[confine, fit=(b) (c0) (c1)] (sistema) {};
    \end{scope}
    \node[nomesistema] at (sistema.north west) {NEXUM Concierge};
    \attore{ap0}{(-4.83,0.00)}{Dipendente}
    \draw[associazione] (ap0) -- (b);
  \end{tikzpicture}
  \end{adjustbox}
  \caption{Diagramma di UC28 -- Esito dell'operazione incerto}\label{fig:uc:op_incerto}
\end{figure}

\begin{scheda}
Descrizione & Un invio interrotto non viene ripetuto alla cieca. \\ \addlinespace
Attore principale & Dipendente \\ \addlinespace
Precondizioni & L'invio si è interrotto e non è noto se il servizio abbia registrato l'operazione. \\ \addlinespace
Postcondizioni & L'esito è verificato; l'operazione non è duplicata. \\ \addlinespace
Scenario principale & \begin{enumerate}[nosep, leftmargin=*, before=\vspace{-\baselineskip}\vspace{2pt}, after=\vspace{-\baselineskip}]\item Il sistema verifica presso il servizio lo stato dell'operazione. \item Se l'operazione risulta registrata, il sistema ne mostra l'esito. \item Se non risulta, il sistema chiede al dipendente se ripetere l'invio. \end{enumerate} \\ \addlinespace
Scenari alternativi & 1a. La verifica non è possibile: il sistema dichiara che l'esito è incerto e invita a controllare più tardi, senza ripetere l'invio. \\ \addlinespace
Estende & \hyperref[uc:registra]{UC20} \{invio interrotto con esito sconosciuto\}; \hyperref[uc:annulla_inv]{UC29} \{invio interrotto con esito sconosciuto\} \\ \addlinespace
Requisiti & \hyperlink{req:RF-53}{RF-53} \\
\end{scheda}

\subsubsection{UC29 -- Cancellazione di una richiesta}\label{uc:annulla_inv}
\begin{figure}[H]
  \centering
  \begin{adjustbox}{max width=\linewidth, max totalheight=0.62\textheight}
  \begin{tikzpicture}
    \node[ucfocus] (b) at (0.00,0.00) {\textbf{UC29}\\Cancellazione di una richiesta};
    \node[uc] (c0) at (7.15,3.56) {\textbf{UC30}\\Individuazione della richiesta interessata};
    \draw[dipendenza] (b) -- node[etich, sloped] {\guillemotleft include\guillemotright} (c0);
    \node[uc] (c1) at (7.15,1.02) {\textbf{UC22}\\Visualizzazione del riepilogo dell'operazione};
    \draw[dipendenza] (b) -- node[etich, sloped] {\guillemotleft include\guillemotright} (c1);
    \node[uc] (c2) at (7.15,-1.27) {\textbf{UC23}\\Conferma esplicita dell'operazione};
    \draw[dipendenza] (b) -- node[etich, sloped] {\guillemotleft include\guillemotright} (c2);
    \node[uc] (c3) at (7.15,-3.56) {\textbf{UC24}\\Visualizzazione dell'esito dell'operazione};
    \draw[dipendenza] (b) -- node[etich, sloped] {\guillemotleft include\guillemotright} (c3);
    \node[uc] (g) at (0.00,2.84) {\textbf{UC6}\\Invio di un messaggio al Concierge};
    \draw[generalizzazione] (b) -- (g);
    \begin{scope}[on background layer]
      \node[confine, fit=(b) (c0) (c1) (c2) (c3) (g)] (sistema) {};
    \end{scope}
    \node[nomesistema] at (sistema.north west) {NEXUM Concierge};
    \attore{ap0}{(-4.83,0.00)}{Dipendente}
    \draw[associazione] (ap0) -- (b);
    \attore{as0}{(0.00,-6.65)}{Servizio\\ferie e permessi NEXUM}
    \draw[associazione] (b) -- (as0);
  \end{tikzpicture}
  \end{adjustbox}
  \caption{Diagramma di UC29 -- Cancellazione di una richiesta}\label{fig:uc:annulla_inv}
\end{figure}

\begin{figure}[H]
  \centering
  \begin{adjustbox}{max width=\linewidth, max totalheight=0.62\textheight}
  \begin{tikzpicture}
    \node[ucfocus] (b) at (0.00,0.00) {\textbf{UC29}\\Cancellazione di una richiesta};
    \node[uc] (e0) at (-5.00,4.00) {\textbf{UC26}\\Rinuncia all'operazione prima della conferma};
    \draw[dipendenza] (e0) -- (b) node[etich, pos=0.3, sloped] {\guillemotleft extend\guillemotright} coordinate[pos=0.5] (m0);
    \node[nota] (n0) at (-6.45,1.75) {\textbf{Condizione:}\\\{il dipendente rinuncia prima della conferma\}};
    \draw[ancora] (n0) -- (m0);
    \node[uc] (e1) at (5.00,4.00) {\textbf{UC27}\\Operazione respinta dal servizio};
    \draw[dipendenza] (e1) -- (b) node[etich, pos=0.3, sloped] {\guillemotleft extend\guillemotright} coordinate[pos=0.5] (m1);
    \node[nota] (n1) at (6.45,1.75) {\textbf{Condizione:}\\\{cancellazione non consentita dal servizio\}};
    \draw[ancora] (n1) -- (m1);
    \node[uc] (e2) at (5.00,-4.00) {\textbf{UC28}\\Esito dell'operazione incerto};
    \draw[dipendenza] (e2) -- (b) node[etich, pos=0.3, sloped] {\guillemotleft extend\guillemotright} coordinate[pos=0.5] (m2);
    \node[nota] (n2) at (6.45,-1.75) {\textbf{Condizione:}\\\{invio interrotto con esito sconosciuto\}};
    \draw[ancora] (n2) -- (m2);
    \begin{scope}[on background layer]
      \node[confine, fit=(b) (e0) (n0) (e1) (n1) (e2) (n2)] (sistema) {};
    \end{scope}
    \node[nomesistema] at (sistema.north west) {NEXUM Concierge};
    \attore{ap0}{(-10.25,0.00)}{Dipendente}
    \draw[associazione] (ap0) -- (b);
  \end{tikzpicture}
  \end{adjustbox}
  \caption{Estensioni di UC29 -- Cancellazione di una richiesta}\label{fig:uc:annulla_inv:ext}
\end{figure}

\begin{scheda}
Descrizione & Il dipendente chiede di cancellare una propria richiesta. Non esistono ritiro o modifica: per correggere una richiesta il dipendente la cancella e ne inserisce una nuova. Le condizioni per le richieste passate dipendono dalle API Eggon. \\ \addlinespace
Attore principale & Dipendente \\ \addlinespace
Attori secondari & Servizio ferie e permessi NEXUM \\ \addlinespace
Precondizioni & Le API Eggon espongono la cancellazione delle richieste. \\ \addlinespace
Postcondizioni & La richiesta è cancellata in NEXUM e il dipendente ne conosce l'esito. \\ \addlinespace
Scenario principale & \begin{enumerate}[nosep, leftmargin=*, before=\vspace{-\baselineskip}\vspace{2pt}, after=\vspace{-\baselineskip}]\item Il dipendente chiede di cancellare una richiesta. \item Il sistema individua la richiesta in modo univoco (\hyperref[uc:sel_rich]{UC30}). \item Il sistema mostra il riepilogo della cancellazione (\hyperref[uc:op_riep]{UC22}). \item Il dipendente conferma (\hyperref[uc:op_conf]{UC23}). \item Il sistema invia la cancellazione e mostra l'esito (\hyperref[uc:op_esito]{UC24}). \end{enumerate} \\ \addlinespace
Scenari alternativi & 2a. La richiesta non è individuata in modo univoco: il sistema non procede finché il dipendente non la identifica (\hyperref[uc:chiarimento]{UC7}). \\ \addlinespace
Inclusioni & \hyperref[uc:sel_rich]{UC30}, \hyperref[uc:op_riep]{UC22}, \hyperref[uc:op_conf]{UC23}, \hyperref[uc:op_esito]{UC24} \\ \addlinespace
Estensioni & \hyperref[uc:op_rinuncia]{UC26} \{il dipendente rinuncia prima della conferma\}; \hyperref[uc:op_rifiuto]{UC27} \{cancellazione non consentita dal servizio\}; \hyperref[uc:op_incerto]{UC28} \{invio interrotto con esito sconosciuto\} \\ \addlinespace
Generalizzazione & specializza \hyperref[uc:msg]{UC6} \\ \addlinespace
Requisiti & \hyperlink{req:RF-54}{RF-54}, \hyperlink{req:RF-96}{RF-96}, \hyperlink{req:RF-102}{RF-102} \\
\end{scheda}

\subsubsection{UC30 -- Individuazione della richiesta interessata}\label{uc:sel_rich}
\begin{figure}[H]
  \centering
  \begin{adjustbox}{max width=\linewidth, max totalheight=0.62\textheight}
  \begin{tikzpicture}
    \node[ucfocus] (b) at (0.00,0.00) {\textbf{UC30}\\Individuazione della richiesta interessata};
    \node[uc] (c0) at (7.15,1.14) {\textbf{UC19}\\Consultazione dello stato di una richiesta};
    \draw[dipendenza] (c0) -- node[etich, sloped] {\guillemotleft include\guillemotright} (b);
    \node[uc] (c1) at (7.15,-1.14) {\textbf{UC29}\\Cancellazione di una richiesta};
    \draw[dipendenza] (c1) -- node[etich, sloped] {\guillemotleft include\guillemotright} (b);
    \begin{scope}[on background layer]
      \node[confine, fit=(b) (c0) (c1)] (sistema) {};
    \end{scope}
    \node[nomesistema] at (sistema.north west) {NEXUM Concierge};
    \attore{ap0}{(-4.83,0.00)}{Dipendente}
    \draw[associazione] (ap0) -- (b);
    \attore{as0}{(0.00,-3.99)}{Servizio\\ferie e permessi NEXUM}
    \draw[associazione] (b) -- (as0);
  \end{tikzpicture}
  \end{adjustbox}
  \caption{Diagramma di UC30 -- Individuazione della richiesta interessata}\label{fig:uc:sel_rich}
\end{figure}

\begin{scheda}
Attore principale & Dipendente \\ \addlinespace
Attori secondari & Servizio ferie e permessi NEXUM \\ \addlinespace
Precondizioni & Il dipendente ha fatto riferimento a una propria richiesta. \\ \addlinespace
Postcondizioni & La richiesta interessata è individuata. \\ \addlinespace
Scenario principale & \begin{enumerate}[nosep, leftmargin=*, before=\vspace{-\baselineskip}\vspace{2pt}, after=\vspace{-\baselineskip}]\item Il sistema cerca tra le richieste del dipendente quella indicata dal messaggio o dal contesto. \item Il sistema mostra la richiesta individuata. \end{enumerate} \\ \addlinespace
Scenari alternativi & 1a. Più richieste corrispondono: il sistema le elenca e chiede quale (\hyperref[uc:chiarimento]{UC7}). \\ \addlinespace
Incluso da & \hyperref[uc:stato_rich]{UC19}, \hyperref[uc:annulla_inv]{UC29} \\ \addlinespace
Requisiti & \hyperlink{req:RF-55}{RF-55}, \hyperlink{req:RF-102}{RF-102} \\
\end{scheda}

\subsection{Gestione delle fonti nel backoffice}\label{area:fonti}
Casi d'uso del backoffice con cui gli amministratori organizzano i bucket e caricano, sostituiscono e ritirano i documenti da cui il Concierge ricava le risposte, compresi i manuali NEXUM.

\begin{figure}[H]
  \centering
  \begin{adjustbox}{max width=\linewidth, max totalheight=0.62\textheight}
  \begin{tikzpicture}
    \node[uc] (p0) at (0.00,3.12) {\textbf{UC31}\\Visualizzazione dell'elenco dei bucket};
    \node[uc] (p1) at (0.00,0.88) {\textbf{UC32}\\Creazione di un bucket};
    \node[uc] (p2) at (0.00,-1.12) {\textbf{UC36}\\Modifica di un bucket};
    \node[uc] (p3) at (0.00,-3.12) {\textbf{UC37}\\Eliminazione di un bucket};
    \begin{scope}[on background layer]
      \node[confine, fit=(p0) (p1) (p2) (p3)] (sistema) {};
    \end{scope}
    \node[nomesistema] at (sistema.north west) {NEXUM Concierge};
    \attore{ap0}{(-5.60,0.00)}{Amministratore}
    \draw[associazione] (ap0) -- (p0);
    \draw[associazione] (ap0) -- (p1);
    \draw[associazione] (ap0) -- (p2);
    \draw[associazione] (ap0) -- (p3);
  \end{tikzpicture}
  \end{adjustbox}
  \caption{Panoramica: gestione delle fonti nel backoffice (1)}\label{fig:pano:fonti:0}
\end{figure}

\begin{figure}[H]
  \centering
  \begin{adjustbox}{max width=\linewidth, max totalheight=0.62\textheight}
  \begin{tikzpicture}
    \node[uc] (p0) at (0.00,5.18) {\textbf{UC38}\\Visualizzazione delle fonti di un bucket};
    \node[uc] (p1) at (0.00,2.94) {\textbf{UC39}\\Caricamento di una fonte};
    \node[uc] (p2) at (0.00,0.70) {\textbf{UC43}\\Visualizzazione del dettaglio di una fonte};
    \node[uc] (p3) at (0.00,-1.54) {\textbf{UC45}\\Sostituzione di una fonte};
    \node[uc] (p4) at (0.00,-3.36) {\textbf{UC46}\\Ritiro di una fonte};
    \node[uc] (p5) at (0.00,-5.18) {\textbf{UC47}\\Nuova elaborazione di una fonte};
    \begin{scope}[on background layer]
      \node[confine, fit=(p0) (p1) (p2) (p3) (p4) (p5)] (sistema) {};
    \end{scope}
    \node[nomesistema] at (sistema.north west) {NEXUM Concierge};
    \attore{ap0}{(-5.60,0.00)}{Amministratore}
    \draw[associazione] (ap0) -- (p0);
    \draw[associazione] (ap0) -- (p1);
    \draw[associazione] (ap0) -- (p2);
    \draw[associazione] (ap0) -- (p3);
    \draw[associazione] (ap0) -- (p4);
    \draw[associazione] (ap0) -- (p5);
    \attore{as0}{(5.60,0.00)}{Fornitore\\di modelli AI}
    \draw[associazione] (as0) -- (p1);
    \draw[associazione] (as0) -- (p3);
    \draw[associazione] (as0) -- (p5);
  \end{tikzpicture}
  \end{adjustbox}
  \caption{Panoramica: gestione delle fonti nel backoffice (2)}\label{fig:pano:fonti:1}
\end{figure}

\subsubsection{UC31 -- Visualizzazione dell'elenco dei bucket}\label{uc:bk_elenco}
\begin{figure}[H]
  \centering
  \begin{adjustbox}{max width=\linewidth, max totalheight=0.62\textheight}
  \begin{tikzpicture}
    \node[ucfocus] (b) at (0.00,0.00) {\textbf{UC31}\\Visualizzazione dell'elenco dei bucket};
    \node[uc] (c0) at (7.15,0.00) {\textbf{UC31.1}\\Visualizzazione di un bucket nell'elenco};
    \draw[dipendenza] (b) -- node[etich, sloped] {\guillemotleft include\guillemotright} (c0);
    \begin{scope}[on background layer]
      \node[confine, fit=(b) (c0)] (sistema) {};
    \end{scope}
    \node[nomesistema] at (sistema.north west) {NEXUM Concierge};
    \attore{ap0}{(-4.83,0.00)}{Amministratore}
    \draw[associazione] (ap0) -- (b);
  \end{tikzpicture}
  \end{adjustbox}
  \caption{Diagramma di UC31 -- Visualizzazione dell'elenco dei bucket}\label{fig:uc:bk_elenco}
\end{figure}

\begin{scheda}
Descrizione & L'amministratore vede tutti i bucket dell'azienda, di conoscenza aziendale e di conoscenza del prodotto: gli amministratori sono tutti allo stesso livello. \\ \addlinespace
Attore principale & Amministratore \\ \addlinespace
Precondizioni & L'amministratore è autenticato nel backoffice. \\ \addlinespace
Postcondizioni & L'elenco dei bucket dell'azienda è visualizzato. \\ \addlinespace
Scenario principale & \begin{enumerate}[nosep, leftmargin=*, before=\vspace{-\baselineskip}\vspace{2pt}, after=\vspace{-\baselineskip}]\item L'amministratore accede alla sezione delle fonti. \item Il sistema mostra i bucket dell'azienda (\hyperref[uc:bk_elem]{UC31.1}). \end{enumerate} \\ \addlinespace
Inclusioni & \hyperref[uc:bk_elem]{UC31.1} \\ \addlinespace
Requisiti & \hyperlink{req:RF-57}{RF-57} \\
\end{scheda}

\subsubsubsection{UC31.1 -- Visualizzazione di un bucket nell'elenco}\label{uc:bk_elem}
\begin{figure}[H]
  \centering
  \begin{adjustbox}{max width=\linewidth, max totalheight=0.62\textheight}
  \begin{tikzpicture}
    \node[ucfocus] (b) at (0.00,0.00) {\textbf{UC31.1}\\Visualizzazione di un bucket nell'elenco};
    \node[uc] (c0) at (7.15,0.00) {\textbf{UC31}\\Visualizzazione dell'elenco dei bucket};
    \draw[dipendenza] (c0) -- node[etich, sloped] {\guillemotleft include\guillemotright} (b);
    \begin{scope}[on background layer]
      \node[confine, fit=(b) (c0)] (sistema) {};
    \end{scope}
    \node[nomesistema] at (sistema.north west) {NEXUM Concierge};
    \attore{ap0}{(-4.83,0.00)}{Amministratore}
    \draw[associazione] (ap0) -- (b);
  \end{tikzpicture}
  \end{adjustbox}
  \caption{Diagramma di UC31.1 -- Visualizzazione di un bucket nell'elenco}\label{fig:uc:bk_elem}
\end{figure}

\begin{scheda}
Attore principale & Amministratore \\ \addlinespace
Precondizioni & L'elenco dei bucket è visualizzato. \\ \addlinespace
Postcondizioni & Le informazioni del bucket sono visibili. \\ \addlinespace
Scenario principale & Il sistema mostra per il bucket nome, descrizione, categoria (conoscenza aziendale o del prodotto) e numero di fonti per stato. \\ \addlinespace
Incluso da & \hyperref[uc:bk_elenco]{UC31} \\ \addlinespace
Requisiti & \hyperlink{req:RF-58}{RF-58} \\
\end{scheda}

\subsubsection{UC32 -- Creazione di un bucket}\label{uc:bk_crea}
\begin{figure}[H]
  \centering
  \begin{adjustbox}{max width=\linewidth, max totalheight=0.62\textheight}
  \begin{tikzpicture}
    \node[ucfocus] (b) at (0.00,0.00) {\textbf{UC32}\\Creazione di un bucket};
    \node[uc] (c0) at (7.15,1.14) {\textbf{UC32.1}\\Inserimento del nome del bucket};
    \draw[dipendenza] (b) -- node[etich, sloped] {\guillemotleft include\guillemotright} (c0);
    \node[uc] (c1) at (7.15,-1.14) {\textbf{UC32.2}\\Inserimento della descrizione del bucket};
    \draw[dipendenza] (b) -- node[etich, sloped] {\guillemotleft include\guillemotright} (c1);
    \begin{scope}[on background layer]
      \node[confine, fit=(b) (c0) (c1)] (sistema) {};
    \end{scope}
    \node[nomesistema] at (sistema.north west) {NEXUM Concierge};
    \attore{ap0}{(-4.83,0.00)}{Amministratore}
    \draw[associazione] (ap0) -- (b);
  \end{tikzpicture}
  \end{adjustbox}
  \caption{Diagramma di UC32 -- Creazione di un bucket}\label{fig:uc:bk_crea}
\end{figure}

\begin{figure}[H]
  \centering
  \begin{adjustbox}{max width=\linewidth, max totalheight=0.62\textheight}
  \begin{tikzpicture}
    \node[ucfocus] (b) at (0.00,0.00) {\textbf{UC32}\\Creazione di un bucket};
    \node[uc] (e0) at (-5.00,4.00) {\textbf{UC34}\\Nome del bucket già in uso};
    \draw[dipendenza] (e0) -- (b) node[etich, pos=0.3, sloped] {\guillemotleft extend\guillemotright} coordinate[pos=0.5] (m0);
    \node[nota] (n0) at (-6.45,1.75) {\textbf{Condizione:}\\\{nome già usato nell'azienda\}};
    \draw[ancora] (n0) -- (m0);
    \node[uc] (e1) at (5.00,4.00) {\textbf{UC35}\\Dati del bucket mancanti};
    \draw[dipendenza] (e1) -- (b) node[etich, pos=0.3, sloped] {\guillemotleft extend\guillemotright} coordinate[pos=0.5] (m1);
    \node[nota] (n1) at (6.45,1.75) {\textbf{Condizione:}\\\{nome o descrizione mancanti\}};
    \draw[ancora] (n1) -- (m1);
    \node[uc] (e2) at (5.00,-4.00) {\textbf{UC33}\\Limitazione della consultazione a gruppi};
    \draw[dipendenza] (e2) -- (b) node[etich, pos=0.3, sloped] {\guillemotleft extend\guillemotright} coordinate[pos=0.5] (m2);
    \node[nota] (n2) at (6.45,-1.75) {\textbf{Condizione:}\\\{l'amministratore vuole limitare la consultazione a gruppi di dipendenti\}};
    \draw[ancora] (n2) -- (m2);
    \begin{scope}[on background layer]
      \node[confine, fit=(b) (e0) (n0) (e1) (n1) (e2) (n2)] (sistema) {};
    \end{scope}
    \node[nomesistema] at (sistema.north west) {NEXUM Concierge};
    \attore{ap0}{(-10.25,0.00)}{Amministratore}
    \draw[associazione] (ap0) -- (b);
  \end{tikzpicture}
  \end{adjustbox}
  \caption{Estensioni di UC32 -- Creazione di un bucket}\label{fig:uc:bk_crea:ext}
\end{figure}

\begin{scheda}
Descrizione & L'amministratore crea un contenitore di fonti con una finalità e una categoria. Le conoscenze del bucket sono consultabili da tutti i dipendenti dell'azienda. \\ \addlinespace
Attore principale & Amministratore \\ \addlinespace
Precondizioni & L'amministratore è autenticato nel backoffice. \\ \addlinespace
Postcondizioni & Il bucket esiste, vuoto, con nome, descrizione e categoria. \\ \addlinespace
Scenario principale & \begin{enumerate}[nosep, leftmargin=*, before=\vspace{-\baselineskip}\vspace{2pt}, after=\vspace{-\baselineskip}]\item L'amministratore sceglie di creare un bucket. \item L'amministratore inserisce il nome (\hyperref[uc:bk_nome]{UC32.1}). \item L'amministratore inserisce la descrizione (\hyperref[uc:bk_desc]{UC32.2}). \item L'amministratore sceglie la categoria: conoscenza aziendale o conoscenza del prodotto. \item L'amministratore conferma. \item Il sistema crea il bucket, consultabile da tutti i dipendenti dell'azienda. \end{enumerate} \\ \addlinespace
Inclusioni & \hyperref[uc:bk_nome]{UC32.1}, \hyperref[uc:bk_desc]{UC32.2} \\ \addlinespace
Estensioni & \hyperref[uc:bk_nome_dup]{UC34} \{nome già usato nell'azienda\}; \hyperref[uc:bk_incompleto]{UC35} \{nome o descrizione mancanti\}; \hyperref[uc:bk_dest]{UC33} \{l'amministratore vuole limitare la consultazione a gruppi di dipendenti\} \\ \addlinespace
Requisiti & \hyperlink{req:RF-59}{RF-59}, \hyperlink{req:RF-60}{RF-60}, \hyperlink{req:RF-63}{RF-63} \\
\end{scheda}

\subsubsubsection{UC32.1 -- Inserimento del nome del bucket}\label{uc:bk_nome}
\begin{figure}[H]
  \centering
  \begin{adjustbox}{max width=\linewidth, max totalheight=0.62\textheight}
  \begin{tikzpicture}
    \node[ucfocus] (b) at (0.00,0.00) {\textbf{UC32.1}\\Inserimento del nome del bucket};
    \node[uc] (c0) at (7.15,1.02) {\textbf{UC32}\\Creazione di un bucket};
    \draw[dipendenza] (c0) -- node[etich, sloped] {\guillemotleft include\guillemotright} (b);
    \node[uc] (c1) at (7.15,-1.02) {\textbf{UC36}\\Modifica di un bucket};
    \draw[dipendenza] (c1) -- node[etich, sloped] {\guillemotleft include\guillemotright} (b);
    \begin{scope}[on background layer]
      \node[confine, fit=(b) (c0) (c1)] (sistema) {};
    \end{scope}
    \node[nomesistema] at (sistema.north west) {NEXUM Concierge};
    \attore{ap0}{(-4.83,0.00)}{Amministratore}
    \draw[associazione] (ap0) -- (b);
  \end{tikzpicture}
  \end{adjustbox}
  \caption{Diagramma di UC32.1 -- Inserimento del nome del bucket}\label{fig:uc:bk_nome}
\end{figure}

\begin{scheda}
Attore principale & Amministratore \\ \addlinespace
Precondizioni & È aperto il modulo di creazione o modifica di un bucket. \\ \addlinespace
Postcondizioni & Il nome è inserito. \\ \addlinespace
Scenario principale & L'amministratore inserisce il nome del bucket. \\ \addlinespace
Incluso da & \hyperref[uc:bk_crea]{UC32}, \hyperref[uc:bk_modifica]{UC36} \\ \addlinespace
Requisiti & \hyperlink{req:RF-59}{RF-59} \\
\end{scheda}

\subsubsubsection{UC32.2 -- Inserimento della descrizione del bucket}\label{uc:bk_desc}
\begin{figure}[H]
  \centering
  \begin{adjustbox}{max width=\linewidth, max totalheight=0.62\textheight}
  \begin{tikzpicture}
    \node[ucfocus] (b) at (0.00,0.00) {\textbf{UC32.2}\\Inserimento della descrizione del bucket};
    \node[uc] (c0) at (7.15,1.02) {\textbf{UC32}\\Creazione di un bucket};
    \draw[dipendenza] (c0) -- node[etich, sloped] {\guillemotleft include\guillemotright} (b);
    \node[uc] (c1) at (7.15,-1.02) {\textbf{UC36}\\Modifica di un bucket};
    \draw[dipendenza] (c1) -- node[etich, sloped] {\guillemotleft include\guillemotright} (b);
    \begin{scope}[on background layer]
      \node[confine, fit=(b) (c0) (c1)] (sistema) {};
    \end{scope}
    \node[nomesistema] at (sistema.north west) {NEXUM Concierge};
    \attore{ap0}{(-4.83,0.00)}{Amministratore}
    \draw[associazione] (ap0) -- (b);
  \end{tikzpicture}
  \end{adjustbox}
  \caption{Diagramma di UC32.2 -- Inserimento della descrizione del bucket}\label{fig:uc:bk_desc}
\end{figure}

\begin{scheda}
Attore principale & Amministratore \\ \addlinespace
Precondizioni & È aperto il modulo di creazione o modifica di un bucket. \\ \addlinespace
Postcondizioni & La descrizione è inserita. \\ \addlinespace
Scenario principale & L'amministratore descrive la finalità del bucket in un campo di testo libero. \\ \addlinespace
Incluso da & \hyperref[uc:bk_crea]{UC32}, \hyperref[uc:bk_modifica]{UC36} \\ \addlinespace
Requisiti & \hyperlink{req:RF-59}{RF-59} \\
\end{scheda}

\subsubsection{UC33 -- Limitazione della consultazione a gruppi}\label{uc:bk_dest}
\begin{figure}[H]
  \centering
  \begin{adjustbox}{max width=\linewidth, max totalheight=0.62\textheight}
  \begin{tikzpicture}
    \node[ucfocus] (b) at (0.00,0.00) {\textbf{UC33}\\Limitazione della consultazione a gruppi};
    \node[uc] (c0) at (7.15,1.02) {\textbf{UC32}\\Creazione di un bucket};
    \draw[dipendenza] (b) -- node[etich, sloped] {\guillemotleft extend\guillemotright} (c0);
    \node[uc] (c1) at (7.15,-1.02) {\textbf{UC36}\\Modifica di un bucket};
    \draw[dipendenza] (b) -- node[etich, sloped] {\guillemotleft extend\guillemotright} (c1);
    \begin{scope}[on background layer]
      \node[confine, fit=(b) (c0) (c1)] (sistema) {};
    \end{scope}
    \node[nomesistema] at (sistema.north west) {NEXUM Concierge};
    \attore{ap0}{(-4.83,0.00)}{Amministratore}
    \draw[associazione] (ap0) -- (b);
  \end{tikzpicture}
  \end{adjustbox}
  \caption{Diagramma di UC33 -- Limitazione della consultazione a gruppi}\label{fig:uc:bk_dest}
\end{figure}

\begin{scheda}
Attore principale & Amministratore \\ \addlinespace
Precondizioni & È aperto il modulo di creazione o modifica di un bucket. \\ \addlinespace
Postcondizioni & Il bucket è consultabile solo dai gruppi selezionati; le modifiche valgono dal recupero successivo. \\ \addlinespace
Scenario principale & L'amministratore seleziona i gruppi di dipendenti che possono consultare il bucket. \\ \addlinespace
Estende & \hyperref[uc:bk_crea]{UC32} \{l'amministratore vuole limitare la consultazione a gruppi di dipendenti\}; \hyperref[uc:bk_modifica]{UC36} \{l'amministratore vuole cambiare i gruppi che consultano il bucket\} \\ \addlinespace
Requisiti & \hyperlink{req:RF-101}{RF-101} \\
\end{scheda}

\subsubsection{UC34 -- Nome del bucket già in uso}\label{uc:bk_nome_dup}
\begin{figure}[H]
  \centering
  \begin{adjustbox}{max width=\linewidth, max totalheight=0.62\textheight}
  \begin{tikzpicture}
    \node[ucfocus] (b) at (0.00,0.00) {\textbf{UC34}\\Nome del bucket già in uso};
    \node[uc] (c0) at (7.15,1.02) {\textbf{UC32}\\Creazione di un bucket};
    \draw[dipendenza] (b) -- node[etich, sloped] {\guillemotleft extend\guillemotright} (c0);
    \node[uc] (c1) at (7.15,-1.02) {\textbf{UC36}\\Modifica di un bucket};
    \draw[dipendenza] (b) -- node[etich, sloped] {\guillemotleft extend\guillemotright} (c1);
    \begin{scope}[on background layer]
      \node[confine, fit=(b) (c0) (c1)] (sistema) {};
    \end{scope}
    \node[nomesistema] at (sistema.north west) {NEXUM Concierge};
    \attore{ap0}{(-4.83,0.00)}{Amministratore}
    \draw[associazione] (ap0) -- (b);
  \end{tikzpicture}
  \end{adjustbox}
  \caption{Diagramma di UC34 -- Nome del bucket già in uso}\label{fig:uc:bk_nome_dup}
\end{figure}

\begin{scheda}
Attore principale & Amministratore \\ \addlinespace
Precondizioni & Il nome inserito è già usato da un altro bucket della stessa azienda. \\ \addlinespace
Postcondizioni & Il bucket non è salvato; i dati inseriti restano nel modulo. \\ \addlinespace
Scenario principale & Il sistema segnala che il nome è già in uso e chiede di cambiarlo. \\ \addlinespace
Estende & \hyperref[uc:bk_crea]{UC32} \{nome già usato nell'azienda\}; \hyperref[uc:bk_modifica]{UC36} \{nome già usato nell'azienda\} \\ \addlinespace
Requisiti & \hyperlink{req:RF-64}{RF-64} \\
\end{scheda}

\subsubsection{UC35 -- Dati del bucket mancanti}\label{uc:bk_incompleto}
\begin{figure}[H]
  \centering
  \begin{adjustbox}{max width=\linewidth, max totalheight=0.62\textheight}
  \begin{tikzpicture}
    \node[ucfocus] (b) at (0.00,0.00) {\textbf{UC35}\\Dati del bucket mancanti};
    \node[uc] (c0) at (7.15,1.02) {\textbf{UC32}\\Creazione di un bucket};
    \draw[dipendenza] (b) -- node[etich, sloped] {\guillemotleft extend\guillemotright} (c0);
    \node[uc] (c1) at (7.15,-1.02) {\textbf{UC36}\\Modifica di un bucket};
    \draw[dipendenza] (b) -- node[etich, sloped] {\guillemotleft extend\guillemotright} (c1);
    \begin{scope}[on background layer]
      \node[confine, fit=(b) (c0) (c1)] (sistema) {};
    \end{scope}
    \node[nomesistema] at (sistema.north west) {NEXUM Concierge};
    \attore{ap0}{(-4.83,0.00)}{Amministratore}
    \draw[associazione] (ap0) -- (b);
  \end{tikzpicture}
  \end{adjustbox}
  \caption{Diagramma di UC35 -- Dati del bucket mancanti}\label{fig:uc:bk_incompleto}
\end{figure}

\begin{scheda}
Attore principale & Amministratore \\ \addlinespace
Precondizioni & Il nome o la descrizione non sono stati inseriti. \\ \addlinespace
Postcondizioni & Il bucket non è salvato; i dati inseriti restano nel modulo. \\ \addlinespace
Scenario principale & Il sistema evidenzia i campi obbligatori mancanti. \\ \addlinespace
Estende & \hyperref[uc:bk_crea]{UC32} \{nome o descrizione mancanti\}; \hyperref[uc:bk_modifica]{UC36} \{nome o descrizione mancanti\} \\ \addlinespace
Requisiti & \hyperlink{req:RF-65}{RF-65} \\
\end{scheda}

\subsubsection{UC36 -- Modifica di un bucket}\label{uc:bk_modifica}
\begin{figure}[H]
  \centering
  \begin{adjustbox}{max width=\linewidth, max totalheight=0.62\textheight}
  \begin{tikzpicture}
    \node[ucfocus] (b) at (0.00,0.00) {\textbf{UC36}\\Modifica di un bucket};
    \node[uc] (c0) at (7.15,1.14) {\textbf{UC32.1}\\Inserimento del nome del bucket};
    \draw[dipendenza] (b) -- node[etich, sloped] {\guillemotleft include\guillemotright} (c0);
    \node[uc] (c1) at (7.15,-1.14) {\textbf{UC32.2}\\Inserimento della descrizione del bucket};
    \draw[dipendenza] (b) -- node[etich, sloped] {\guillemotleft include\guillemotright} (c1);
    \begin{scope}[on background layer]
      \node[confine, fit=(b) (c0) (c1)] (sistema) {};
    \end{scope}
    \node[nomesistema] at (sistema.north west) {NEXUM Concierge};
    \attore{ap0}{(-4.83,0.00)}{Amministratore}
    \draw[associazione] (ap0) -- (b);
  \end{tikzpicture}
  \end{adjustbox}
  \caption{Diagramma di UC36 -- Modifica di un bucket}\label{fig:uc:bk_modifica}
\end{figure}

\begin{figure}[H]
  \centering
  \begin{adjustbox}{max width=\linewidth, max totalheight=0.62\textheight}
  \begin{tikzpicture}
    \node[ucfocus] (b) at (0.00,0.00) {\textbf{UC36}\\Modifica di un bucket};
    \node[uc] (e0) at (-5.00,4.00) {\textbf{UC34}\\Nome del bucket già in uso};
    \draw[dipendenza] (e0) -- (b) node[etich, pos=0.3, sloped] {\guillemotleft extend\guillemotright} coordinate[pos=0.5] (m0);
    \node[nota] (n0) at (-6.45,1.75) {\textbf{Condizione:}\\\{nome già usato nell'azienda\}};
    \draw[ancora] (n0) -- (m0);
    \node[uc] (e1) at (5.00,4.00) {\textbf{UC35}\\Dati del bucket mancanti};
    \draw[dipendenza] (e1) -- (b) node[etich, pos=0.3, sloped] {\guillemotleft extend\guillemotright} coordinate[pos=0.5] (m1);
    \node[nota] (n1) at (6.45,1.75) {\textbf{Condizione:}\\\{nome o descrizione mancanti\}};
    \draw[ancora] (n1) -- (m1);
    \node[uc] (e2) at (5.00,-4.00) {\textbf{UC33}\\Limitazione della consultazione a gruppi};
    \draw[dipendenza] (e2) -- (b) node[etich, pos=0.3, sloped] {\guillemotleft extend\guillemotright} coordinate[pos=0.5] (m2);
    \node[nota] (n2) at (6.45,-1.75) {\textbf{Condizione:}\\\{l'amministratore vuole cambiare i gruppi che consultano il bucket\}};
    \draw[ancora] (n2) -- (m2);
    \begin{scope}[on background layer]
      \node[confine, fit=(b) (e0) (n0) (e1) (n1) (e2) (n2)] (sistema) {};
    \end{scope}
    \node[nomesistema] at (sistema.north west) {NEXUM Concierge};
    \attore{ap0}{(-10.25,0.00)}{Amministratore}
    \draw[associazione] (ap0) -- (b);
  \end{tikzpicture}
  \end{adjustbox}
  \caption{Estensioni di UC36 -- Modifica di un bucket}\label{fig:uc:bk_modifica:ext}
\end{figure}

\begin{scheda}
Attore principale & Amministratore \\ \addlinespace
Precondizioni & Il bucket esiste. \\ \addlinespace
Postcondizioni & Il bucket ha il nuovo nome o la nuova descrizione. \\ \addlinespace
Scenario principale & \begin{enumerate}[nosep, leftmargin=*, before=\vspace{-\baselineskip}\vspace{2pt}, after=\vspace{-\baselineskip}]\item L'amministratore sceglie di modificare un bucket. \item L'amministratore cambia nome (\hyperref[uc:bk_nome]{UC32.1}) o descrizione (\hyperref[uc:bk_desc]{UC32.2}). \item L'amministratore conferma. \item Il sistema salva le modifiche. \end{enumerate} \\ \addlinespace
Inclusioni & \hyperref[uc:bk_nome]{UC32.1}, \hyperref[uc:bk_desc]{UC32.2} \\ \addlinespace
Estensioni & \hyperref[uc:bk_nome_dup]{UC34} \{nome già usato nell'azienda\}; \hyperref[uc:bk_incompleto]{UC35} \{nome o descrizione mancanti\}; \hyperref[uc:bk_dest]{UC33} \{l'amministratore vuole cambiare i gruppi che consultano il bucket\} \\ \addlinespace
Requisiti & \hyperlink{req:RF-66}{RF-66}, \hyperlink{req:RF-101}{RF-101} \\
\end{scheda}

\subsubsection{UC37 -- Eliminazione di un bucket}\label{uc:bk_elimina}
\begin{figure}[H]
  \centering
  \begin{adjustbox}{max width=\linewidth, max totalheight=0.62\textheight}
  \begin{tikzpicture}
    \node[ucfocus] (b) at (0.00,0.00) {\textbf{UC37}\\Eliminazione di un bucket};
    \begin{scope}[on background layer]
      \node[confine, fit=(b)] (sistema) {};
    \end{scope}
    \node[nomesistema] at (sistema.north west) {NEXUM Concierge};
    \attore{ap0}{(-4.83,0.00)}{Amministratore}
    \draw[associazione] (ap0) -- (b);
  \end{tikzpicture}
  \end{adjustbox}
  \caption{Diagramma di UC37 -- Eliminazione di un bucket}\label{fig:uc:bk_elimina}
\end{figure}

\begin{scheda}
Attore principale & Amministratore \\ \addlinespace
Precondizioni & Il bucket esiste. \\ \addlinespace
Postcondizioni & Il bucket non è più visibile; le sue fonti sono ritirate ed escluse dal recupero, e restano nella traccia. \\ \addlinespace
Scenario principale & \begin{enumerate}[nosep, leftmargin=*, before=\vspace{-\baselineskip}\vspace{2pt}, after=\vspace{-\baselineskip}]\item L'amministratore sceglie di eliminare un bucket. \item Il sistema avvisa che le fonti contenute saranno ritirate e chiede conferma. \item L'amministratore conferma. \item Il sistema ritira le fonti ed elimina il bucket. \end{enumerate} \\ \addlinespace
Scenari alternativi & 3a. L'amministratore annulla: il bucket resta invariato. \\ \addlinespace
Requisiti & \hyperlink{req:RF-67}{RF-67} \\
\end{scheda}

\subsubsection{UC38 -- Visualizzazione delle fonti di un bucket}\label{uc:fn_elenco}
\begin{figure}[H]
  \centering
  \begin{adjustbox}{max width=\linewidth, max totalheight=0.62\textheight}
  \begin{tikzpicture}
    \node[ucfocus] (b) at (0.00,0.00) {\textbf{UC38}\\Visualizzazione delle fonti di un bucket};
    \node[uc] (c0) at (7.15,0.00) {\textbf{UC38.1}\\Visualizzazione di una fonte nell'elenco};
    \draw[dipendenza] (b) -- node[etich, sloped] {\guillemotleft include\guillemotright} (c0);
    \begin{scope}[on background layer]
      \node[confine, fit=(b) (c0)] (sistema) {};
    \end{scope}
    \node[nomesistema] at (sistema.north west) {NEXUM Concierge};
    \attore{ap0}{(-4.83,0.00)}{Amministratore}
    \draw[associazione] (ap0) -- (b);
  \end{tikzpicture}
  \end{adjustbox}
  \caption{Diagramma di UC38 -- Visualizzazione delle fonti di un bucket}\label{fig:uc:fn_elenco}
\end{figure}

\begin{scheda}
Attore principale & Amministratore \\ \addlinespace
Precondizioni & Il bucket esiste. \\ \addlinespace
Postcondizioni & L'elenco delle fonti del bucket è visualizzato. \\ \addlinespace
Scenario principale & \begin{enumerate}[nosep, leftmargin=*, before=\vspace{-\baselineskip}\vspace{2pt}, after=\vspace{-\baselineskip}]\item L'amministratore apre un bucket. \item Il sistema mostra le fonti contenute (\hyperref[uc:fn_elem]{UC38.1}). \end{enumerate} \\ \addlinespace
Inclusioni & \hyperref[uc:fn_elem]{UC38.1} \\ \addlinespace
Requisiti & \hyperlink{req:RF-68}{RF-68} \\
\end{scheda}

\subsubsubsection{UC38.1 -- Visualizzazione di una fonte nell'elenco}\label{uc:fn_elem}
\begin{figure}[H]
  \centering
  \begin{adjustbox}{max width=\linewidth, max totalheight=0.62\textheight}
  \begin{tikzpicture}
    \node[ucfocus] (b) at (0.00,0.00) {\textbf{UC38.1}\\Visualizzazione di una fonte nell'elenco};
    \node[uc] (c0) at (7.15,0.00) {\textbf{UC38}\\Visualizzazione delle fonti di un bucket};
    \draw[dipendenza] (c0) -- node[etich, sloped] {\guillemotleft include\guillemotright} (b);
    \begin{scope}[on background layer]
      \node[confine, fit=(b) (c0)] (sistema) {};
    \end{scope}
    \node[nomesistema] at (sistema.north west) {NEXUM Concierge};
    \attore{ap0}{(-4.83,0.00)}{Amministratore}
    \draw[associazione] (ap0) -- (b);
  \end{tikzpicture}
  \end{adjustbox}
  \caption{Diagramma di UC38.1 -- Visualizzazione di una fonte nell'elenco}\label{fig:uc:fn_elem}
\end{figure}

\begin{scheda}
Attore principale & Amministratore \\ \addlinespace
Precondizioni & L'elenco delle fonti è visualizzato. \\ \addlinespace
Postcondizioni & Le informazioni principali della fonte sono visibili. \\ \addlinespace
Scenario principale & Il sistema mostra per la fonte titolo, versione, data di caricamento, stato di elaborazione e disponibilità alla consultazione. \\ \addlinespace
Incluso da & \hyperref[uc:fn_elenco]{UC38} \\ \addlinespace
Requisiti & \hyperlink{req:RF-69}{RF-69} \\
\end{scheda}

\subsubsection{UC39 -- Caricamento di una fonte}\label{uc:fn_carica}
\begin{figure}[H]
  \centering
  \begin{adjustbox}{max width=\linewidth, max totalheight=0.62\textheight}
  \begin{tikzpicture}
    \node[ucfocus] (b) at (0.00,0.00) {\textbf{UC39}\\Caricamento di una fonte};
    \node[uc] (c0) at (7.15,1.06) {\textbf{UC39.1}\\Selezione del file};
    \draw[dipendenza] (b) -- node[etich, sloped] {\guillemotleft include\guillemotright} (c0);
    \node[uc] (c1) at (7.15,-1.06) {\textbf{UC39.2}\\Inserimento dei metadati della fonte};
    \draw[dipendenza] (b) -- node[etich, sloped] {\guillemotleft include\guillemotright} (c1);
    \begin{scope}[on background layer]
      \node[confine, fit=(b) (c0) (c1)] (sistema) {};
    \end{scope}
    \node[nomesistema] at (sistema.north west) {NEXUM Concierge};
    \attore{ap0}{(-4.83,0.00)}{Amministratore}
    \draw[associazione] (ap0) -- (b);
    \attore{as0}{(0.00,-4.15)}{Fornitore\\di modelli AI}
    \draw[associazione] (b) -- (as0);
  \end{tikzpicture}
  \end{adjustbox}
  \caption{Diagramma di UC39 -- Caricamento di una fonte}\label{fig:uc:fn_carica}
\end{figure}

\begin{figure}[H]
  \centering
  \begin{adjustbox}{max width=\linewidth, max totalheight=0.62\textheight}
  \begin{tikzpicture}
    \node[ucfocus] (b) at (0.00,0.00) {\textbf{UC39}\\Caricamento di una fonte};
    \node[uc] (e0) at (-5.00,4.00) {\textbf{UC40}\\Formato del file non supportato};
    \draw[dipendenza] (e0) -- (b) node[etich, pos=0.3, sloped] {\guillemotleft extend\guillemotright} coordinate[pos=0.5] (m0);
    \node[nota] (n0) at (-6.45,1.75) {\textbf{Condizione:}\\\{file non in formato PDF\}};
    \draw[ancora] (n0) -- (m0);
    \node[uc] (e1) at (5.00,4.00) {\textbf{UC41}\\Dimensione massima superata};
    \draw[dipendenza] (e1) -- (b) node[etich, pos=0.3, sloped] {\guillemotleft extend\guillemotright} coordinate[pos=0.5] (m1);
    \node[nota] (n1) at (6.45,1.75) {\textbf{Condizione:}\\\{file oltre la dimensione massima\}};
    \draw[ancora] (n1) -- (m1);
    \node[uc] (e2) at (5.00,-4.00) {\textbf{UC42}\\Fonte già presente};
    \draw[dipendenza] (e2) -- (b) node[etich, pos=0.3, sloped] {\guillemotleft extend\guillemotright} coordinate[pos=0.5] (m2);
    \node[nota] (n2) at (6.45,-1.75) {\textbf{Condizione:}\\\{file già presente nel bucket\}};
    \draw[ancora] (n2) -- (m2);
    \begin{scope}[on background layer]
      \node[confine, fit=(b) (e0) (n0) (e1) (n1) (e2) (n2)] (sistema) {};
    \end{scope}
    \node[nomesistema] at (sistema.north west) {NEXUM Concierge};
    \attore{ap0}{(-10.25,0.00)}{Amministratore}
    \draw[associazione] (ap0) -- (b);
  \end{tikzpicture}
  \end{adjustbox}
  \caption{Estensioni di UC39 -- Caricamento di una fonte}\label{fig:uc:fn_carica:ext}
\end{figure}

\begin{scheda}
Descrizione & L'amministratore carica un documento PDF in un bucket. L'elaborazione avviene in modo asincrono: il documento viene trasformato in porzioni ricercabili che conservano provenienza (fonte e bucket) e versione. \\ \addlinespace
Attore principale & Amministratore \\ \addlinespace
Attori secondari & Fornitore di modelli AI \\ \addlinespace
Precondizioni & Il bucket esiste. \\ \addlinespace
Postcondizioni & \begin{itemize}[nosep, leftmargin=*, before=\vspace{-\baselineskip}\vspace{2pt}, after=\vspace{-\baselineskip}]\item La fonte è registrata nel bucket e la sua elaborazione è avviata; lo stato è consultabile. \item Al termine dell'elaborazione la fonte è disponibile alla consultazione. \end{itemize} \\ \addlinespace
Scenario principale & \begin{enumerate}[nosep, leftmargin=*, before=\vspace{-\baselineskip}\vspace{2pt}, after=\vspace{-\baselineskip}]\item L'amministratore sceglie di caricare una fonte nel bucket. \item L'amministratore seleziona il file (\hyperref[uc:fn_file]{UC39.1}). \item L'amministratore inserisce i metadati (\hyperref[uc:fn_meta]{UC39.2}). \item L'amministratore conferma il caricamento. \item Il sistema registra la fonte con stato «in coda» e avvia l'elaborazione asincrona: estrazione del testo, suddivisione in porzioni e indicizzazione con i metadati. \item Al termine il sistema imposta lo stato «disponibile». \end{enumerate} \\ \addlinespace
Scenari alternativi & 6a. L'elaborazione non riesce: il sistema imposta lo stato «errore» e ne registra il motivo (\hyperref[uc:fn_errore]{UC44}). \\ \addlinespace
Inclusioni & \hyperref[uc:fn_file]{UC39.1}, \hyperref[uc:fn_meta]{UC39.2} \\ \addlinespace
Estensioni & \hyperref[uc:fn_formato]{UC40} \{file non in formato PDF\}; \hyperref[uc:fn_dim]{UC41} \{file oltre la dimensione massima\}; \hyperref[uc:fn_dup]{UC42} \{file già presente nel bucket\} \\ \addlinespace
Requisiti & \hyperlink{req:RF-70}{RF-70}, \hyperlink{req:RF-71}{RF-71}, \hyperlink{req:RF-72}{RF-72} \\
\end{scheda}

\subsubsubsection{UC39.1 -- Selezione del file}\label{uc:fn_file}
\begin{figure}[H]
  \centering
  \begin{adjustbox}{max width=\linewidth, max totalheight=0.62\textheight}
  \begin{tikzpicture}
    \node[ucfocus] (b) at (0.00,0.00) {\textbf{UC39.1}\\Selezione del file};
    \node[uc] (c0) at (7.15,1.02) {\textbf{UC39}\\Caricamento di una fonte};
    \draw[dipendenza] (c0) -- node[etich, sloped] {\guillemotleft include\guillemotright} (b);
    \node[uc] (c1) at (7.15,-1.02) {\textbf{UC45}\\Sostituzione di una fonte};
    \draw[dipendenza] (c1) -- node[etich, sloped] {\guillemotleft include\guillemotright} (b);
    \begin{scope}[on background layer]
      \node[confine, fit=(b) (c0) (c1)] (sistema) {};
    \end{scope}
    \node[nomesistema] at (sistema.north west) {NEXUM Concierge};
    \attore{ap0}{(-4.83,0.00)}{Amministratore}
    \draw[associazione] (ap0) -- (b);
  \end{tikzpicture}
  \end{adjustbox}
  \caption{Diagramma di UC39.1 -- Selezione del file}\label{fig:uc:fn_file}
\end{figure}

\begin{scheda}
Attore principale & Amministratore \\ \addlinespace
Precondizioni & È aperto il modulo di caricamento o sostituzione. \\ \addlinespace
Postcondizioni & Il file è selezionato. \\ \addlinespace
Scenario principale & L'amministratore seleziona dal proprio dispositivo il file PDF da caricare. \\ \addlinespace
Incluso da & \hyperref[uc:fn_carica]{UC39}, \hyperref[uc:fn_sostituisci]{UC45} \\ \addlinespace
Requisiti & \hyperlink{req:RF-73}{RF-73} \\
\end{scheda}

\subsubsubsection{UC39.2 -- Inserimento dei metadati della fonte}\label{uc:fn_meta}
\begin{figure}[H]
  \centering
  \begin{adjustbox}{max width=\linewidth, max totalheight=0.62\textheight}
  \begin{tikzpicture}
    \node[ucfocus] (b) at (0.00,0.00) {\textbf{UC39.2}\\Inserimento dei metadati della fonte};
    \node[uc] (c0) at (7.15,1.02) {\textbf{UC39}\\Caricamento di una fonte};
    \draw[dipendenza] (c0) -- node[etich, sloped] {\guillemotleft include\guillemotright} (b);
    \node[uc] (c1) at (7.15,-1.02) {\textbf{UC45}\\Sostituzione di una fonte};
    \draw[dipendenza] (c1) -- node[etich, sloped] {\guillemotleft include\guillemotright} (b);
    \begin{scope}[on background layer]
      \node[confine, fit=(b) (c0) (c1)] (sistema) {};
    \end{scope}
    \node[nomesistema] at (sistema.north west) {NEXUM Concierge};
    \attore{ap0}{(-4.83,0.00)}{Amministratore}
    \draw[associazione] (ap0) -- (b);
  \end{tikzpicture}
  \end{adjustbox}
  \caption{Diagramma di UC39.2 -- Inserimento dei metadati della fonte}\label{fig:uc:fn_meta}
\end{figure}

\begin{scheda}
Attore principale & Amministratore \\ \addlinespace
Precondizioni & È aperto il modulo di caricamento o sostituzione. \\ \addlinespace
Postcondizioni & I metadati sono inseriti. \\ \addlinespace
Scenario principale & \begin{enumerate}[nosep, leftmargin=*, before=\vspace{-\baselineskip}\vspace{2pt}, after=\vspace{-\baselineskip}]\item L'amministratore inserisce titolo e versione della fonte. \item L'amministratore indica, se rilevanti, la data di entrata in vigore e le condizioni di applicabilità (per esempio sede o tipo di contratto). \end{enumerate} \\ \addlinespace
Scenari alternativi & 1a. L'amministratore non inserisce il titolo: il sistema propone il nome del file. \\ \addlinespace
Incluso da & \hyperref[uc:fn_carica]{UC39}, \hyperref[uc:fn_sostituisci]{UC45} \\ \addlinespace
Requisiti & \hyperlink{req:RF-74}{RF-74}, \hyperlink{req:RF-75}{RF-75} \\
\end{scheda}

\subsubsection{UC40 -- Formato del file non supportato}\label{uc:fn_formato}
\begin{figure}[H]
  \centering
  \begin{adjustbox}{max width=\linewidth, max totalheight=0.62\textheight}
  \begin{tikzpicture}
    \node[ucfocus] (b) at (0.00,0.00) {\textbf{UC40}\\Formato del file non supportato};
    \node[uc] (c0) at (7.15,1.02) {\textbf{UC39}\\Caricamento di una fonte};
    \draw[dipendenza] (b) -- node[etich, sloped] {\guillemotleft extend\guillemotright} (c0);
    \node[uc] (c1) at (7.15,-1.02) {\textbf{UC45}\\Sostituzione di una fonte};
    \draw[dipendenza] (b) -- node[etich, sloped] {\guillemotleft extend\guillemotright} (c1);
    \begin{scope}[on background layer]
      \node[confine, fit=(b) (c0) (c1)] (sistema) {};
    \end{scope}
    \node[nomesistema] at (sistema.north west) {NEXUM Concierge};
    \attore{ap0}{(-4.83,0.00)}{Amministratore}
    \draw[associazione] (ap0) -- (b);
  \end{tikzpicture}
  \end{adjustbox}
  \caption{Diagramma di UC40 -- Formato del file non supportato}\label{fig:uc:fn_formato}
\end{figure}

\begin{scheda}
Attore principale & Amministratore \\ \addlinespace
Precondizioni & Il file selezionato non è un PDF. \\ \addlinespace
Postcondizioni & Il file non è caricato. \\ \addlinespace
Scenario principale & Il sistema segnala che è accettato solo il formato PDF. \\ \addlinespace
Estende & \hyperref[uc:fn_carica]{UC39} \{file non in formato PDF\}; \hyperref[uc:fn_sostituisci]{UC45} \{file non in formato PDF\} \\ \addlinespace
Requisiti & \hyperlink{req:RF-76}{RF-76} \\
\end{scheda}

\subsubsection{UC41 -- Dimensione massima superata}\label{uc:fn_dim}
\begin{figure}[H]
  \centering
  \begin{adjustbox}{max width=\linewidth, max totalheight=0.62\textheight}
  \begin{tikzpicture}
    \node[ucfocus] (b) at (0.00,0.00) {\textbf{UC41}\\Dimensione massima superata};
    \node[uc] (c0) at (7.15,1.02) {\textbf{UC39}\\Caricamento di una fonte};
    \draw[dipendenza] (b) -- node[etich, sloped] {\guillemotleft extend\guillemotright} (c0);
    \node[uc] (c1) at (7.15,-1.02) {\textbf{UC45}\\Sostituzione di una fonte};
    \draw[dipendenza] (b) -- node[etich, sloped] {\guillemotleft extend\guillemotright} (c1);
    \begin{scope}[on background layer]
      \node[confine, fit=(b) (c0) (c1)] (sistema) {};
    \end{scope}
    \node[nomesistema] at (sistema.north west) {NEXUM Concierge};
    \attore{ap0}{(-4.83,0.00)}{Amministratore}
    \draw[associazione] (ap0) -- (b);
  \end{tikzpicture}
  \end{adjustbox}
  \caption{Diagramma di UC41 -- Dimensione massima superata}\label{fig:uc:fn_dim}
\end{figure}

\begin{scheda}
Attore principale & Amministratore \\ \addlinespace
Precondizioni & Il file supera la dimensione massima consentita. \\ \addlinespace
Postcondizioni & Il file non è caricato. \\ \addlinespace
Scenario principale & Il sistema segnala la dimensione massima consentita. \\ \addlinespace
Estende & \hyperref[uc:fn_carica]{UC39} \{file oltre la dimensione massima\}; \hyperref[uc:fn_sostituisci]{UC45} \{file oltre la dimensione massima\} \\ \addlinespace
Requisiti & \hyperlink{req:RF-77}{RF-77} \\
\end{scheda}

\subsubsection{UC42 -- Fonte già presente}\label{uc:fn_dup}
\begin{figure}[H]
  \centering
  \begin{adjustbox}{max width=\linewidth, max totalheight=0.62\textheight}
  \begin{tikzpicture}
    \node[ucfocus] (b) at (0.00,0.00) {\textbf{UC42}\\Fonte già presente};
    \node[uc] (c0) at (7.15,0.00) {\textbf{UC39}\\Caricamento di una fonte};
    \draw[dipendenza] (b) -- node[etich, sloped] {\guillemotleft extend\guillemotright} (c0);
    \begin{scope}[on background layer]
      \node[confine, fit=(b) (c0)] (sistema) {};
    \end{scope}
    \node[nomesistema] at (sistema.north west) {NEXUM Concierge};
    \attore{ap0}{(-4.83,0.00)}{Amministratore}
    \draw[associazione] (ap0) -- (b);
  \end{tikzpicture}
  \end{adjustbox}
  \caption{Diagramma di UC42 -- Fonte già presente}\label{fig:uc:fn_dup}
\end{figure}

\begin{scheda}
Attore principale & Amministratore \\ \addlinespace
Precondizioni & Il file ha lo stesso contenuto di una fonte attiva del bucket. \\ \addlinespace
Postcondizioni & L'amministratore sceglie se annullare o sostituire la fonte esistente. \\ \addlinespace
Scenario principale & \begin{enumerate}[nosep, leftmargin=*, before=\vspace{-\baselineskip}\vspace{2pt}, after=\vspace{-\baselineskip}]\item Il sistema segnala la fonte già presente. \item Il sistema propone di annullare il caricamento o di sostituire la fonte esistente (\hyperref[uc:fn_sostituisci]{UC45}). \end{enumerate} \\ \addlinespace
Estende & \hyperref[uc:fn_carica]{UC39} \{file già presente nel bucket\} \\ \addlinespace
Requisiti & \hyperlink{req:RF-78}{RF-78} \\
\end{scheda}

\subsubsection{UC43 -- Visualizzazione del dettaglio di una fonte}\label{uc:fn_dettaglio}
\begin{figure}[H]
  \centering
  \begin{adjustbox}{max width=\linewidth, max totalheight=0.62\textheight}
  \begin{tikzpicture}
    \node[ucfocus] (b) at (0.00,0.00) {\textbf{UC43}\\Visualizzazione del dettaglio di una fonte};
    \node[uc] (c0) at (7.15,1.14) {\textbf{UC43.1}\\Visualizzazione dei metadati};
    \draw[dipendenza] (b) -- node[etich, sloped] {\guillemotleft include\guillemotright} (c0);
    \node[uc] (c1) at (7.15,-1.14) {\textbf{UC43.2}\\Visualizzazione dello stato di elaborazione};
    \draw[dipendenza] (b) -- node[etich, sloped] {\guillemotleft include\guillemotright} (c1);
    \begin{scope}[on background layer]
      \node[confine, fit=(b) (c0) (c1)] (sistema) {};
    \end{scope}
    \node[nomesistema] at (sistema.north west) {NEXUM Concierge};
    \attore{ap0}{(-4.83,0.00)}{Amministratore}
    \draw[associazione] (ap0) -- (b);
  \end{tikzpicture}
  \end{adjustbox}
  \caption{Diagramma di UC43 -- Visualizzazione del dettaglio di una fonte}\label{fig:uc:fn_dettaglio}
\end{figure}

\begin{figure}[H]
  \centering
  \begin{adjustbox}{max width=\linewidth, max totalheight=0.62\textheight}
  \begin{tikzpicture}
    \node[ucfocus] (b) at (0.00,0.00) {\textbf{UC43}\\Visualizzazione del dettaglio di una fonte};
    \node[uc] (e0) at (6.60,0.00) {\textbf{UC44}\\Visualizzazione dell'errore di elaborazione};
    \draw[dipendenza] (e0) -- (b) node[etich, pos=0.3, sloped] {\guillemotleft extend\guillemotright} coordinate[pos=0.5] (m0);
    \node[nota] (n0) at (3.3,-1.9) {\textbf{Condizione:}\\\{elaborazione non riuscita\}};
    \draw[ancora] (n0) -- (m0);
    \begin{scope}[on background layer]
      \node[confine, fit=(b) (e0) (n0)] (sistema) {};
    \end{scope}
    \node[nomesistema] at (sistema.north west) {NEXUM Concierge};
    \attore{ap0}{(-4.23,0.00)}{Amministratore}
    \draw[associazione] (ap0) -- (b);
  \end{tikzpicture}
  \end{adjustbox}
  \caption{Estensioni di UC43 -- Visualizzazione del dettaglio di una fonte}\label{fig:uc:fn_dettaglio:ext}
\end{figure}

\begin{scheda}
Attore principale & Amministratore \\ \addlinespace
Precondizioni & La fonte esiste in un bucket dell'azienda. \\ \addlinespace
Postcondizioni & Il dettaglio della fonte è visualizzato. \\ \addlinespace
Scenario principale & \begin{enumerate}[nosep, leftmargin=*, before=\vspace{-\baselineskip}\vspace{2pt}, after=\vspace{-\baselineskip}]\item L'amministratore seleziona una fonte. \item Il sistema mostra i metadati (\hyperref[uc:fn_meta_v]{UC43.1}) e lo stato di elaborazione (\hyperref[uc:fn_stato]{UC43.2}). \end{enumerate} \\ \addlinespace
Inclusioni & \hyperref[uc:fn_meta_v]{UC43.1}, \hyperref[uc:fn_stato]{UC43.2} \\ \addlinespace
Estensioni & \hyperref[uc:fn_errore]{UC44} \{elaborazione non riuscita\} \\ \addlinespace
Requisiti & nei sottocasi \hyperref[uc:fn_meta_v]{UC43.1}, \hyperref[uc:fn_stato]{UC43.2} \\
\end{scheda}

\subsubsubsection{UC43.1 -- Visualizzazione dei metadati}\label{uc:fn_meta_v}
\begin{figure}[H]
  \centering
  \begin{adjustbox}{max width=\linewidth, max totalheight=0.62\textheight}
  \begin{tikzpicture}
    \node[ucfocus] (b) at (0.00,0.00) {\textbf{UC43.1}\\Visualizzazione dei metadati};
    \node[uc] (c0) at (7.15,0.00) {\textbf{UC43}\\Visualizzazione del dettaglio di una fonte};
    \draw[dipendenza] (c0) -- node[etich, sloped] {\guillemotleft include\guillemotright} (b);
    \begin{scope}[on background layer]
      \node[confine, fit=(b) (c0)] (sistema) {};
    \end{scope}
    \node[nomesistema] at (sistema.north west) {NEXUM Concierge};
    \attore{ap0}{(-4.83,0.00)}{Amministratore}
    \draw[associazione] (ap0) -- (b);
  \end{tikzpicture}
  \end{adjustbox}
  \caption{Diagramma di UC43.1 -- Visualizzazione dei metadati}\label{fig:uc:fn_meta_v}
\end{figure}

\begin{scheda}
Attore principale & Amministratore \\ \addlinespace
Precondizioni & Il dettaglio della fonte è aperto. \\ \addlinespace
Postcondizioni & I metadati sono visibili. \\ \addlinespace
Scenario principale & Il sistema mostra titolo, versione, bucket, categoria, data di caricamento, autore del caricamento, data di entrata in vigore e condizioni di applicabilità. \\ \addlinespace
Incluso da & \hyperref[uc:fn_dettaglio]{UC43} \\ \addlinespace
Requisiti & \hyperlink{req:RF-80}{RF-80} \\
\end{scheda}

\subsubsubsection{UC43.2 -- Visualizzazione dello stato di elaborazione}\label{uc:fn_stato}
\begin{figure}[H]
  \centering
  \begin{adjustbox}{max width=\linewidth, max totalheight=0.62\textheight}
  \begin{tikzpicture}
    \node[ucfocus] (b) at (0.00,0.00) {\textbf{UC43.2}\\Visualizzazione dello stato di elaborazione};
    \node[uc] (c0) at (7.15,0.00) {\textbf{UC43}\\Visualizzazione del dettaglio di una fonte};
    \draw[dipendenza] (c0) -- node[etich, sloped] {\guillemotleft include\guillemotright} (b);
    \begin{scope}[on background layer]
      \node[confine, fit=(b) (c0)] (sistema) {};
    \end{scope}
    \node[nomesistema] at (sistema.north west) {NEXUM Concierge};
    \attore{ap0}{(-4.83,0.00)}{Amministratore}
    \draw[associazione] (ap0) -- (b);
  \end{tikzpicture}
  \end{adjustbox}
  \caption{Diagramma di UC43.2 -- Visualizzazione dello stato di elaborazione}\label{fig:uc:fn_stato}
\end{figure}

\begin{scheda}
Attore principale & Amministratore \\ \addlinespace
Precondizioni & Il dettaglio della fonte è aperto. \\ \addlinespace
Postcondizioni & Lo stato è visibile. \\ \addlinespace
Scenario principale & Il sistema mostra lo stato (in coda, in elaborazione, disponibile, errore, ritirata) e la data dell'ultimo cambiamento. \\ \addlinespace
Incluso da & \hyperref[uc:fn_dettaglio]{UC43} \\ \addlinespace
Requisiti & \hyperlink{req:RF-81}{RF-81} \\
\end{scheda}

\subsubsection{UC44 -- Visualizzazione dell'errore di elaborazione}\label{uc:fn_errore}
\begin{figure}[H]
  \centering
  \begin{adjustbox}{max width=\linewidth, max totalheight=0.62\textheight}
  \begin{tikzpicture}
    \node[ucfocus] (b) at (0.00,0.00) {\textbf{UC44}\\Visualizzazione dell'errore di elaborazione};
    \node[uc] (c0) at (7.15,0.00) {\textbf{UC43}\\Visualizzazione del dettaglio di una fonte};
    \draw[dipendenza] (b) -- node[etich, sloped] {\guillemotleft extend\guillemotright} (c0);
    \begin{scope}[on background layer]
      \node[confine, fit=(b) (c0)] (sistema) {};
    \end{scope}
    \node[nomesistema] at (sistema.north west) {NEXUM Concierge};
    \attore{ap0}{(-4.83,0.00)}{Amministratore}
    \draw[associazione] (ap0) -- (b);
  \end{tikzpicture}
  \end{adjustbox}
  \caption{Diagramma di UC44 -- Visualizzazione dell'errore di elaborazione}\label{fig:uc:fn_errore}
\end{figure}

\begin{scheda}
Descrizione & L'errore di elaborazione si verifica al caricamento, quando la fonte viene trasformata in porzioni ricercabili; gli errori durante una risposta sono trattati in \hyperref[uc:errore_tecnico]{UC10}. \\ \addlinespace
Attore principale & Amministratore \\ \addlinespace
Precondizioni & L'elaborazione della fonte non è riuscita (per esempio un PDF senza testo estraibile). \\ \addlinespace
Postcondizioni & L'amministratore conosce il motivo e le azioni possibili. \\ \addlinespace
Scenario principale & \begin{enumerate}[nosep, leftmargin=*, before=\vspace{-\baselineskip}\vspace{2pt}, after=\vspace{-\baselineskip}]\item Il sistema mostra il motivo dell'errore in modo comprensibile. \item Il sistema propone di sostituire il file o di ripetere l'elaborazione (\hyperref[uc:fn_riprova]{UC47}). \end{enumerate} \\ \addlinespace
Estende & \hyperref[uc:fn_dettaglio]{UC43} \{elaborazione non riuscita\} \\ \addlinespace
Requisiti & \hyperlink{req:RF-83}{RF-83} \\
\end{scheda}

\subsubsection{UC45 -- Sostituzione di una fonte}\label{uc:fn_sostituisci}
\begin{figure}[H]
  \centering
  \begin{adjustbox}{max width=\linewidth, max totalheight=0.62\textheight}
  \begin{tikzpicture}
    \node[ucfocus] (b) at (0.00,0.00) {\textbf{UC45}\\Sostituzione di una fonte};
    \node[uc] (c0) at (7.15,1.06) {\textbf{UC39.1}\\Selezione del file};
    \draw[dipendenza] (b) -- node[etich, sloped] {\guillemotleft include\guillemotright} (c0);
    \node[uc] (c1) at (7.15,-1.06) {\textbf{UC39.2}\\Inserimento dei metadati della fonte};
    \draw[dipendenza] (b) -- node[etich, sloped] {\guillemotleft include\guillemotright} (c1);
    \begin{scope}[on background layer]
      \node[confine, fit=(b) (c0) (c1)] (sistema) {};
    \end{scope}
    \node[nomesistema] at (sistema.north west) {NEXUM Concierge};
    \attore{ap0}{(-4.83,0.00)}{Amministratore}
    \draw[associazione] (ap0) -- (b);
    \attore{as0}{(0.00,-4.15)}{Fornitore\\di modelli AI}
    \draw[associazione] (b) -- (as0);
  \end{tikzpicture}
  \end{adjustbox}
  \caption{Diagramma di UC45 -- Sostituzione di una fonte}\label{fig:uc:fn_sostituisci}
\end{figure}

\begin{figure}[H]
  \centering
  \begin{adjustbox}{max width=\linewidth, max totalheight=0.62\textheight}
  \begin{tikzpicture}
    \node[ucfocus] (b) at (0.00,0.00) {\textbf{UC45}\\Sostituzione di una fonte};
    \node[uc] (e0) at (6.60,2.30) {\textbf{UC40}\\Formato del file non supportato};
    \draw[dipendenza] (e0) -- (b) node[etich, pos=0.3, sloped] {\guillemotleft extend\guillemotright} coordinate[pos=0.5] (m0);
    \node[nota] (n0) at (2.3,2.75) {\textbf{Condizione:}\\\{file non in formato PDF\}};
    \draw[ancora] (n0) -- (m0);
    \node[uc] (e1) at (6.60,-2.30) {\textbf{UC41}\\Dimensione massima superata};
    \draw[dipendenza] (e1) -- (b) node[etich, pos=0.3, sloped] {\guillemotleft extend\guillemotright} coordinate[pos=0.5] (m1);
    \node[nota] (n1) at (2.3,-2.75) {\textbf{Condizione:}\\\{file oltre la dimensione massima\}};
    \draw[ancora] (n1) -- (m1);
    \begin{scope}[on background layer]
      \node[confine, fit=(b) (e0) (n0) (e1) (n1)] (sistema) {};
    \end{scope}
    \node[nomesistema] at (sistema.north west) {NEXUM Concierge};
    \attore{ap0}{(-4.23,0.00)}{Amministratore}
    \draw[associazione] (ap0) -- (b);
  \end{tikzpicture}
  \end{adjustbox}
  \caption{Estensioni di UC45 -- Sostituzione di una fonte}\label{fig:uc:fn_sostituisci:ext}
\end{figure}

\begin{scheda}
Descrizione & L'amministratore carica una nuova versione di una fonte. La versione precedente resta in uso finché la nuova non è disponibile; poi è esclusa dal recupero ma resta nello storico, così la traccia permette di ricostruire quale versione ha supportato ogni risposta passata. \\ \addlinespace
Attore principale & Amministratore \\ \addlinespace
Attori secondari & Fornitore di modelli AI \\ \addlinespace
Precondizioni & La fonte esiste ed è attiva. \\ \addlinespace
Postcondizioni & La nuova versione è elaborata e, una volta disponibile, è l'unica usata nelle risposte successive; la precedente resta nello storico, collegata alle interazioni in cui è stata usata. \\ \addlinespace
Scenario principale & \begin{enumerate}[nosep, leftmargin=*, before=\vspace{-\baselineskip}\vspace{2pt}, after=\vspace{-\baselineskip}]\item L'amministratore sceglie di sostituire una fonte. \item L'amministratore seleziona il nuovo file (\hyperref[uc:fn_file]{UC39.1}). \item L'amministratore aggiorna i metadati e la versione (\hyperref[uc:fn_meta]{UC39.2}). \item L'amministratore conferma. \item Il sistema elabora la nuova versione in modo asincrono. \item Quando la nuova versione è disponibile, il sistema esclude la precedente dal recupero. \end{enumerate} \\ \addlinespace
Scenari alternativi & 5a. L'elaborazione non riesce: la versione precedente resta in uso e il sistema segnala l'errore (\hyperref[uc:fn_errore]{UC44}). \\ \addlinespace
Inclusioni & \hyperref[uc:fn_file]{UC39.1}, \hyperref[uc:fn_meta]{UC39.2} \\ \addlinespace
Estensioni & \hyperref[uc:fn_formato]{UC40} \{file non in formato PDF\}; \hyperref[uc:fn_dim]{UC41} \{file oltre la dimensione massima\} \\ \addlinespace
Requisiti & \hyperlink{req:RF-84}{RF-84}, \hyperlink{req:RF-85}{RF-85} \\
\end{scheda}

\subsubsection{UC46 -- Ritiro di una fonte}\label{uc:fn_ritira}
\begin{figure}[H]
  \centering
  \begin{adjustbox}{max width=\linewidth, max totalheight=0.62\textheight}
  \begin{tikzpicture}
    \node[ucfocus] (b) at (0.00,0.00) {\textbf{UC46}\\Ritiro di una fonte};
    \begin{scope}[on background layer]
      \node[confine, fit=(b)] (sistema) {};
    \end{scope}
    \node[nomesistema] at (sistema.north west) {NEXUM Concierge};
    \attore{ap0}{(-4.83,0.00)}{Amministratore}
    \draw[associazione] (ap0) -- (b);
  \end{tikzpicture}
  \end{adjustbox}
  \caption{Diagramma di UC46 -- Ritiro di una fonte}\label{fig:uc:fn_ritira}
\end{figure}

\begin{scheda}
Attore principale & Amministratore \\ \addlinespace
Precondizioni & La fonte esiste ed è attiva. \\ \addlinespace
Postcondizioni & La fonte è esclusa dal recupero entro il tempo concordato; resta nello storico e nella traccia. \\ \addlinespace
Scenario principale & \begin{enumerate}[nosep, leftmargin=*, before=\vspace{-\baselineskip}\vspace{2pt}, after=\vspace{-\baselineskip}]\item L'amministratore sceglie di ritirare una fonte. \item Il sistema chiede conferma. \item L'amministratore conferma. \item Il sistema imposta lo stato «ritirata» ed esclude la fonte dal recupero. \end{enumerate} \\ \addlinespace
Scenari alternativi & 3a. L'amministratore annulla: la fonte resta attiva. \\ \addlinespace
Requisiti & \hyperlink{req:RF-86}{RF-86} \\
\end{scheda}

\subsubsection{UC47 -- Nuova elaborazione di una fonte}\label{uc:fn_riprova}
\begin{figure}[H]
  \centering
  \begin{adjustbox}{max width=\linewidth, max totalheight=0.62\textheight}
  \begin{tikzpicture}
    \node[ucfocus] (b) at (0.00,0.00) {\textbf{UC47}\\Nuova elaborazione di una fonte};
    \begin{scope}[on background layer]
      \node[confine, fit=(b)] (sistema) {};
    \end{scope}
    \node[nomesistema] at (sistema.north west) {NEXUM Concierge};
    \attore{ap0}{(-4.83,0.00)}{Amministratore}
    \draw[associazione] (ap0) -- (b);
    \attore{as0}{(0.00,-2.85)}{Fornitore\\di modelli AI}
    \draw[associazione] (b) -- (as0);
  \end{tikzpicture}
  \end{adjustbox}
  \caption{Diagramma di UC47 -- Nuova elaborazione di una fonte}\label{fig:uc:fn_riprova}
\end{figure}

\begin{scheda}
Attore principale & Amministratore \\ \addlinespace
Attori secondari & Fornitore di modelli AI \\ \addlinespace
Precondizioni & La fonte è in stato di errore. \\ \addlinespace
Postcondizioni & L'elaborazione è riavviata. \\ \addlinespace
Scenario principale & \begin{enumerate}[nosep, leftmargin=*, before=\vspace{-\baselineskip}\vspace{2pt}, after=\vspace{-\baselineskip}]\item L'amministratore chiede di ripetere l'elaborazione. \item Il sistema riporta la fonte nello stato «in coda» e la rielabora. \end{enumerate} \\ \addlinespace
Requisiti & \hyperlink{req:RF-87}{RF-87} \\
\end{scheda}

\subsection{Tracciabilità}\label{area:traccia}
Casi d'uso con cui gli amministratori ricostruiscono le interazioni avvenute nella propria azienda: fonti consultate con la versione, capacità o abilità invocate, conferme ed esiti.

\begin{figure}[H]
  \centering
  \begin{adjustbox}{max width=\linewidth, max totalheight=0.62\textheight}
  \begin{tikzpicture}
    \node[uc] (p0) at (0.00,2.24) {\textbf{UC48}\\Consultazione del registro delle interazioni};
    \node[uc] (p1) at (0.00,0.00) {\textbf{UC49}\\Filtraggio del registro};
    \node[uc] (p2) at (0.00,-2.24) {\textbf{UC50}\\Visualizzazione del dettaglio di un'interazione};
    \begin{scope}[on background layer]
      \node[confine, fit=(p0) (p1) (p2)] (sistema) {};
    \end{scope}
    \node[nomesistema] at (sistema.north west) {NEXUM Concierge};
    \attore{ap0}{(-5.60,0.00)}{Amministratore}
    \draw[associazione] (ap0) -- (p0);
    \draw[associazione] (ap0) -- (p1);
    \draw[associazione] (ap0) -- (p2);
  \end{tikzpicture}
  \end{adjustbox}
  \caption{Panoramica: tracciabilità}\label{fig:pano:traccia:0}
\end{figure}

\subsubsection{UC48 -- Consultazione del registro delle interazioni}\label{uc:tr_elenco}
\begin{figure}[H]
  \centering
  \begin{adjustbox}{max width=\linewidth, max totalheight=0.62\textheight}
  \begin{tikzpicture}
    \node[ucfocus] (b) at (0.00,0.00) {\textbf{UC48}\\Consultazione del registro delle interazioni};
    \node[uc] (c0) at (7.15,0.00) {\textbf{UC48.1}\\Visualizzazione di una voce del registro};
    \draw[dipendenza] (b) -- node[etich, sloped] {\guillemotleft include\guillemotright} (c0);
    \begin{scope}[on background layer]
      \node[confine, fit=(b) (c0)] (sistema) {};
    \end{scope}
    \node[nomesistema] at (sistema.north west) {NEXUM Concierge};
    \attore{ap0}{(-4.83,0.00)}{Amministratore}
    \draw[associazione] (ap0) -- (b);
  \end{tikzpicture}
  \end{adjustbox}
  \caption{Diagramma di UC48 -- Consultazione del registro delle interazioni}\label{fig:uc:tr_elenco}
\end{figure}

\begin{scheda}
Descrizione & Gli amministratori consultano la traccia delle interazioni e delle operazioni avvenute nella propria azienda. \\ \addlinespace
Attore principale & Amministratore \\ \addlinespace
Precondizioni & L'amministratore è autenticato nel backoffice. \\ \addlinespace
Postcondizioni & Il registro delle interazioni dell'azienda è visualizzato. \\ \addlinespace
Scenario principale & \begin{enumerate}[nosep, leftmargin=*, before=\vspace{-\baselineskip}\vspace{2pt}, after=\vspace{-\baselineskip}]\item L'amministratore apre il registro. \item Il sistema mostra le interazioni dell'azienda dalla più recente (\hyperref[uc:tr_voce]{UC48.1}). \end{enumerate} \\ \addlinespace
Inclusioni & \hyperref[uc:tr_voce]{UC48.1} \\ \addlinespace
Requisiti & \hyperlink{req:RF-89}{RF-89} \\
\end{scheda}

\subsubsubsection{UC48.1 -- Visualizzazione di una voce del registro}\label{uc:tr_voce}
\begin{figure}[H]
  \centering
  \begin{adjustbox}{max width=\linewidth, max totalheight=0.62\textheight}
  \begin{tikzpicture}
    \node[ucfocus] (b) at (0.00,0.00) {\textbf{UC48.1}\\Visualizzazione di una voce del registro};
    \node[uc] (c0) at (7.15,0.00) {\textbf{UC48}\\Consultazione del registro delle interazioni};
    \draw[dipendenza] (c0) -- node[etich, sloped] {\guillemotleft include\guillemotright} (b);
    \begin{scope}[on background layer]
      \node[confine, fit=(b) (c0)] (sistema) {};
    \end{scope}
    \node[nomesistema] at (sistema.north west) {NEXUM Concierge};
    \attore{ap0}{(-4.83,0.00)}{Amministratore}
    \draw[associazione] (ap0) -- (b);
  \end{tikzpicture}
  \end{adjustbox}
  \caption{Diagramma di UC48.1 -- Visualizzazione di una voce del registro}\label{fig:uc:tr_voce}
\end{figure}

\begin{scheda}
Attore principale & Amministratore \\ \addlinespace
Precondizioni & Il registro è visualizzato. \\ \addlinespace
Postcondizioni & Le informazioni principali della voce sono visibili. \\ \addlinespace
Scenario principale & Il sistema mostra data e ora, persona, azienda, capacità o abilità invocata e, quando applicabili, conferma ed esito. \\ \addlinespace
Incluso da & \hyperref[uc:tr_elenco]{UC48} \\ \addlinespace
Requisiti & \hyperlink{req:RF-90}{RF-90} \\
\end{scheda}

\subsubsection{UC49 -- Filtraggio del registro}\label{uc:tr_filtro}
\begin{figure}[H]
  \centering
  \begin{adjustbox}{max width=\linewidth, max totalheight=0.62\textheight}
  \begin{tikzpicture}
    \node[ucfocus] (b) at (0.00,0.00) {\textbf{UC49}\\Filtraggio del registro};
    \node[uc] (c0) at (7.15,2.05) {\textbf{UC49.1}\\Filtraggio per persona};
    \draw[dipendenza] (b) -- node[etich, sloped] {\guillemotleft include\guillemotright} (c0);
    \node[uc] (c1) at (7.15,0.00) {\textbf{UC49.2}\\Filtraggio per periodo};
    \draw[dipendenza] (b) -- node[etich, sloped] {\guillemotleft include\guillemotright} (c1);
    \node[uc] (c2) at (7.15,-2.05) {\textbf{UC49.3}\\Filtraggio per capacità ed esito};
    \draw[dipendenza] (b) -- node[etich, sloped] {\guillemotleft include\guillemotright} (c2);
    \begin{scope}[on background layer]
      \node[confine, fit=(b) (c0) (c1) (c2)] (sistema) {};
    \end{scope}
    \node[nomesistema] at (sistema.north west) {NEXUM Concierge};
    \attore{ap0}{(-4.83,0.00)}{Amministratore}
    \draw[associazione] (ap0) -- (b);
  \end{tikzpicture}
  \end{adjustbox}
  \caption{Diagramma di UC49 -- Filtraggio del registro}\label{fig:uc:tr_filtro}
\end{figure}

\begin{scheda}
Attore principale & Amministratore \\ \addlinespace
Precondizioni & Il registro è visualizzato. \\ \addlinespace
Postcondizioni & Il registro mostra solo le voci che rispettano i filtri. \\ \addlinespace
Scenario principale & \begin{enumerate}[nosep, leftmargin=*, before=\vspace{-\baselineskip}\vspace{2pt}, after=\vspace{-\baselineskip}]\item L'amministratore imposta uno o più filtri: persona (\hyperref[uc:tr_f_persona]{UC49.1}), periodo (\hyperref[uc:tr_f_periodo]{UC49.2}), capacità o esito (\hyperref[uc:tr_f_capacita]{UC49.3}). \item Il sistema mostra le voci corrispondenti. \end{enumerate} \\ \addlinespace
Scenari alternativi & 2a. Nessuna voce corrisponde: il sistema lo dichiara. \\ \addlinespace
Inclusioni & \hyperref[uc:tr_f_persona]{UC49.1}, \hyperref[uc:tr_f_periodo]{UC49.2}, \hyperref[uc:tr_f_capacita]{UC49.3} \\ \addlinespace
Requisiti & \hyperlink{req:RF-91}{RF-91} \\
\end{scheda}

\subsubsubsection{UC49.1 -- Filtraggio per persona}\label{uc:tr_f_persona}
\begin{figure}[H]
  \centering
  \begin{adjustbox}{max width=\linewidth, max totalheight=0.62\textheight}
  \begin{tikzpicture}
    \node[ucfocus] (b) at (0.00,0.00) {\textbf{UC49.1}\\Filtraggio per persona};
    \node[uc] (c0) at (7.15,0.00) {\textbf{UC49}\\Filtraggio del registro};
    \draw[dipendenza] (c0) -- node[etich, sloped] {\guillemotleft include\guillemotright} (b);
    \begin{scope}[on background layer]
      \node[confine, fit=(b) (c0)] (sistema) {};
    \end{scope}
    \node[nomesistema] at (sistema.north west) {NEXUM Concierge};
    \attore{ap0}{(-4.83,0.00)}{Amministratore}
    \draw[associazione] (ap0) -- (b);
  \end{tikzpicture}
  \end{adjustbox}
  \caption{Diagramma di UC49.1 -- Filtraggio per persona}\label{fig:uc:tr_f_persona}
\end{figure}

\begin{scheda}
Attore principale & Amministratore \\ \addlinespace
Precondizioni & Il registro è visualizzato. \\ \addlinespace
Postcondizioni & Sono mostrate le voci della persona scelta. \\ \addlinespace
Scenario principale & L'amministratore seleziona una persona dell'azienda. \\ \addlinespace
Incluso da & \hyperref[uc:tr_filtro]{UC49} \\ \addlinespace
Requisiti & \hyperlink{req:RF-92}{RF-92} \\
\end{scheda}

\subsubsubsection{UC49.2 -- Filtraggio per periodo}\label{uc:tr_f_periodo}
\begin{figure}[H]
  \centering
  \begin{adjustbox}{max width=\linewidth, max totalheight=0.62\textheight}
  \begin{tikzpicture}
    \node[ucfocus] (b) at (0.00,0.00) {\textbf{UC49.2}\\Filtraggio per periodo};
    \node[uc] (c0) at (7.15,0.00) {\textbf{UC49}\\Filtraggio del registro};
    \draw[dipendenza] (c0) -- node[etich, sloped] {\guillemotleft include\guillemotright} (b);
    \begin{scope}[on background layer]
      \node[confine, fit=(b) (c0)] (sistema) {};
    \end{scope}
    \node[nomesistema] at (sistema.north west) {NEXUM Concierge};
    \attore{ap0}{(-4.83,0.00)}{Amministratore}
    \draw[associazione] (ap0) -- (b);
  \end{tikzpicture}
  \end{adjustbox}
  \caption{Diagramma di UC49.2 -- Filtraggio per periodo}\label{fig:uc:tr_f_periodo}
\end{figure}

\begin{scheda}
Attore principale & Amministratore \\ \addlinespace
Precondizioni & Il registro è visualizzato. \\ \addlinespace
Postcondizioni & Sono mostrate le voci del periodo scelto. \\ \addlinespace
Scenario principale & L'amministratore indica data iniziale e finale. \\ \addlinespace
Incluso da & \hyperref[uc:tr_filtro]{UC49} \\ \addlinespace
Requisiti & \hyperlink{req:RF-93}{RF-93} \\
\end{scheda}

\subsubsubsection{UC49.3 -- Filtraggio per capacità ed esito}\label{uc:tr_f_capacita}
\begin{figure}[H]
  \centering
  \begin{adjustbox}{max width=\linewidth, max totalheight=0.62\textheight}
  \begin{tikzpicture}
    \node[ucfocus] (b) at (0.00,0.00) {\textbf{UC49.3}\\Filtraggio per capacità ed esito};
    \node[uc] (c0) at (7.15,0.00) {\textbf{UC49}\\Filtraggio del registro};
    \draw[dipendenza] (c0) -- node[etich, sloped] {\guillemotleft include\guillemotright} (b);
    \begin{scope}[on background layer]
      \node[confine, fit=(b) (c0)] (sistema) {};
    \end{scope}
    \node[nomesistema] at (sistema.north west) {NEXUM Concierge};
    \attore{ap0}{(-4.83,0.00)}{Amministratore}
    \draw[associazione] (ap0) -- (b);
  \end{tikzpicture}
  \end{adjustbox}
  \caption{Diagramma di UC49.3 -- Filtraggio per capacità ed esito}\label{fig:uc:tr_f_capacita}
\end{figure}

\begin{scheda}
Attore principale & Amministratore \\ \addlinespace
Precondizioni & Il registro è visualizzato. \\ \addlinespace
Postcondizioni & Sono mostrate le voci con la capacità o l'esito scelti. \\ \addlinespace
Scenario principale & L'amministratore seleziona una capacità e, se vuole, un esito. \\ \addlinespace
Incluso da & \hyperref[uc:tr_filtro]{UC49} \\ \addlinespace
Requisiti & \hyperlink{req:RF-94}{RF-94} \\
\end{scheda}

\subsubsection{UC50 -- Visualizzazione del dettaglio di un'interazione}\label{uc:tr_dettaglio}
\begin{figure}[H]
  \centering
  \begin{adjustbox}{max width=\linewidth, max totalheight=0.62\textheight}
  \begin{tikzpicture}
    \node[ucfocus] (b) at (0.00,0.00) {\textbf{UC50}\\Visualizzazione del dettaglio di un'interazione};
    \begin{scope}[on background layer]
      \node[confine, fit=(b)] (sistema) {};
    \end{scope}
    \node[nomesistema] at (sistema.north west) {NEXUM Concierge};
    \attore{ap0}{(-4.83,0.00)}{Amministratore}
    \draw[associazione] (ap0) -- (b);
  \end{tikzpicture}
  \end{adjustbox}
  \caption{Diagramma di UC50 -- Visualizzazione del dettaglio di un'interazione}\label{fig:uc:tr_dettaglio}
\end{figure}

\begin{scheda}
Attore principale & Amministratore \\ \addlinespace
Precondizioni & Il registro è visualizzato. \\ \addlinespace
Postcondizioni & L'interazione è ricostruibile. \\ \addlinespace
Scenario principale & \begin{enumerate}[nosep, leftmargin=*, before=\vspace{-\baselineskip}\vspace{2pt}, after=\vspace{-\baselineskip}]\item L'amministratore seleziona una voce. \item Il sistema mostra il messaggio, le fonti consultate con la versione, la capacità o abilità invocata con i dati inviati e, quando applicabili, la conferma e l'esito restituito dal servizio. \end{enumerate} \\ \addlinespace
Requisiti & \hyperlink{req:RF-95}{RF-95}, \hyperlink{req:RF-99}{RF-99} \\
\end{scheda}

\section{Requisiti}\label{sec:requisiti}

I requisiti traducono i casi d'uso, il capitolato e gli accordi con il
proponente in proprietà verificabili del sistema.

\subsection{Classificazione e codifica}
Ogni requisito ha un codice nella forma \textbf{R\textit{T}-\textit{nn}}, dove
\textit{nn} è un numero progressivo che non viene riassegnato e \textit{T} indica
il tipo:
\begin{itemize}
  \item \textbf{F}, \textbf{funzionale}: un comportamento che il sistema deve
        offrire;
  \item \textbf{Q}, \textbf{di qualità}: una caratteristica del prodotto o della
        documentazione da misurare o da garantire;
  \item \textbf{P}, \textbf{prestazionale}: un vincolo sui tempi, sull'uso delle
        risorse o sui costi;
  \item \textbf{V}, \textbf{di vincolo}: un limite imposto dal capitolato, dal
        proponente o dal contesto.
\end{itemize}
La priorità è indicata in una colonna a parte, così che possa cambiare senza
cambiare il codice:
\begin{itemize}
  \item \textbf{obbligatorio}: necessario per l'accettazione dell'MVP;
  \item \textbf{desiderabile}: atteso dal proponente, da soddisfare se il budget
        lo consente;
  \item \textbf{facoltativo}: utile ma non richiesto.
\end{itemize}
La colonna \textit{Fonti} indica i casi d'uso in cui il requisito interviene e
da dove nasce. Le origini possibili sono i requisiti di business del capitolato
(\textit{RB-nn}, con la numerazione del capitolato) e i criteri di accettazione
(\textit{CA-nn}, numerati nell'ordine in cui il capitolato li elenca). Sono
origini anche le sezioni del capitolato, i verbali degli incontri con il
proponente, il riscontro scritto del proponente e le proposte del gruppo
(\textit{Interno}). Le fonti esterne
sono elencate per intero nella Tabella~\ref{tab:tr-fonti}.

La colonna \textit{Ver.} indica come il requisito sarà verificato:
\begin{itemize}
  \item \textbf{T}, \textbf{test}: test di sistema o di accettazione;
  \item \textbf{A}, \textbf{analisi}: misura su un insieme di prova, con le soglie
        del \textit{Piano di Qualifica};
  \item \textbf{I}, \textbf{ispezione}: revisione del codice o dei documenti;
  \item \textbf{D}, \textbf{dimostrazione}: esecuzione completa davanti al
        proponente.
\end{itemize}

\subsection{Requisiti funzionali}\label{sec:rf}

I requisiti funzionali derivano dai casi d'uso e sono raggruppati per area
funzionale; l'attore coinvolto è quello dei casi d'uso indicati nella colonna
\textit{Fonti} (Tabella~\ref{tab:rf}).

{\small
\begin{longtable}{@{}>{\raggedright\arraybackslash}p{1.35cm}>{\raggedright\arraybackslash}p{7.3cm}>{\raggedright\arraybackslash}p{1.9cm}>{\centering\arraybackslash}p{0.75cm}>{\raggedright\arraybackslash}p{3.1cm}@{}}
\caption{Requisiti funzionali}\label{tab:rf}\\
\toprule
\intestazione{Codice} & \intestazione{Descrizione} & \intestazione{Priorità} & \intestazione{Ver.} & \intestazione{Fonti} \\
\midrule\endfirsthead
\toprule
\intestazione{Codice} & \intestazione{Descrizione} & \intestazione{Priorità} & \intestazione{Ver.} & \intestazione{Fonti} \\
\midrule\endhead
\midrule\multicolumn{5}{r@{}}{\footnotesize\itshape continua nella pagina successiva}\\ \endfoot \bottomrule\endlastfoot
\addlinespace[3pt]\multicolumn{5}{@{}l}{\textbf{\sffamily Accesso e profilo}} \\ \addlinespace[2pt]
\leavevmode\hypertarget{req:RF-01}{\textbf{RF-01}} & Il sistema deve riconoscere la persona che usa il Concierge e il suo ruolo (dipendente o amministratore) prima di consentire l’accesso alla chat o al backoffice. Il riconoscimento è gestito dall’applicativo, autonomo dal sistema di accesso di NEXUM. & Obbligatorio & T & \hyperref[uc:login]{UC1}, \mbox{RB-01}, \mbox{VE-2026-11-18}, \mbox{Riscontro Eggon} \\
\leavevmode\hypertarget{req:RF-02}{\textbf{RF-02}} & Il sistema deve associare a ogni utente autenticato l’azienda e l’utente NEXUM corrispondente, secondo la mappatura degli utenti di prova concordata con Eggon, e usare quell’identificativo nelle chiamate alle API Eggon. & Obbligatorio & T & \hyperref[uc:login]{UC1}, \mbox{RB-01}, \mbox{CA-09}, \mbox{VE-2026-11-18}, \mbox{Riscontro Eggon} \\
\leavevmode\hypertarget{req:RF-03}{\textbf{RF-03}} & Il sistema deve permettere l’inserimento dell’identificativo nella pagina di accesso. & Obbligatorio & T & \hyperref[uc:login_id]{UC1.1}, \mbox{Interno} \\
\leavevmode\hypertarget{req:RF-04}{\textbf{RF-04}} & Il sistema deve permettere l’inserimento della password senza mostrarla in chiaro. & Obbligatorio & T & \hyperref[uc:login_pw]{UC1.2}, \mbox{Interno} \\
\leavevmode\hypertarget{req:RF-05}{\textbf{RF-05}} & Il sistema deve segnalare le credenziali non valide con un messaggio che non riveli quale dato è errato. & Obbligatorio & T & \hyperref[uc:cred_err]{UC2}, \mbox{RB-14}, \mbox{Interno} \\
\leavevmode\hypertarget{req:RF-07}{\textbf{RF-07}} & Il sistema deve limitare conoscenze, dati e operazioni all’azienda a cui l’utente è associato. & Obbligatorio & T & \hyperref[uc:login]{UC1}, \mbox{RB-01}, \mbox{CA-09} \\
\leavevmode\hypertarget{req:RF-08}{\textbf{RF-08}} & Il sistema deve permettere all’utente autenticato di chiudere la sessione, eliminando dalla sessione il contesto della conversazione in corso; la traccia delle interazioni (RF-15) e il contenuto conservato secondo RF-99 non vengono cancellati. & Obbligatorio & T & \hyperref[uc:logout]{UC3}, \mbox{Interno} \\
\addlinespace[3pt]\multicolumn{5}{@{}l}{\textbf{\sffamily Conversazione}} \\ \addlinespace[2pt]
\leavevmode\hypertarget{req:RF-09}{\textbf{RF-09}} & Il sistema deve avviare una nuova conversazione, senza il contesto delle precedenti, a ogni apertura della chat e quando il dipendente lo chiede. & Obbligatorio & T & \hyperref[uc:nuova_conv]{UC4}, \mbox{RB-04}, \mbox{VE-2026-11-18}, \mbox{Riscontro Eggon} \\
\leavevmode\hypertarget{req:RF-10}{\textbf{RF-10}} & Il sistema deve mostrare i messaggi della conversazione in corso in ordine cronologico, rendendo riconoscibile il mittente di ciascuno. & Obbligatorio & T & \hyperref[uc:vis_conv]{UC5}, \mbox{Interno}, \mbox{Riscontro Eggon} \\
\leavevmode\hypertarget{req:RF-11}{\textbf{RF-11}} & Il sistema deve permettere al dipendente di inviare messaggi in linguaggio naturale attraverso una chat testuale. & Obbligatorio & T & \hyperref[uc:msg]{UC6}, \mbox{Capitolato} \\
\leavevmode\hypertarget{req:RF-12}{\textbf{RF-12}} & Il sistema deve interpretare ogni messaggio nel contesto della conversazione corrente, conservando le informazioni raccolte fino alla sua chiusura. & Obbligatorio & T & \hyperref[uc:msg]{UC6}, \mbox{RB-04}, \mbox{CA-05} \\
\leavevmode\hypertarget{req:RF-13}{\textbf{RF-13}} & Il sistema deve distinguere le richieste sulle regole aziendali, sulle istruzioni di NEXUM, sui dati personali e sulle operazioni, e scegliere di conseguenza la fonte o l’abilità. & Obbligatorio & A & \hyperref[uc:msg]{UC6}, \mbox{RB-02} \\
\leavevmode\hypertarget{req:RF-14}{\textbf{RF-14}} & Il sistema deve poter combinare più capacità nella stessa risposta, indicando per ogni parte la fonte utilizzata. & Obbligatorio & T & \hyperref[uc:msg]{UC6}, \mbox{RB-03}, \mbox{CA-13} \\
\leavevmode\hypertarget{req:RF-15}{\textbf{RF-15}} & Il sistema deve registrare ogni interazione con persona, azienda, data e ora, fonti utilizzate con le relative versioni, capacità o abilità invocata e, quando applicabili, conferma ed esito dell’operazione. & Obbligatorio & T & \hyperref[uc:msg]{UC6}, \mbox{RB-15}, \mbox{CA-12}, \mbox{Riscontro Eggon} \\
\leavevmode\hypertarget{req:RF-16}{\textbf{RF-16}} & Il sistema deve chiedere un chiarimento quando un messaggio ammette interpretazioni diverse su periodo, anno, tipo di richiesta o durata. & Obbligatorio & T & \hyperref[uc:chiarimento]{UC7}, \mbox{RB-04}, \mbox{CA-06} \\
\leavevmode\hypertarget{req:RF-17}{\textbf{RF-17}} & Il sistema deve completare una data priva di anno con l’anno già esplicitato nella conversazione oppure chiederlo al dipendente. & Obbligatorio & T & \hyperref[uc:chiarimento]{UC7}, \mbox{CA-06} \\
\leavevmode\hypertarget{req:RF-18}{\textbf{RF-18}} & Il sistema deve rispondere a una richiesta estranea al perimetro del servizio con un messaggio che ne spiega il limite e indica che cosa il Concierge può fare. & Obbligatorio & T & \hyperref[uc:fuori_ambito]{UC8}, \mbox{Vincoli}, \mbox{RB-14}, \mbox{VE-2026-11-18}, \mbox{Riscontro Eggon} \\
\leavevmode\hypertarget{req:RF-19}{\textbf{RF-19}} & Il sistema deve impedire l’invio di messaggi oltre la lunghezza massima, indicandola al dipendente. & Desiderabile & T & \hyperref[uc:msg_lungo]{UC9}, \mbox{Interno} \\
\leavevmode\hypertarget{req:RF-20}{\textbf{RF-20}} & Il sistema deve comunicare gli errori tecnici con un messaggio comprensibile, distinto dall’insufficienza di informazioni, e registrarne i dettagli nella traccia per l’amministratore. & Obbligatorio & T & \hyperref[uc:errore_tecnico]{UC10}, \mbox{RB-14}, \mbox{CA-11}, \mbox{VE-2026-11-18} \\
\addlinespace[3pt]\multicolumn{5}{@{}l}{\textbf{\sffamily Conoscenza aziendale}} \\ \addlinespace[2pt]
\leavevmode\hypertarget{req:RF-21}{\textbf{RF-21}} & Il sistema deve rispondere alle domande sulle regole aziendali usando le fonti caricate nel backoffice per l’azienda dell’utente. & Obbligatorio & A & \hyperref[uc:regole]{UC11}, \mbox{CA-01} \\
\leavevmode\hypertarget{req:RF-22}{\textbf{RF-22}} & Il sistema deve recuperare contenuti solo dalle fonti attive dei bucket dell’azienda dell’utente. & Obbligatorio & T & \hyperref[uc:regole]{UC11}, \mbox{RB-06}, \mbox{CA-09}, \mbox{VE-2026-11-18} \\
\leavevmode\hypertarget{req:RF-23}{\textbf{RF-23}} & Il sistema deve basare le risposte documentali esclusivamente sui contenuti recuperati. & Obbligatorio & A & \hyperref[uc:regole]{UC11}, \mbox{RB-08} \\
\leavevmode\hypertarget{req:RF-24}{\textbf{RF-24}} & Il sistema deve identificare per ogni risposta documentale le fonti utilizzate, con titolo, categoria e versione. & Obbligatorio & T & \hyperref[uc:fonti]{UC12}, \mbox{RB-07}, \mbox{RB-08}, \mbox{CA-07}, \mbox{Riscontro Eggon} \\
\leavevmode\hypertarget{req:RF-25}{\textbf{RF-25}} & Il sistema deve astenersi dal rispondere, dichiarandolo, quando le fonti disponibili sono insufficienti o in contrasto non risolvibile. & Obbligatorio & T & \hyperref[uc:astensione]{UC13}, \mbox{RB-08}, \mbox{CA-07} \\
\addlinespace[3pt]\multicolumn{5}{@{}l}{\textbf{\sffamily Istruzioni NEXUM}} \\ \addlinespace[2pt]
\leavevmode\hypertarget{req:RF-28}{\textbf{RF-28}} & Il sistema deve rispondere alle domande sull’uso delle funzioni di ferie e permessi sulla base dei manuali NEXUM. & Obbligatorio & A & \hyperref[uc:prodotto]{UC14}, \mbox{CA-02} \\
\leavevmode\hypertarget{req:RF-29}{\textbf{RF-29}} & Il sistema deve mantenere distinte le istruzioni dei manuali NEXUM dalle regole aziendali, nelle fonti e nelle risposte. & Obbligatorio & T & \hyperref[uc:prodotto]{UC14}, \mbox{RB-07} \\
\addlinespace[3pt]\multicolumn{5}{@{}l}{\textbf{\sffamily Lettura dei dati}} \\ \addlinespace[2pt]
\leavevmode\hypertarget{req:RF-31}{\textbf{RF-31}} & Il sistema deve recuperare i dati personali di ferie e permessi solo tramite l’abilità di lettura. & Obbligatorio & T & \hyperref[uc:dati]{UC15}, \mbox{RB-09}, \mbox{CA-03} \\
\leavevmode\hypertarget{req:RF-32}{\textbf{RF-32}} & Il sistema deve presentare conteggi e valori così come restituiti dal servizio, senza che il modello li calcoli o li deduca. & Obbligatorio & T & \hyperref[uc:dati]{UC15}, \mbox{Capitolato}, \mbox{CA-03} \\
\leavevmode\hypertarget{req:RF-33}{\textbf{RF-33}} & Il sistema deve distinguere nei risultati il periodo delle richieste (passate o future) dal loro stato (per esempio approvata, in attesa o rifiutata), usando il significato dei dati concordato con Eggon. & Obbligatorio & T & \hyperref[uc:dati_riep]{UC15.1}, \mbox{RB-10}, \mbox{CA-03}, \mbox{Riscontro Eggon} \\
\leavevmode\hypertarget{req:RF-34}{\textbf{RF-34}} & Il sistema deve specificare nei risultati sui dati il periodo, l’unità di misura e il significato degli stati considerati. & Obbligatorio & T & \hyperref[uc:dati_criteri]{UC15.2}, \mbox{RB-10} \\
\leavevmode\hypertarget{req:RF-35}{\textbf{RF-35}} & Il sistema deve mostrare il saldo residuo delle ferie quando il servizio lo fornisce, con il riferimento al suo aggiornamento. & Obbligatorio & T & \hyperref[uc:saldo]{UC16}, \mbox{VE-2026-10-26}, \mbox{Riscontro Eggon} \\
\leavevmode\hypertarget{req:RF-36}{\textbf{RF-36}} & Il sistema non deve ricostruire autonomamente il saldo residuo quando il servizio non lo fornisce. & Obbligatorio & T & \hyperref[uc:saldo]{UC16}, \mbox{Capitolato}, \mbox{Riscontro Eggon} \\
\leavevmode\hypertarget{req:RF-37}{\textbf{RF-37}} & Il sistema deve dichiarare l’assenza di dati per il periodo richiesto, distinguendola da un errore tecnico. & Obbligatorio & T & \hyperref[uc:nessun_dato]{UC17}, \mbox{RB-14} \\
\leavevmode\hypertarget{req:RF-38}{\textbf{RF-38}} & Il sistema deve rispondere con un messaggio di rifiuto alle richieste su dati o operazioni di altre persone o aziende, senza mostrare alcun dato di altri. & Obbligatorio & T & \hyperref[uc:non_aut]{UC18}, \mbox{RB-01}, \mbox{CA-09} \\
\leavevmode\hypertarget{req:RF-39}{\textbf{RF-39}} & Il sistema deve consultare e mostrare lo stato effettivo in NEXUM di una richiesta del dipendente, anche dopo l’inserimento, distinguendo per esempio registrata, in attesa di approvazione e approvata. & Obbligatorio & T & \hyperref[uc:stato_rich]{UC19}, \mbox{RB-13}, \mbox{CA-10}, \mbox{VE-2026-11-18}, \mbox{Riscontro Eggon} \\
\addlinespace[3pt]\multicolumn{5}{@{}l}{\textbf{\sffamily Richieste di ferie e permessi}} \\ \addlinespace[2pt]
\leavevmode\hypertarget{req:RF-40}{\textbf{RF-40}} & Il sistema deve permettere al dipendente di preparare e registrare una richiesta di ferie o permesso tramite la chat. & Obbligatorio & T & \hyperref[uc:registra]{UC20}, \mbox{CA-04} \\
\leavevmode\hypertarget{req:RF-41}{\textbf{RF-41}} & Il sistema deve trattare la richiesta in linguaggio naturale come avvio della preparazione, senza inviarla. & Obbligatorio & T & \hyperref[uc:registra]{UC20}, \mbox{RB-12} \\
\leavevmode\hypertarget{req:RF-42}{\textbf{RF-42}} & Il sistema deve eseguire le operazioni solo attraverso abilità definite e servizi autorizzati. & Obbligatorio & I & \hyperref[uc:registra]{UC20}, \mbox{RB-11} \\
\leavevmode\hypertarget{req:RF-43}{\textbf{RF-43}} & Per ogni operazione sulle richieste, oltre ai dati registrati per ogni interazione (RF-15), il sistema deve registrare sempre la conferma del dipendente e l’esito restituito dal servizio. & Obbligatorio & T & \hyperref[uc:registra]{UC20}, \mbox{RB-15}, \mbox{CA-12}, \mbox{Riscontro Eggon} \\
\leavevmode\hypertarget{req:RF-44}{\textbf{RF-44}} & Il sistema deve determinare il tipo di richiesta tra quelli previsti dal servizio NEXUM. & Obbligatorio & T & \hyperref[uc:reg_tipo]{UC20.1}, \mbox{Capitolato} \\
\leavevmode\hypertarget{req:RF-45}{\textbf{RF-45}} & Il sistema deve determinare date di inizio e fine e, per i permessi, orari o durata, completi di anno. & Obbligatorio & T & \hyperref[uc:periodo]{UC21}, \mbox{Capitolato}, \mbox{CA-06} \\
\leavevmode\hypertarget{req:RF-46}{\textbf{RF-46}} & Il sistema deve raccogliere tutte le informazioni obbligatorie previste dal servizio NEXUM prima del riepilogo. & Obbligatorio & T & \hyperref[uc:reg_campi]{UC20.2}, \mbox{Capitolato} \\
\leavevmode\hypertarget{req:RF-47}{\textbf{RF-47}} & Il sistema deve mostrare prima di ogni invio un riepilogo completo, con date e orari per esteso. & Obbligatorio & T & \hyperref[uc:op_riep]{UC22}, \mbox{RB-12}, \mbox{CA-04} \\
\leavevmode\hypertarget{req:RF-48}{\textbf{RF-48}} & Il sistema deve inviare un’operazione solo dopo una conferma esplicita del dipendente sul riepilogo corrente. & Obbligatorio & T & \hyperref[uc:op_conf]{UC23}, \mbox{RB-12}, \mbox{CA-04}, \mbox{Vincoli} \\
\leavevmode\hypertarget{req:RF-49}{\textbf{RF-49}} & Il sistema deve mostrare l’esito effettivo restituito dal servizio, distinguendo richiesta registrata, in attesa di approvazione e approvata. & Obbligatorio & T & \hyperref[uc:op_esito]{UC24}, \mbox{RB-13}, \mbox{CA-10} \\
\leavevmode\hypertarget{req:RF-50}{\textbf{RF-50}} & Il sistema deve richiedere un nuovo riepilogo e una nuova conferma quando cambiano i dati da inviare. & Obbligatorio & T & \hyperref[uc:op_modifica]{UC25}, \mbox{RB-12}, \mbox{CA-10} \\
\leavevmode\hypertarget{req:RF-51}{\textbf{RF-51}} & Il sistema deve permettere di rinunciare a un’operazione prima della conferma, senza inviare nulla. & Obbligatorio & T & \hyperref[uc:op_rinuncia]{UC26}, \mbox{RB-12} \\
\leavevmode\hypertarget{req:RF-52}{\textbf{RF-52}} & Il sistema deve comunicare il motivo delle validazioni negative restituite dal servizio NEXUM. & Obbligatorio & T & \hyperref[uc:op_rifiuto]{UC27}, \mbox{RB-14}, \mbox{CA-11}, \mbox{Vincoli} \\
\leavevmode\hypertarget{req:RF-53}{\textbf{RF-53}} & Il sistema deve verificare l’esito di un invio interrotto prima di ripeterlo, evitando operazioni duplicate. & Obbligatorio & T & \hyperref[uc:op_incerto]{UC28}, \mbox{Capitolato}, \mbox{CA-11} \\
\leavevmode\hypertarget{req:RF-54}{\textbf{RF-54}} & Il sistema deve permettere al dipendente di cancellare una propria richiesta, con riepilogo e conferma. La correzione di una richiesta avviene tramite cancellazione e nuovo inserimento. & Obbligatorio & T & \hyperref[uc:annulla_inv]{UC29}, \mbox{VE-2026-10-26}, \mbox{VE-2026-11-18}, \mbox{Riscontro Eggon} \\
\leavevmode\hypertarget{req:RF-55}{\textbf{RF-55}} & Il sistema deve individuare in modo conversazionale la richiesta a cui si riferisce il messaggio, per esempio da date o tipo, chiedendo quale quando ve ne sono più. & Desiderabile & T & \hyperref[uc:sel_rich]{UC30}, \mbox{VE-2026-10-26}, \mbox{Riscontro Eggon} \\
\addlinespace[3pt]\multicolumn{5}{@{}l}{\textbf{\sffamily Gestione dei bucket}} \\ \addlinespace[2pt]
\leavevmode\hypertarget{req:RF-57}{\textbf{RF-57}} & Il sistema deve mostrare agli amministratori tutti i bucket dell’azienda. & Obbligatorio & T & \hyperref[uc:bk_elenco]{UC31}, \mbox{RB-05}, \mbox{VE-2026-11-18}, \mbox{Riscontro Eggon} \\
\leavevmode\hypertarget{req:RF-58}{\textbf{RF-58}} & Il sistema deve mostrare per ogni bucket nome, descrizione, categoria e numero di fonti per stato. & Obbligatorio & T & \hyperref[uc:bk_elem]{UC31.1}, \mbox{Capitolato}, \mbox{VE-2026-11-18}, \mbox{Riscontro Eggon} \\
\leavevmode\hypertarget{req:RF-59}{\textbf{RF-59}} & Il sistema deve permettere all’amministratore di creare un bucket inserendo nome, descrizione della finalità (campo di testo libero) e categoria. & Obbligatorio & T & \hyperref[uc:bk_crea]{UC32}, \hyperref[uc:bk_nome]{UC32.1}, \hyperref[uc:bk_desc]{UC32.2}, \mbox{RB-05}, \mbox{Capitolato}, \mbox{VE-2026-11-18}, \mbox{Riscontro Eggon} \\
\leavevmode\hypertarget{req:RF-60}{\textbf{RF-60}} & Il sistema deve rendere riconoscibile la categoria di ogni bucket, conoscenza aziendale o conoscenza del prodotto, scelta dall’amministratore alla creazione. & Obbligatorio & T & \hyperref[uc:bk_crea]{UC32}, \mbox{RB-07}, \mbox{Riscontro Eggon}, \mbox{Interno} \\
\leavevmode\hypertarget{req:RF-63}{\textbf{RF-63}} & Il sistema deve rendere consultabili le conoscenze di ogni bucket attivo a tutti i dipendenti dell’azienda, senza segmentazione per gruppi, dipartimenti o sedi. & Obbligatorio & T & \hyperref[uc:bk_crea]{UC32}, \mbox{CA-09}, \mbox{VE-2026-11-18}, \mbox{Riscontro Eggon} \\
\leavevmode\hypertarget{req:RF-64}{\textbf{RF-64}} & Il sistema deve impedire due bucket con lo stesso nome nella stessa azienda. & Desiderabile & T & \hyperref[uc:bk_nome_dup]{UC34}, \mbox{Interno} \\
\leavevmode\hypertarget{req:RF-65}{\textbf{RF-65}} & Il sistema deve impedire il salvataggio di un bucket senza nome o senza descrizione, evidenziando i campi mancanti. & Desiderabile & T & \hyperref[uc:bk_incompleto]{UC35}, \mbox{Interno}, \mbox{VE-2026-11-18}, \mbox{Riscontro Eggon} \\
\leavevmode\hypertarget{req:RF-66}{\textbf{RF-66}} & Il sistema deve permettere di modificare nome e descrizione di un bucket. & Desiderabile & T & \hyperref[uc:bk_modifica]{UC36}, \mbox{RB-05}, \mbox{VE-2026-11-18}, \mbox{Riscontro Eggon} \\
\leavevmode\hypertarget{req:RF-67}{\textbf{RF-67}} & Il sistema deve permettere di eliminare un bucket, previa conferma, ritirando tutte le fonti che contiene. & Obbligatorio & T & \hyperref[uc:bk_elimina]{UC37}, \mbox{RB-05} \\
\addlinespace[3pt]\multicolumn{5}{@{}l}{\textbf{\sffamily Gestione delle fonti}} \\ \addlinespace[2pt]
\leavevmode\hypertarget{req:RF-68}{\textbf{RF-68}} & Il sistema deve mostrare le fonti contenute in un bucket. & Obbligatorio & T & \hyperref[uc:fn_elenco]{UC38}, \mbox{RB-05} \\
\leavevmode\hypertarget{req:RF-69}{\textbf{RF-69}} & Il sistema deve mostrare per ogni fonte titolo, versione, data di caricamento, stato di elaborazione e disponibilità. & Obbligatorio & T & \hyperref[uc:fn_elem]{UC38.1}, \mbox{Capitolato} \\
\leavevmode\hypertarget{req:RF-70}{\textbf{RF-70}} & Il sistema deve permettere all’amministratore di caricare documenti PDF testuali in un bucket. & Obbligatorio & T & \hyperref[uc:fn_carica]{UC39}, \mbox{RB-05}, \mbox{CA-01} \\
\leavevmode\hypertarget{req:RF-71}{\textbf{RF-71}} & Il sistema deve elaborare le fonti in modo asincrono, rendendone consultabile lo stato. & Obbligatorio & T & \hyperref[uc:fn_carica]{UC39}, \mbox{Prestazioni} \\
\leavevmode\hypertarget{req:RF-72}{\textbf{RF-72}} & Il sistema deve trasformare ogni fonte in porzioni ricercabili che conservano provenienza (fonte e bucket) e versione. & Obbligatorio & I & \hyperref[uc:fn_carica]{UC39}, \mbox{RB-06} \\
\leavevmode\hypertarget{req:RF-73}{\textbf{RF-73}} & Il sistema deve permettere di selezionare il file da caricare dal dispositivo dell’amministratore. & Obbligatorio & T & \hyperref[uc:fn_file]{UC39.1}, \mbox{Interno} \\
\leavevmode\hypertarget{req:RF-74}{\textbf{RF-74}} & Il sistema deve permettere di associare a ogni fonte titolo e versione. & Obbligatorio & T & \hyperref[uc:fn_meta]{UC39.2}, \mbox{RB-06} \\
\leavevmode\hypertarget{req:RF-75}{\textbf{RF-75}} & Il sistema deve permettere di associare a una fonte data di entrata in vigore e condizioni di applicabilità. & Desiderabile & T & \hyperref[uc:fn_meta]{UC39.2}, \mbox{Capitolato} \\
\leavevmode\hypertarget{req:RF-76}{\textbf{RF-76}} & Il sistema deve rifiutare i file non in formato PDF, indicando il formato accettato. & Obbligatorio & T & \hyperref[uc:fn_formato]{UC40}, \mbox{Vincoli} \\
\leavevmode\hypertarget{req:RF-77}{\textbf{RF-77}} & Il sistema deve rifiutare i file oltre la dimensione massima, indicandola. & Desiderabile & T & \hyperref[uc:fn_dim]{UC41}, \mbox{Interno} \\
\leavevmode\hypertarget{req:RF-78}{\textbf{RF-78}} & Il sistema deve segnalare il caricamento di un file già presente nel bucket. & Desiderabile & T & \hyperref[uc:fn_dup]{UC42}, \mbox{Interno} \\
\leavevmode\hypertarget{req:RF-80}{\textbf{RF-80}} & Il sistema deve mostrare i metadati di una fonte, compresi autore e data del caricamento. & Obbligatorio & T & \hyperref[uc:fn_meta_v]{UC43.1}, \mbox{RB-06} \\
\leavevmode\hypertarget{req:RF-81}{\textbf{RF-81}} & Il sistema deve mostrare lo stato di elaborazione di una fonte e la sua disponibilità alla consultazione. & Obbligatorio & T & \hyperref[uc:fn_stato]{UC43.2}, \mbox{Capitolato} \\
\leavevmode\hypertarget{req:RF-83}{\textbf{RF-83}} & Il sistema deve mostrare all’amministratore il motivo degli errori di elaborazione di una fonte, che avviene al caricamento. & Desiderabile & T & \hyperref[uc:fn_errore]{UC44}, \mbox{Capitolato}, \mbox{RB-14}, \mbox{VE-2026-11-18} \\
\leavevmode\hypertarget{req:RF-84}{\textbf{RF-84}} & Il sistema deve permettere di sostituire una fonte con una nuova versione, conservando la precedente nello storico, così da poter ricostruire quale versione abbia supportato un’interazione precedente. & Obbligatorio & T & \hyperref[uc:fn_sostituisci]{UC45}, \mbox{RB-05}, \mbox{RB-06}, \mbox{CA-08}, \mbox{VE-2026-11-18}, \mbox{Riscontro Eggon} \\
\leavevmode\hypertarget{req:RF-85}{\textbf{RF-85}} & Il sistema deve usare nelle risposte successive la nuova versione di una fonte sostituita, dal momento in cui è disponibile ed entro il tempo concordato con Eggon. & Obbligatorio & T & \hyperref[uc:fn_sostituisci]{UC45}, \mbox{Capitolato}, \mbox{CA-08}, \mbox{VE-2026-11-18}, \mbox{Riscontro Eggon} \\
\leavevmode\hypertarget{req:RF-86}{\textbf{RF-86}} & Il sistema deve permettere di ritirare una fonte, previa conferma, escludendola dal recupero entro il tempo concordato con Eggon. & Obbligatorio & T & \hyperref[uc:fn_ritira]{UC46}, \mbox{RB-05}, \mbox{CA-08} \\
\leavevmode\hypertarget{req:RF-87}{\textbf{RF-87}} & Il sistema deve permettere di ripetere l’elaborazione di una fonte non riuscita. & Desiderabile & T & \hyperref[uc:fn_riprova]{UC47}, \mbox{Interno} \\
\addlinespace[3pt]\multicolumn{5}{@{}l}{\textbf{\sffamily Tracciabilità}} \\ \addlinespace[2pt]
\leavevmode\hypertarget{req:RF-89}{\textbf{RF-89}} & Il sistema deve permettere agli amministratori di consultare nel backoffice la traccia delle interazioni della propria azienda; nessun altro profilo può accedervi. & Obbligatorio & T & \hyperref[uc:tr_elenco]{UC48}, \mbox{RB-15}, \mbox{CA-12}, \mbox{Riscontro Eggon} \\
\leavevmode\hypertarget{req:RF-90}{\textbf{RF-90}} & Il sistema deve mostrare nell’elenco del registro, per ogni voce, data e ora, persona, azienda, capacità o abilità invocata e, quando applicabili, conferma ed esito. & Obbligatorio & T & \hyperref[uc:tr_voce]{UC48.1}, \mbox{RB-15}, \mbox{VE-2026-11-18}, \mbox{Riscontro Eggon} \\
\leavevmode\hypertarget{req:RF-91}{\textbf{RF-91}} & Il sistema deve permettere di filtrare il registro delle interazioni. & Facoltativo & T & \hyperref[uc:tr_filtro]{UC49}, \mbox{Interno} \\
\leavevmode\hypertarget{req:RF-92}{\textbf{RF-92}} & Il sistema deve permettere di filtrare il registro per persona. & Facoltativo & T & \hyperref[uc:tr_f_persona]{UC49.1}, \mbox{Interno} \\
\leavevmode\hypertarget{req:RF-93}{\textbf{RF-93}} & Il sistema deve permettere di filtrare il registro per periodo. & Facoltativo & T & \hyperref[uc:tr_f_periodo]{UC49.2}, \mbox{Interno} \\
\leavevmode\hypertarget{req:RF-94}{\textbf{RF-94}} & Il sistema deve permettere di filtrare il registro per capacità ed esito. & Facoltativo & T & \hyperref[uc:tr_f_capacita]{UC49.3}, \mbox{Interno} \\
\leavevmode\hypertarget{req:RF-95}{\textbf{RF-95}} & Il sistema deve mostrare nel dettaglio di ogni interazione le fonti consultate con la versione, la capacità o abilità invocata e, quando applicabili, conferma ed esito. & Obbligatorio & T & \hyperref[uc:tr_dettaglio]{UC50}, \mbox{RB-15}, \mbox{CA-12}, \mbox{VE-2026-11-18}, \mbox{Riscontro Eggon} \\
\addlinespace[3pt]\multicolumn{5}{@{}l}{\textbf{\sffamily Accesso e profilo}} \\ \addlinespace[2pt]
\leavevmode\hypertarget{req:RF-96}{\textbf{RF-96}} & Il sistema deve autenticare con OAuth, secondo il contratto definito con Eggon, ogni chiamata verso le API Eggon, senza esporre le credenziali al modello né al dipendente. & Obbligatorio & T & \hyperref[uc:dati]{UC15}, \hyperref[uc:stato_rich]{UC19}, \hyperref[uc:registra]{UC20}, \hyperref[uc:annulla_inv]{UC29}, \mbox{VE-2026-11-18}, \mbox{Riscontro Eggon} \\
\addlinespace[3pt]\multicolumn{5}{@{}l}{\textbf{\sffamily Conversazione}} \\ \addlinespace[2pt]
\leavevmode\hypertarget{req:RF-97}{\textbf{RF-97}} & Il sistema deve mostrare nell’interfaccia della chat l’orario di ogni messaggio. & Desiderabile & T & \hyperref[uc:vis_conv]{UC5}, \mbox{VE-2026-11-18}, \mbox{Riscontro Eggon} \\
\leavevmode\hypertarget{req:RF-98}{\textbf{RF-98}} & Il sistema deve mostrare all’apertura della chat un avviso che spiega che cosa può fare il Concierge e invita a non inserire informazioni riservate. & Desiderabile & T & \hyperref[uc:nuova_conv]{UC4}, \mbox{VE-2026-11-18} \\
\addlinespace[3pt]\multicolumn{5}{@{}l}{\textbf{\sffamily Tracciabilità}} \\ \addlinespace[2pt]
\leavevmode\hypertarget{req:RF-99}{\textbf{RF-99}} & Il sistema deve conservare il contenuto delle conversazioni, separato dall’interfaccia del dipendente, per la verifica e il miglioramento della qualità delle risposte. & Desiderabile & T & \hyperref[uc:msg]{UC6}, \hyperref[uc:tr_dettaglio]{UC50}, \mbox{VE-2026-11-18}, \mbox{Riscontro Eggon} \\
\addlinespace[3pt]\multicolumn{5}{@{}l}{\textbf{\sffamily Conoscenza aziendale}} \\ \addlinespace[2pt]
\leavevmode\hypertarget{req:RF-100}{\textbf{RF-100}} & Il sistema deve indicare per ogni fonte citata la pagina o la sezione utilizzata per la risposta. & Desiderabile & T & \hyperref[uc:fonti]{UC12}, \mbox{VE-2026-11-18}, \mbox{Riscontro Eggon} \\
\addlinespace[3pt]\multicolumn{5}{@{}l}{\textbf{\sffamily Gestione dei bucket}} \\ \addlinespace[2pt]
\leavevmode\hypertarget{req:RF-101}{\textbf{RF-101}} & Il sistema deve permettere di limitare la consultazione di un bucket a gruppi di dipendenti e di modificarli, applicando i nuovi destinatari al recupero successivo. & Desiderabile & T & \hyperref[uc:bk_dest]{UC33}, \hyperref[uc:bk_modifica]{UC36}, \mbox{Capitolato}, \mbox{CA-09}, \mbox{VE-2026-11-18}, \mbox{Riscontro Eggon} \\
\addlinespace[3pt]\multicolumn{5}{@{}l}{\textbf{\sffamily Richieste di ferie e permessi}} \\ \addlinespace[2pt]
\leavevmode\hypertarget{req:RF-102}{\textbf{RF-102}} & Il sistema non deve eseguire un’operazione su una richiesta non identificata in modo univoco: prima di cancellarla deve identificarla e ottenere la conferma del dipendente. & Obbligatorio & T & \hyperref[uc:annulla_inv]{UC29}, \hyperref[uc:sel_rich]{UC30}, \mbox{Riscontro Eggon} \\
\end{longtable}}

\subsection{Requisiti di qualità}\label{sec:rq}

Il capitolato chiede di valutare separatamente il recupero, la generazione della
risposta e l'esecuzione delle operazioni, e rinvia le soglie quantitative ai
volumi di prova: le soglie saranno fissate nel \textit{Piano di Qualifica} e
concordate con il proponente (Tabella~\ref{tab:rq}).

{\small
\begin{longtable}{@{}>{\raggedright\arraybackslash}p{1.35cm}>{\raggedright\arraybackslash}p{7.3cm}>{\raggedright\arraybackslash}p{1.9cm}>{\centering\arraybackslash}p{0.75cm}>{\raggedright\arraybackslash}p{3.1cm}@{}}
\caption{Requisiti di qualità}\label{tab:rq}\\
\toprule
\intestazione{Codice} & \intestazione{Descrizione} & \intestazione{Priorità} & \intestazione{Ver.} & \intestazione{Fonti} \\
\midrule\endfirsthead
\toprule
\intestazione{Codice} & \intestazione{Descrizione} & \intestazione{Priorità} & \intestazione{Ver.} & \intestazione{Fonti} \\
\midrule\endhead
\midrule\multicolumn{5}{r@{}}{\footnotesize\itshape continua nella pagina successiva}\\ \endfoot \bottomrule\endlastfoot
\leavevmode\hypertarget{req:RQ-01}{\textbf{RQ-01}} & La qualità delle risposte documentali deve essere misurata come correttezza delle risposte e delle citazioni, con soglie definite nel Piano di Qualifica. & Obbligatorio & A & \hyperref[uc:regole]{UC11}, \hyperref[uc:fonti]{UC12}, \mbox{Prestazioni}, \mbox{RB-08} \\
\leavevmode\hypertarget{req:RQ-02}{\textbf{RQ-02}} & La pertinenza del recupero deve essere misurata separatamente dalla generazione della risposta. & Obbligatorio & A & \hyperref[uc:regole]{UC11}, \mbox{Prestazioni} \\
\leavevmode\hypertarget{req:RQ-03}{\textbf{RQ-03}} & La scelta della capacità (conoscenza o abilità applicative) deve essere misurata su un insieme di domande di prova. & Obbligatorio & A & \hyperref[uc:msg]{UC6}, \mbox{Prestazioni}, \mbox{RB-02} \\
\leavevmode\hypertarget{req:RQ-04}{\textbf{RQ-04}} & Il mantenimento del contesto nella conversazione corrente e la richiesta di chiarimenti devono essere verificati con dialoghi di prova a più messaggi. & Obbligatorio & A & \hyperref[uc:msg]{UC6}, \hyperref[uc:chiarimento]{UC7}, \mbox{Prestazioni}, \mbox{CA-05}, \mbox{CA-06} \\
\leavevmode\hypertarget{req:RQ-05}{\textbf{RQ-05}} & La correttezza dei dati restituiti e il completamento delle operazioni devono essere misurati separatamente dalle risposte documentali. & Obbligatorio & A & \hyperref[uc:dati]{UC15}, \hyperref[uc:registra]{UC20}, \mbox{Prestazioni} \\
\leavevmode\hypertarget{req:RQ-06}{\textbf{RQ-06}} & La gestione di errori e tentativi ripetuti deve essere verificata con prove che simulano indisponibilità, validazioni negative e interruzioni. & Obbligatorio & T & \hyperref[uc:errore_tecnico]{UC10}, \hyperref[uc:op_rifiuto]{UC27}, \hyperref[uc:op_incerto]{UC28}, \mbox{Prestazioni}, \mbox{CA-11} \\
\leavevmode\hypertarget{req:RQ-07}{\textbf{RQ-07}} & La completezza dell’elaborazione delle fonti deve essere misurata sui documenti di prova. & Obbligatorio & A & \hyperref[uc:fn_carica]{UC39}, \mbox{Prestazioni} \\
\leavevmode\hypertarget{req:RQ-08}{\textbf{RQ-08}} & Domande, dialoghi, dati e scenari di prova devono essere conservati e rieseguibili, così che la valutazione sia riproducibile. & Obbligatorio & I & \mbox{Deliverables} \\
\leavevmode\hypertarget{req:RQ-09}{\textbf{RQ-09}} & Deve essere prodotto un report di valutazione con i risultati delle misure. & Obbligatorio & I & \mbox{Deliverables} \\
\leavevmode\hypertarget{req:RQ-10}{\textbf{RQ-10}} & L’architettura deve permettere di aggiungere fonti e abilità senza riprogettare il Concierge, con interfacce delle abilità documentate. & Obbligatorio & I & \mbox{RB-16}, \mbox{Riscontro Eggon} \\
\leavevmode\hypertarget{req:RQ-11}{\textbf{RQ-11}} & La documentazione e la dimostrazione devono identificare chiaramente le componenti reali e quelle simulate. & Obbligatorio & I & \mbox{Vincoli} \\
\leavevmode\hypertarget{req:RQ-12}{\textbf{RQ-12}} & Devono essere consegnati codice sorgente versionato, istruzioni di installazione, Manuale Utente e documentazione tecnica. & Obbligatorio & I & \mbox{Deliverables} \\
\leavevmode\hypertarget{req:RQ-13}{\textbf{RQ-13}} & Devono essere documentati i limiti noti e le indicazioni per le estensioni. & Obbligatorio & I & \mbox{Deliverables} \\
\leavevmode\hypertarget{req:RQ-15}{\textbf{RQ-15}} & Il codice deve essere sottoposto a revisione di sicurezza con riferimento alle linee guida OWASP. & Desiderabile & I & \mbox{VE-2026-10-26} \\
\leavevmode\hypertarget{req:RQ-16}{\textbf{RQ-16}} & Le attività di sviluppo e documentazione devono rispettare le Norme di Progetto. & Obbligatorio & I & \mbox{Interno} \\
\end{longtable}}

\subsection{Requisiti prestazionali}\label{sec:rp}

I requisiti prestazionali riguardano tempi e consumi (Tabella~\ref{tab:rp}).
Anche in questo caso le soglie numeriche saranno definite nel \textit{Piano di
Qualifica}.

{\small
\begin{longtable}{@{}>{\raggedright\arraybackslash}p{1.35cm}>{\raggedright\arraybackslash}p{7.3cm}>{\raggedright\arraybackslash}p{1.9cm}>{\centering\arraybackslash}p{0.75cm}>{\raggedright\arraybackslash}p{3.1cm}@{}}
\caption{Requisiti prestazionali}\label{tab:rp}\\
\toprule
\intestazione{Codice} & \intestazione{Descrizione} & \intestazione{Priorità} & \intestazione{Ver.} & \intestazione{Fonti} \\
\midrule\endfirsthead
\toprule
\intestazione{Codice} & \intestazione{Descrizione} & \intestazione{Priorità} & \intestazione{Ver.} & \intestazione{Fonti} \\
\midrule\endhead
\midrule\multicolumn{5}{r@{}}{\footnotesize\itshape continua nella pagina successiva}\\ \endfoot \bottomrule\endlastfoot
\leavevmode\hypertarget{req:RP-01}{\textbf{RP-01}} & Le elaborazioni documentali lunghe devono avvenire in modo asincrono, senza bloccare il backoffice. & Obbligatorio & T & \hyperref[uc:fn_carica]{UC39}, \mbox{Prestazioni} \\
\leavevmode\hypertarget{req:RP-02}{\textbf{RP-02}} & I tempi di risposta della chat devono essere misurati per capacità e rispettare le soglie definite nel Piano di Qualifica. & Obbligatorio & A & \hyperref[uc:msg]{UC6}, \mbox{Prestazioni} \\
\leavevmode\hypertarget{req:RP-03}{\textbf{RP-03}} & Il tempo tra la sostituzione o il ritiro di una fonte e il suo effetto sulle risposte deve essere misurato e rispettare la soglia concordata con Eggon. & Obbligatorio & A & \hyperref[uc:fn_sostituisci]{UC45}, \hyperref[uc:fn_ritira]{UC46}, \mbox{Prestazioni}, \mbox{CA-08}, \mbox{Riscontro Eggon} \\
\leavevmode\hypertarget{req:RP-04}{\textbf{RP-04}} & Il consumo di token e il costo dei servizi AI devono essere misurati per interazione e restare entro il budget concordato con il proponente. & Obbligatorio & A & \mbox{Prestazioni}, \mbox{VE-2026-10-26} \\
\end{longtable}}

\subsection{Requisiti di vincolo}\label{sec:rv}

I requisiti di vincolo raccolgono i vincoli del capitolato e gli accordi con il
proponente (Tabella~\ref{tab:rv}). Lo \glo{Stack tecnologico} di Eggon è
consigliato e non obbligatorio, mentre l'indipendenza dal fornitore dei
modelli è l'unico vincolo tecnologico tassativo emerso dall'incontro.

{\small
\begin{longtable}{@{}>{\raggedright\arraybackslash}p{1.35cm}>{\raggedright\arraybackslash}p{7.3cm}>{\raggedright\arraybackslash}p{1.9cm}>{\centering\arraybackslash}p{0.75cm}>{\raggedright\arraybackslash}p{3.1cm}@{}}
\caption{Requisiti di vincolo}\label{tab:rv}\\
\toprule
\intestazione{Codice} & \intestazione{Descrizione} & \intestazione{Priorità} & \intestazione{Ver.} & \intestazione{Fonti} \\
\midrule\endfirsthead
\toprule
\intestazione{Codice} & \intestazione{Descrizione} & \intestazione{Priorità} & \intestazione{Ver.} & \intestazione{Fonti} \\
\midrule\endhead
\midrule\multicolumn{5}{r@{}}{\footnotesize\itshape continua nella pagina successiva}\\ \endfoot \bottomrule\endlastfoot
\leavevmode\hypertarget{req:RV-01}{\textbf{RV-01}} & Il dominio operativo dell’MVP deve essere limitato a ferie e permessi. & Obbligatorio & I & \mbox{Vincoli}, \mbox{VE-2026-10-26} \\
\leavevmode\hypertarget{req:RV-02}{\textbf{RV-02}} & L’interfaccia conversazionale deve essere una chat testuale, senza voce né allegati. & Obbligatorio & I & \hyperref[uc:msg]{UC6}, \mbox{Vincoli} \\
\leavevmode\hypertarget{req:RV-03}{\textbf{RV-03}} & La memoria personale tra conversazioni diverse e uno storico di chat che il dipendente possa riaprire e proseguire non sono richiesti nell’MVP: il contesto vale solo nella conversazione corrente. & Obbligatorio & I & \hyperref[uc:nuova_conv]{UC4}, \mbox{Vincoli}, \mbox{VE-2026-11-18}, \mbox{Riscontro Eggon} \\
\leavevmode\hypertarget{req:RV-04}{\textbf{RV-04}} & Il sistema non deve implementare un sistema generale di workflow o plugin, un marketplace o un editor di abilità. & Obbligatorio & I & \mbox{Vincoli} \\
\leavevmode\hypertarget{req:RV-05}{\textbf{RV-05}} & Il sistema deve rispettare i processi e le validazioni di NEXUM; la gestione dell’approvazione resta in NEXUM. & Obbligatorio & I & \hyperref[uc:stato_rich]{UC19}, \hyperref[uc:op_esito]{UC24}, \mbox{Vincoli}, \mbox{RB-13} \\
\leavevmode\hypertarget{req:RV-06}{\textbf{RV-06}} & I controlli su dati e operazioni devono essere svolti dai servizi applicativi e dalle API Eggon in base all’utente autenticato e alla sua azienda; il testo di un messaggio non deve poter cambiare identità o autorizzazioni. & Obbligatorio & T & \hyperref[uc:non_aut]{UC18}, \mbox{Vincoli}, \mbox{RB-01}, \mbox{CA-09}, \mbox{Riscontro Eggon} \\
\leavevmode\hypertarget{req:RV-07}{\textbf{RV-07}} & Il contenuto delle fonti e dei messaggi non deve poter modificare istruzioni operative o autorizzazioni. & Obbligatorio & T & \hyperref[uc:regole]{UC11}, \hyperref[uc:non_aut]{UC18}, \mbox{Vincoli} \\
\leavevmode\hypertarget{req:RV-08}{\textbf{RV-08}} & Gli ambienti di sviluppo e prova devono usare solo dati fittizi o autorizzati. & Obbligatorio & I & \mbox{Vincoli}, \mbox{VE-2026-10-26} \\
\leavevmode\hypertarget{req:RV-09}{\textbf{RV-09}} & La traccia e le conversazioni conservate devono essere mantenute secondo le regole concordate con il proponente. & Obbligatorio & I & \hyperref[uc:tr_elenco]{UC48}, \mbox{Vincoli}, \mbox{RB-15}, \mbox{Riscontro Eggon} \\
\leavevmode\hypertarget{req:RV-10}{\textbf{RV-10}} & Le fonti accettate nell’MVP devono essere documenti PDF testuali. & Obbligatorio & T & \hyperref[uc:fn_carica]{UC39}, \hyperref[uc:fn_formato]{UC40}, \mbox{Capitolato}, \mbox{VE-2026-10-26} \\
\leavevmode\hypertarget{req:RV-11}{\textbf{RV-11}} & Il Concierge deve operare in lingua italiana. & Obbligatorio & T & \mbox{VE-2026-10-26} \\
\leavevmode\hypertarget{req:RV-12}{\textbf{RV-12}} & La soluzione deve essere indipendente dal fornitore dei modelli, che deve poter essere sostituito senza modificare il resto del sistema. & Obbligatorio & I & \mbox{VE-2026-10-26} \\
\leavevmode\hypertarget{req:RV-13}{\textbf{RV-13}} & L’integrazione con NEXUM deve seguire il contratto API concordato con Eggon per lettura dei dati, consultazione dello stato, inserimento e cancellazione delle richieste e verifica dell’esito dopo un’interruzione, con chiamate autenticate tramite OAuth. L’applicativo è autonomo e non riproduce il sistema di accesso di NEXUM. & Obbligatorio & I & \hyperref[uc:dati]{UC15}, \hyperref[uc:stato_rich]{UC19}, \hyperref[uc:registra]{UC20}, \hyperref[uc:op_incerto]{UC28}, \hyperref[uc:annulla_inv]{UC29}, \mbox{Vincoli}, \mbox{VE-2026-11-18}, \mbox{Riscontro Eggon} \\
\leavevmode\hypertarget{req:RV-14}{\textbf{RV-14}} & Con servizi simulati il percorso deve restare eseguibile end to end, con registrazione persistente delle richieste e verifica dello stato. & Obbligatorio & D & \hyperref[uc:registra]{UC20}, \hyperref[uc:stato_rich]{UC19}, \mbox{Vincoli}, \mbox{CA-13} \\
\leavevmode\hypertarget{req:RV-15}{\textbf{RV-15}} & Le abilità di lettura e di operazione devono essere esposte tramite il protocollo \glo{MCP}, con parametri d’ingresso e d’uscita definiti nel registro delle abilità. & Desiderabile & I & \mbox{VE-2026-10-26} \\
\leavevmode\hypertarget{req:RV-16}{\textbf{RV-16}} & Il sistema dovrebbe adottare lo stack di Eggon: Ruby on Rails per il back-end, Angular per il front-end e RubyLLM come interfaccia ai modelli. & Desiderabile & I & \mbox{VE-2026-10-26} \\
\leavevmode\hypertarget{req:RV-17}{\textbf{RV-17}} & I modelli dovrebbero essere usati tramite l’infrastruttura AWS fornita da Eggon, entro i limiti di budget impostati. & Desiderabile & I & \mbox{VE-2026-10-26} \\
\leavevmode\hypertarget{req:RV-18}{\textbf{RV-18}} & I dati restituiti dai modelli dovrebbero essere salvati cifrati. & Desiderabile & I & \mbox{VE-2026-10-26} \\
\end{longtable}}

\subsection{Questioni aperte e proposte}\label{sec:aperte}

Alcuni aspetti richiedono una decisione condivisa con il proponente o dati che
il proponente deve ancora fornire. La Tabella~\ref{tab:aperte} li raccoglie,
con i requisiti interessati; la Tabella~\ref{tab:proposte} riporta le soluzioni
che il gruppo propone per i punti che il proponente ha lasciato alla sua
scelta. Le decisioni prese saranno riportate nei requisiti nella versione
successiva del documento.

{\small
\begin{longtable}{@{}p{1.25cm}>{\raggedright\arraybackslash}p{2.2cm}>{\raggedright\arraybackslash}p{2.1cm}>{\raggedright\arraybackslash}p{6.3cm}>{\raggedright\arraybackslash}p{2.4cm}@{}}
\caption{Questioni aperte da concordare con il proponente}\label{tab:aperte}\\
\toprule
\intestazione{Codice} & \intestazione{Tema} & \intestazione{Requisiti} & \intestazione{Questione} & \intestazione{Chi decide} \\
\midrule\endfirsthead
\toprule
\intestazione{Codice} & \intestazione{Tema} & \intestazione{Requisiti} & \intestazione{Questione} & \intestazione{Chi decide} \\
\midrule\endhead
\midrule\multicolumn{5}{r@{}}{\footnotesize\itshape continua nella pagina successiva}\\ \endfoot \bottomrule\endlastfoot
\textbf{Q1} & Accesso e integrazione & \hyperlink{req:RF-01}{RF-01}, \hyperlink{req:RF-02}{RF-02}, \hyperlink{req:RF-96}{RF-96} & Quali utenti di prova verranno mappati (quanti dipendenti e amministratori, di quali aziende) e come otteniamo i token OAuth (flusso, durata, rinnovo)? & Eggon \\
\textbf{Q2} & Cancellazione & \hyperlink{req:RF-54}{RF-54} & Quali richieste si possono cancellare via API: solo future, anche approvate, anche passate? Le passate restano riservate all’amministratore? & Eggon \\
\textbf{Q3} & Lettura dati & \hyperlink{req:RF-35}{RF-35} & Il servizio fornisce il saldo residuo delle ferie? Con quale data o riferimento di aggiornamento? & Eggon \\
\textbf{Q4} & Lettura dati & \hyperlink{req:RF-33}{RF-33}, \hyperlink{req:RF-34}{RF-34} & Quali stati può avere una richiesta e che cosa significano? Le unità di misura sono giorni, ore o entrambe? & Eggon \\
\textbf{Q5} & Stato delle richieste & \hyperlink{req:RF-39}{RF-39}, \hyperlink{req:RF-49}{RF-49} & Nei dati di prova l’approvazione è automatica o manuale? Lo stato si interroga a richiesta o è previsto un avviso? & Eggon \\
\textbf{Q6} & Tracciabilità & \hyperlink{req:RF-99}{RF-99}, \hyperlink{req:RV-09}{RV-09} & Per quanto tempo e con quale dettaglio vanno conservate traccia e conversazioni, e chi può consultarle? La conservazione delle conversazioni è obbligatoria nell’MVP? & Eggon \\
\textbf{Q7} & Soglie & \hyperlink{req:RQ-01}{RQ-01}, \hyperlink{req:RQ-07}{RQ-07}, \hyperlink{req:RP-02}{RP-02}, \hyperlink{req:RP-03}{RP-03}, \hyperlink{req:RP-04}{RP-04} & Soglie di correttezza, tempi di risposta, tempo di effetto della sostituzione o del ritiro di una fonte, budget per i servizi AI. & Eggon e gruppo \\
\textbf{Q8} & Materiali & \hyperlink{req:RF-28}{RF-28}, \hyperlink{req:RQ-07}{RQ-07}, \hyperlink{req:RV-08}{RV-08} & Documenti aziendali di esempio, manuali NEXUM e dati di prova da usare: quando sono disponibili? & Eggon \\
\textbf{Q9} & Richieste & \hyperlink{req:RF-44}{RF-44}, \hyperlink{req:RF-46}{RF-46} & Tipi di richiesta e campi obbligatori del servizio di ferie e permessi. & Eggon, nella proposta API \\
\textbf{Q10} & Perimetro PoC & — & Quali capacità e requisiti includere nel \glo{PoC} per restare nel monte ore del gruppo. & Gruppo (proposta) e Eggon \\
\textbf{Q11} & Limiti tecnici & \hyperlink{req:RF-19}{RF-19}, \hyperlink{req:RF-77}{RF-77} & Lunghezza massima dei messaggi e dimensione massima dei file: il gruppo propone i valori, Eggon conferma. & Gruppo \\
\end{longtable}}

{\small
\begin{longtable}{@{}p{1.25cm}>{\raggedright\arraybackslash}p{2.4cm}>{\raggedright\arraybackslash}p{2.1cm}>{\raggedright\arraybackslash}p{8.5cm}@{}}
\caption{Proposte del gruppo}\label{tab:proposte}\\
\toprule
\intestazione{Codice} & \intestazione{Tema} & \intestazione{Requisiti} & \intestazione{Proposta} \\
\midrule\endfirsthead
\toprule
\intestazione{Codice} & \intestazione{Tema} & \intestazione{Requisiti} & \intestazione{Proposta} \\
\midrule\endhead
\midrule\multicolumn{4}{r@{}}{\footnotesize\itshape continua nella pagina successiva}\\ \endfoot \bottomrule\endlastfoot
\textbf{P1} & Versionamento delle fonti & \hyperlink{req:RF-84}{RF-84}, \hyperlink{req:RF-85}{RF-85}, \hyperlink{req:RF-15}{RF-15} & Ogni caricamento di una nuova versione crea un nuovo record della fonte con numero di versione progressivo; la versione precedente viene marcata «sostituita», resta nello storico ed è esclusa dal recupero. Ogni porzione ricercabile porta l’identificativo di fonte e versione, e la traccia di ogni interazione registra quelli usati: così si ricostruisce quale versione ha supportato una risposta passata. \\
\textbf{P2} & Riconoscimento degli utenti & \hyperlink{req:RF-01}{RF-01}–\hyperlink{req:RF-05}{RF-05} & Accesso proprio dell’applicativo con identificativo e password per gli utenti di prova, ciascuno con ruolo (dipendente o amministratore) e associato all’utente NEXUM della mappatura. Le chiamate alle API Eggon usano OAuth (RF-96). \\
\textbf{P3} & Categoria dei bucket & \hyperlink{req:RF-60}{RF-60} & La categoria (conoscenza aziendale o conoscenza del prodotto) è scelta dall’amministratore alla creazione del bucket ed è mostrata in elenco e nelle citazioni. \\
\textbf{P4} & Proposta API & \hyperlink{req:RV-13}{RV-13}, \hyperlink{req:RF-53}{RF-53} & Proposta in \glo{Bruno} con lettura dei dati, consultazione dello stato, inserimento, cancellazione e verifica dell’esito dopo un’interruzione: per ogni operazione dati in ingresso, risposta, errori ed esempio. Documento separato, da presentare a Eggon. \\
\end{longtable}}

\clearpage
\section{Tracciamento}\label{sec:tracciamento}

Il tracciamento verifica due proprietà: che ogni fonte del capitolato sia
coperta da almeno un requisito, e che ogni caso d'uso abbia prodotto almeno un
requisito funzionale. Il percorso inverso, da ogni requisito alle sue fonti, è
nella colonna \textit{Fonti} delle tabelle della sezione~\ref{sec:requisiti}.

\subsection{Dalle fonti ai requisiti}

La Tabella~\ref{tab:tr-fonti} riporta per ogni requisito di business e per ogni
criterio di accettazione del capitolato i requisiti che lo soddisfano. Nessuna
riga è vuota: ogni fonte è coperta.

{\small
\begin{longtable}{@{}p{2.6cm}>{\raggedright\arraybackslash}p{7.0cm}>{\raggedright\arraybackslash}p{5.2cm}@{}}
\caption{Tracciamento dalle fonti ai requisiti}\label{tab:tr-fonti}\\
\toprule \intestazione{Fonte} & \intestazione{Contenuto} & \intestazione{Requisiti} \\ \midrule\endfirsthead
\toprule \intestazione{Fonte} & \intestazione{Contenuto} & \intestazione{Requisiti} \\ \midrule\endhead
\midrule\multicolumn{3}{r@{}}{\footnotesize\itshape continua nella pagina successiva}\\ \endfoot \bottomrule\endlastfoot
\multicolumn{3}{@{}l}{\textbf{Requisiti di business del capitolato}} \\ \addlinespace[2pt]
\mbox{RB-01} & Il Concierge opera nel contesto della persona autenticata e dell'azienda selezionata. & \hyperlink{req:RF-01}{RF-01}, \hyperlink{req:RF-02}{RF-02}, \hyperlink{req:RF-07}{RF-07}, \hyperlink{req:RF-38}{RF-38}, \hyperlink{req:RV-06}{RV-06} \\
\mbox{RB-02} & Distingue richieste su regole aziendali, istruzioni di prodotto, dati personali e operazioni. & \hyperlink{req:RF-13}{RF-13}, \hyperlink{req:RQ-03}{RQ-03} \\
\mbox{RB-03} & Combina le capacità necessarie nella stessa conversazione. & \hyperlink{req:RF-14}{RF-14} \\
\mbox{RB-04} & Mantiene il contesto del dialogo e chiede chiarimenti quando mancano informazioni. & \hyperlink{req:RF-09}{RF-09}, \hyperlink{req:RF-12}{RF-12}, \hyperlink{req:RF-16}{RF-16} \\
\mbox{RB-05} & Il backoffice consente caricamento, organizzazione, sostituzione e ritiro delle fonti aziendali. & \hyperlink{req:RF-57}{RF-57}, \hyperlink{req:RF-59}{RF-59}, \hyperlink{req:RF-66}{RF-66}, \hyperlink{req:RF-67}{RF-67}, \hyperlink{req:RF-68}{RF-68}, \hyperlink{req:RF-70}{RF-70}, \hyperlink{req:RF-84}{RF-84}, \hyperlink{req:RF-86}{RF-86} \\
\mbox{RB-06} & Le fonti diventano conoscenza consultabile mantenendo provenienza e versione. & \hyperlink{req:RF-22}{RF-22}, \hyperlink{req:RF-72}{RF-72}, \hyperlink{req:RF-74}{RF-74}, \hyperlink{req:RF-80}{RF-80}, \hyperlink{req:RF-84}{RF-84} \\
\mbox{RB-07} & I manuali NEXUM sono una fonte riconoscibile e distinta dai regolamenti aziendali. & \hyperlink{req:RF-24}{RF-24}, \hyperlink{req:RF-29}{RF-29}, \hyperlink{req:RF-60}{RF-60} \\
\mbox{RB-08} & Le risposte documentali indicano le fonti; con fonti insufficienti o in contrasto il sistema si astiene. & \hyperlink{req:RF-23}{RF-23}, \hyperlink{req:RF-24}{RF-24}, \hyperlink{req:RF-25}{RF-25}, \hyperlink{req:RQ-01}{RQ-01} \\
\mbox{RB-09} & I dati personali sono recuperati tramite servizi applicativi autorizzati. & \hyperlink{req:RF-31}{RF-31} \\
\mbox{RB-10} & I risultati sui dati specificano periodo, unità di misura e significato degli stati. & \hyperlink{req:RF-33}{RF-33}, \hyperlink{req:RF-34}{RF-34} \\
\mbox{RB-11} & Le operazioni sono eseguite tramite abilità definite e servizi autorizzati. & \hyperlink{req:RF-42}{RF-42} \\
\mbox{RB-12} & Prima dell'invio, riepilogo completo e conferma esplicita; una modifica richiede nuovo riepilogo e nuova conferma. & \hyperlink{req:RF-41}{RF-41}, \hyperlink{req:RF-47}{RF-47}, \hyperlink{req:RF-48}{RF-48}, \hyperlink{req:RF-50}{RF-50}, \hyperlink{req:RF-51}{RF-51} \\
\mbox{RB-13} & Il sistema distingue registrazione, attesa di approvazione e approvazione. & \hyperlink{req:RF-39}{RF-39}, \hyperlink{req:RF-49}{RF-49}, \hyperlink{req:RV-05}{RV-05} \\
\mbox{RB-14} & Errori tecnici e insufficienza informativa sono comunicati chiaramente. & \hyperlink{req:RF-05}{RF-05}, \hyperlink{req:RF-18}{RF-18}, \hyperlink{req:RF-20}{RF-20}, \hyperlink{req:RF-37}{RF-37}, \hyperlink{req:RF-52}{RF-52}, \hyperlink{req:RF-83}{RF-83} \\
\mbox{RB-15} & Fonti consultate ed esiti delle operazioni sono tracciabili secondo le autorizzazioni. & \hyperlink{req:RF-15}{RF-15}, \hyperlink{req:RF-43}{RF-43}, \hyperlink{req:RF-89}{RF-89}, \hyperlink{req:RF-90}{RF-90}, \hyperlink{req:RF-95}{RF-95}, \hyperlink{req:RV-09}{RV-09} \\
\mbox{RB-16} & Si possono aggiungere fonti e abilità senza riprogettare il Concierge. & \hyperlink{req:RQ-10}{RQ-10} \\
\addlinespace
\multicolumn{3}{@{}l}{\textbf{Criteri di accettazione del capitolato}} \\ \addlinespace[2pt]
\mbox{CA-01} & Risposta a una domanda sul regolamento aziendale usando una fonte caricata dal backoffice. & \hyperlink{req:RF-21}{RF-21}, \hyperlink{req:RF-70}{RF-70} \\
\mbox{CA-02} & Spiegazione dell'uso di una funzione sulla base dei manuali NEXUM. & \hyperlink{req:RF-28}{RF-28} \\
\mbox{CA-03} & Recupero e presentazione corretta dei dati personali su ferie o permessi. & \hyperlink{req:RF-31}{RF-31}, \hyperlink{req:RF-32}{RF-32}, \hyperlink{req:RF-33}{RF-33} \\
\mbox{CA-04} & Preparazione e registrazione di una richiesta dopo conferma. & \hyperlink{req:RF-40}{RF-40}, \hyperlink{req:RF-47}{RF-47}, \hyperlink{req:RF-48}{RF-48} \\
\mbox{CA-05} & Continuità tra domande successive. & \hyperlink{req:RF-12}{RF-12}, \hyperlink{req:RQ-04}{RQ-04} \\
\mbox{CA-06} & Chiarimenti su informazioni ambigue, comprese le date prive di un anno ricavabile dal contesto. & \hyperlink{req:RF-16}{RF-16}, \hyperlink{req:RF-17}{RF-17}, \hyperlink{req:RF-45}{RF-45}, \hyperlink{req:RQ-04}{RQ-04} \\
\mbox{CA-07} & Visibilità delle fonti e astensione con informazioni insufficienti o in contrasto. & \hyperlink{req:RF-24}{RF-24}, \hyperlink{req:RF-25}{RF-25} \\
\mbox{CA-08} & Aggiornamento delle risposte dopo sostituzione o ritiro di un documento. & \hyperlink{req:RF-84}{RF-84}, \hyperlink{req:RF-85}{RF-85}, \hyperlink{req:RF-86}{RF-86}, \hyperlink{req:RP-03}{RP-03} \\
\mbox{CA-09} & Rispetto delle autorizzazioni. & \hyperlink{req:RF-02}{RF-02}, \hyperlink{req:RF-07}{RF-07}, \hyperlink{req:RF-22}{RF-22}, \hyperlink{req:RF-38}{RF-38}, \hyperlink{req:RF-63}{RF-63}, \hyperlink{req:RF-101}{RF-101}, \hyperlink{req:RV-06}{RV-06} \\
\mbox{CA-10} & Distinzione tra richiesta registrata e approvata; nuova conferma se cambiano i dati. & \hyperlink{req:RF-39}{RF-39}, \hyperlink{req:RF-49}{RF-49}, \hyperlink{req:RF-50}{RF-50} \\
\mbox{CA-11} & Gestione di errori e invii ripetuti. & \hyperlink{req:RF-20}{RF-20}, \hyperlink{req:RF-52}{RF-52}, \hyperlink{req:RF-53}{RF-53}, \hyperlink{req:RQ-06}{RQ-06} \\
\mbox{CA-12} & Disponibilità della traccia delle operazioni. & \hyperlink{req:RF-15}{RF-15}, \hyperlink{req:RF-43}{RF-43}, \hyperlink{req:RF-89}{RF-89}, \hyperlink{req:RF-95}{RF-95} \\
\mbox{CA-13} & Percorso unitario dalla domanda informativa alla consultazione dei dati e all'invio della richiesta. & \hyperlink{req:RF-14}{RF-14}, \hyperlink{req:RV-14}{RV-14} \\
\addlinespace
\multicolumn{3}{@{}l}{\textbf{Altre fonti}} \\ \addlinespace[2pt]
\mbox{Capitolato} & Capitolato C2, sezione «Ambito funzionale». & \hyperlink{req:RF-11}{RF-11}, \hyperlink{req:RF-32}{RF-32}, \hyperlink{req:RF-36}{RF-36}, \hyperlink{req:RF-44}{RF-44}, \hyperlink{req:RF-45}{RF-45}, \hyperlink{req:RF-46}{RF-46}, \hyperlink{req:RF-53}{RF-53}, \hyperlink{req:RF-58}{RF-58}, \hyperlink{req:RF-59}{RF-59}, \hyperlink{req:RF-69}{RF-69}, \hyperlink{req:RF-75}{RF-75}, \hyperlink{req:RF-81}{RF-81}, \hyperlink{req:RF-83}{RF-83}, \hyperlink{req:RF-85}{RF-85}, \hyperlink{req:RF-101}{RF-101}, \hyperlink{req:RV-10}{RV-10} \\
\mbox{Vincoli} & Capitolato C2, sezione «Vincoli». & \hyperlink{req:RF-18}{RF-18}, \hyperlink{req:RF-48}{RF-48}, \hyperlink{req:RF-52}{RF-52}, \hyperlink{req:RF-76}{RF-76}, \hyperlink{req:RQ-11}{RQ-11}, \hyperlink{req:RV-01}{RV-01}, \hyperlink{req:RV-02}{RV-02}, \hyperlink{req:RV-03}{RV-03}, \hyperlink{req:RV-04}{RV-04}, \hyperlink{req:RV-05}{RV-05}, \hyperlink{req:RV-06}{RV-06}, \hyperlink{req:RV-07}{RV-07}, \hyperlink{req:RV-08}{RV-08}, \hyperlink{req:RV-09}{RV-09}, \hyperlink{req:RV-13}{RV-13}, \hyperlink{req:RV-14}{RV-14} \\
\mbox{Prestazioni} & Capitolato C2, sezione «Requisiti di prestazione». & \hyperlink{req:RF-71}{RF-71}, \hyperlink{req:RQ-01}{RQ-01}, \hyperlink{req:RQ-02}{RQ-02}, \hyperlink{req:RQ-03}{RQ-03}, \hyperlink{req:RQ-04}{RQ-04}, \hyperlink{req:RQ-05}{RQ-05}, \hyperlink{req:RQ-06}{RQ-06}, \hyperlink{req:RQ-07}{RQ-07}, \hyperlink{req:RP-01}{RP-01}, \hyperlink{req:RP-02}{RP-02}, \hyperlink{req:RP-03}{RP-03}, \hyperlink{req:RP-04}{RP-04} \\
\mbox{Deliverables} & Capitolato C2, sezione «Risultati attesi». & \hyperlink{req:RQ-08}{RQ-08}, \hyperlink{req:RQ-09}{RQ-09}, \hyperlink{req:RQ-12}{RQ-12}, \hyperlink{req:RQ-13}{RQ-13} \\
\mbox{VE-2026-10-26} & Verbale esterno dell'incontro con Eggon del 2026-10-26. & \hyperlink{req:RF-35}{RF-35}, \hyperlink{req:RF-54}{RF-54}, \hyperlink{req:RF-55}{RF-55}, \hyperlink{req:RQ-15}{RQ-15}, \hyperlink{req:RP-04}{RP-04}, \hyperlink{req:RV-01}{RV-01}, \hyperlink{req:RV-08}{RV-08}, \hyperlink{req:RV-10}{RV-10}, \hyperlink{req:RV-11}{RV-11}, \hyperlink{req:RV-12}{RV-12}, \hyperlink{req:RV-15}{RV-15}, \hyperlink{req:RV-16}{RV-16}, \hyperlink{req:RV-17}{RV-17}, \hyperlink{req:RV-18}{RV-18} \\
\mbox{Interno} & Proposta del gruppo, a tutela dell'usabilità o della robustezza. & \hyperlink{req:RF-03}{RF-03}, \hyperlink{req:RF-04}{RF-04}, \hyperlink{req:RF-05}{RF-05}, \hyperlink{req:RF-08}{RF-08}, \hyperlink{req:RF-10}{RF-10}, \hyperlink{req:RF-19}{RF-19}, \hyperlink{req:RF-60}{RF-60}, \hyperlink{req:RF-64}{RF-64}, \hyperlink{req:RF-65}{RF-65}, \hyperlink{req:RF-73}{RF-73}, \hyperlink{req:RF-77}{RF-77}, \hyperlink{req:RF-78}{RF-78}, \hyperlink{req:RF-87}{RF-87}, \hyperlink{req:RF-91}{RF-91}, \hyperlink{req:RF-92}{RF-92}, \hyperlink{req:RF-93}{RF-93}, \hyperlink{req:RF-94}{RF-94}, \hyperlink{req:RQ-16}{RQ-16} \\
\mbox{VE-2026-11-18} & Verbale esterno dell'incontro con Eggon del 2026-11-18. & \hyperlink{req:RF-01}{RF-01}, \hyperlink{req:RF-02}{RF-02}, \hyperlink{req:RF-09}{RF-09}, \hyperlink{req:RF-18}{RF-18}, \hyperlink{req:RF-20}{RF-20}, \hyperlink{req:RF-22}{RF-22}, \hyperlink{req:RF-39}{RF-39}, \hyperlink{req:RF-54}{RF-54}, \hyperlink{req:RF-57}{RF-57}, \hyperlink{req:RF-58}{RF-58}, \hyperlink{req:RF-59}{RF-59}, \hyperlink{req:RF-63}{RF-63}, \hyperlink{req:RF-65}{RF-65}, \hyperlink{req:RF-66}{RF-66}, \hyperlink{req:RF-83}{RF-83}, \hyperlink{req:RF-84}{RF-84}, \hyperlink{req:RF-85}{RF-85}, \hyperlink{req:RF-90}{RF-90}, \hyperlink{req:RF-95}{RF-95}, \hyperlink{req:RF-96}{RF-96}, \hyperlink{req:RF-97}{RF-97}, \hyperlink{req:RF-98}{RF-98}, \hyperlink{req:RF-99}{RF-99}, \hyperlink{req:RF-100}{RF-100}, \hyperlink{req:RF-101}{RF-101}, \hyperlink{req:RV-03}{RV-03}, \hyperlink{req:RV-13}{RV-13} \\
\mbox{Riscontro Eggon} & Osservazioni scritte di Eggon sulla bozza dei requisiti (novembre 2026). & \hyperlink{req:RF-01}{RF-01}, \hyperlink{req:RF-02}{RF-02}, \hyperlink{req:RF-09}{RF-09}, \hyperlink{req:RF-10}{RF-10}, \hyperlink{req:RF-15}{RF-15}, \hyperlink{req:RF-18}{RF-18}, \hyperlink{req:RF-24}{RF-24}, \hyperlink{req:RF-33}{RF-33}, \hyperlink{req:RF-35}{RF-35}, \hyperlink{req:RF-36}{RF-36}, \hyperlink{req:RF-39}{RF-39}, \hyperlink{req:RF-43}{RF-43}, \hyperlink{req:RF-54}{RF-54}, \hyperlink{req:RF-55}{RF-55}, \hyperlink{req:RF-57}{RF-57}, \hyperlink{req:RF-58}{RF-58}, \hyperlink{req:RF-59}{RF-59}, \hyperlink{req:RF-60}{RF-60}, \hyperlink{req:RF-63}{RF-63}, \hyperlink{req:RF-65}{RF-65}, \hyperlink{req:RF-66}{RF-66}, \hyperlink{req:RF-84}{RF-84}, \hyperlink{req:RF-85}{RF-85}, \hyperlink{req:RF-89}{RF-89}, \hyperlink{req:RF-90}{RF-90}, \hyperlink{req:RF-95}{RF-95}, \hyperlink{req:RF-96}{RF-96}, \hyperlink{req:RF-97}{RF-97}, \hyperlink{req:RF-99}{RF-99}, \hyperlink{req:RF-100}{RF-100}, \hyperlink{req:RF-101}{RF-101}, \hyperlink{req:RF-102}{RF-102}, \hyperlink{req:RQ-10}{RQ-10}, \hyperlink{req:RP-03}{RP-03}, \hyperlink{req:RV-03}{RV-03}, \hyperlink{req:RV-06}{RV-06}, \hyperlink{req:RV-09}{RV-09}, \hyperlink{req:RV-13}{RV-13} \\
\addlinespace
\end{longtable}}

\subsection{Dai casi d'uso ai requisiti}

La Tabella~\ref{tab:tr-uc} riporta per ogni caso d'uso i requisiti funzionali
che ne derivano.

{\small
\begin{longtable}{@{}p{1.4cm}>{\raggedright\arraybackslash}p{9.2cm}>{\raggedright\arraybackslash}p{4.2cm}@{}}
\caption{Tracciamento dai casi d'uso ai requisiti funzionali}\label{tab:tr-uc}\\
\toprule \intestazione{Codice} & \intestazione{Caso d'uso} & \intestazione{Requisiti} \\ \midrule\endfirsthead
\toprule \intestazione{Codice} & \intestazione{Caso d'uso} & \intestazione{Requisiti} \\ \midrule\endhead
\midrule\multicolumn{3}{r@{}}{\footnotesize\itshape continua nella pagina successiva}\\ \endfoot \bottomrule\endlastfoot
\hyperref[uc:login]{UC1} & Autenticazione & \hyperlink{req:RF-01}{RF-01}, \hyperlink{req:RF-02}{RF-02}, \hyperlink{req:RF-07}{RF-07} \\
\hyperref[uc:login_id]{UC1.1} & Inserimento dell'identificativo & \hyperlink{req:RF-03}{RF-03} \\
\hyperref[uc:login_pw]{UC1.2} & Inserimento della password & \hyperlink{req:RF-04}{RF-04} \\
\hyperref[uc:cred_err]{UC2} & Credenziali non valide & \hyperlink{req:RF-05}{RF-05} \\
\hyperref[uc:logout]{UC3} & Uscita dal sistema & \hyperlink{req:RF-08}{RF-08} \\
\hyperref[uc:nuova_conv]{UC4} & Avvio di una nuova conversazione & \hyperlink{req:RF-09}{RF-09}, \hyperlink{req:RF-98}{RF-98} \\
\hyperref[uc:vis_conv]{UC5} & Visualizzazione della conversazione & \hyperlink{req:RF-10}{RF-10}, \hyperlink{req:RF-97}{RF-97} \\
\hyperref[uc:msg]{UC6} & Invio di un messaggio al Concierge & \hyperlink{req:RF-11}{RF-11}, \hyperlink{req:RF-12}{RF-12}, \hyperlink{req:RF-13}{RF-13}, \hyperlink{req:RF-14}{RF-14}, \hyperlink{req:RF-15}{RF-15}, \hyperlink{req:RF-99}{RF-99} \\
\hyperref[uc:chiarimento]{UC7} & Richiesta di chiarimento & \hyperlink{req:RF-16}{RF-16}, \hyperlink{req:RF-17}{RF-17} \\
\hyperref[uc:fuori_ambito]{UC8} & Richiesta fuori dal perimetro & \hyperlink{req:RF-18}{RF-18} \\
\hyperref[uc:msg_lungo]{UC9} & Messaggio troppo lungo & \hyperlink{req:RF-19}{RF-19} \\
\hyperref[uc:errore_tecnico]{UC10} & Comunicazione di un errore tecnico & \hyperlink{req:RF-20}{RF-20} \\
\hyperref[uc:regole]{UC11} & Domanda sulle regole aziendali & \hyperlink{req:RF-21}{RF-21}, \hyperlink{req:RF-22}{RF-22}, \hyperlink{req:RF-23}{RF-23} \\
\hyperref[uc:fonti]{UC12} & Visualizzazione delle fonti citate & \hyperlink{req:RF-24}{RF-24}, \hyperlink{req:RF-100}{RF-100} \\
\hyperref[uc:astensione]{UC13} & Astensione dalla risposta & \hyperlink{req:RF-25}{RF-25} \\
\hyperref[uc:prodotto]{UC14} & Domanda sull'utilizzo di NEXUM & \hyperlink{req:RF-28}{RF-28}, \hyperlink{req:RF-29}{RF-29} \\
\hyperref[uc:dati]{UC15} & Consultazione dei propri dati di ferie e permessi & \hyperlink{req:RF-31}{RF-31}, \hyperlink{req:RF-32}{RF-32}, \hyperlink{req:RF-96}{RF-96} \\
\hyperref[uc:dati_riep]{UC15.1} & Visualizzazione delle richieste per periodo e stato & \hyperlink{req:RF-33}{RF-33} \\
\hyperref[uc:dati_criteri]{UC15.2} & Visualizzazione di periodo, unità di misura e criterio di conteggio & \hyperlink{req:RF-34}{RF-34} \\
\hyperref[uc:saldo]{UC16} & Visualizzazione del saldo residuo & \hyperlink{req:RF-35}{RF-35}, \hyperlink{req:RF-36}{RF-36} \\
\hyperref[uc:nessun_dato]{UC17} & Nessun dato corrispondente & \hyperlink{req:RF-37}{RF-37} \\
\hyperref[uc:non_aut]{UC18} & Rifiuto di una richiesta non autorizzata & \hyperlink{req:RF-38}{RF-38} \\
\hyperref[uc:stato_rich]{UC19} & Consultazione dello stato di una richiesta & \hyperlink{req:RF-39}{RF-39}, \hyperlink{req:RF-96}{RF-96} \\
\hyperref[uc:registra]{UC20} & Registrazione di una richiesta di ferie o permesso & \hyperlink{req:RF-40}{RF-40}, \hyperlink{req:RF-41}{RF-41}, \hyperlink{req:RF-42}{RF-42}, \hyperlink{req:RF-43}{RF-43}, \hyperlink{req:RF-96}{RF-96} \\
\hyperref[uc:reg_tipo]{UC20.1} & Determinazione del tipo di richiesta & \hyperlink{req:RF-44}{RF-44} \\
\hyperref[uc:periodo]{UC21} & Determinazione del periodo & \hyperlink{req:RF-45}{RF-45} \\
\hyperref[uc:reg_campi]{UC20.2} & Raccolta degli altri dati richiesti & \hyperlink{req:RF-46}{RF-46} \\
\hyperref[uc:op_riep]{UC22} & Visualizzazione del riepilogo dell'operazione & \hyperlink{req:RF-47}{RF-47} \\
\hyperref[uc:op_conf]{UC23} & Conferma esplicita dell'operazione & \hyperlink{req:RF-48}{RF-48} \\
\hyperref[uc:op_esito]{UC24} & Visualizzazione dell'esito dell'operazione & \hyperlink{req:RF-49}{RF-49} \\
\hyperref[uc:op_modifica]{UC25} & Modifica dei dati prima della conferma & \hyperlink{req:RF-50}{RF-50} \\
\hyperref[uc:op_rinuncia]{UC26} & Rinuncia all'operazione prima della conferma & \hyperlink{req:RF-51}{RF-51} \\
\hyperref[uc:op_rifiuto]{UC27} & Operazione respinta dal servizio & \hyperlink{req:RF-52}{RF-52} \\
\hyperref[uc:op_incerto]{UC28} & Esito dell'operazione incerto & \hyperlink{req:RF-53}{RF-53} \\
\hyperref[uc:annulla_inv]{UC29} & Cancellazione di una richiesta & \hyperlink{req:RF-54}{RF-54}, \hyperlink{req:RF-96}{RF-96}, \hyperlink{req:RF-102}{RF-102} \\
\hyperref[uc:sel_rich]{UC30} & Individuazione della richiesta interessata & \hyperlink{req:RF-55}{RF-55}, \hyperlink{req:RF-102}{RF-102} \\
\hyperref[uc:bk_elenco]{UC31} & Visualizzazione dell'elenco dei bucket & \hyperlink{req:RF-57}{RF-57} \\
\hyperref[uc:bk_elem]{UC31.1} & Visualizzazione di un bucket nell'elenco & \hyperlink{req:RF-58}{RF-58} \\
\hyperref[uc:bk_crea]{UC32} & Creazione di un bucket & \hyperlink{req:RF-59}{RF-59}, \hyperlink{req:RF-60}{RF-60}, \hyperlink{req:RF-63}{RF-63} \\
\hyperref[uc:bk_nome]{UC32.1} & Inserimento del nome del bucket & \hyperlink{req:RF-59}{RF-59} \\
\hyperref[uc:bk_desc]{UC32.2} & Inserimento della descrizione del bucket & \hyperlink{req:RF-59}{RF-59} \\
\hyperref[uc:bk_dest]{UC33} & Limitazione della consultazione a gruppi & \hyperlink{req:RF-101}{RF-101} \\
\hyperref[uc:bk_nome_dup]{UC34} & Nome del bucket già in uso & \hyperlink{req:RF-64}{RF-64} \\
\hyperref[uc:bk_incompleto]{UC35} & Dati del bucket mancanti & \hyperlink{req:RF-65}{RF-65} \\
\hyperref[uc:bk_modifica]{UC36} & Modifica di un bucket & \hyperlink{req:RF-66}{RF-66}, \hyperlink{req:RF-101}{RF-101} \\
\hyperref[uc:bk_elimina]{UC37} & Eliminazione di un bucket & \hyperlink{req:RF-67}{RF-67} \\
\hyperref[uc:fn_elenco]{UC38} & Visualizzazione delle fonti di un bucket & \hyperlink{req:RF-68}{RF-68} \\
\hyperref[uc:fn_elem]{UC38.1} & Visualizzazione di una fonte nell'elenco & \hyperlink{req:RF-69}{RF-69} \\
\hyperref[uc:fn_carica]{UC39} & Caricamento di una fonte & \hyperlink{req:RF-70}{RF-70}, \hyperlink{req:RF-71}{RF-71}, \hyperlink{req:RF-72}{RF-72} \\
\hyperref[uc:fn_file]{UC39.1} & Selezione del file & \hyperlink{req:RF-73}{RF-73} \\
\hyperref[uc:fn_meta]{UC39.2} & Inserimento dei metadati della fonte & \hyperlink{req:RF-74}{RF-74}, \hyperlink{req:RF-75}{RF-75} \\
\hyperref[uc:fn_formato]{UC40} & Formato del file non supportato & \hyperlink{req:RF-76}{RF-76} \\
\hyperref[uc:fn_dim]{UC41} & Dimensione massima superata & \hyperlink{req:RF-77}{RF-77} \\
\hyperref[uc:fn_dup]{UC42} & Fonte già presente & \hyperlink{req:RF-78}{RF-78} \\
\hyperref[uc:fn_dettaglio]{UC43} & Visualizzazione del dettaglio di una fonte &  \\
\hyperref[uc:fn_meta_v]{UC43.1} & Visualizzazione dei metadati & \hyperlink{req:RF-80}{RF-80} \\
\hyperref[uc:fn_stato]{UC43.2} & Visualizzazione dello stato di elaborazione & \hyperlink{req:RF-81}{RF-81} \\
\hyperref[uc:fn_errore]{UC44} & Visualizzazione dell'errore di elaborazione & \hyperlink{req:RF-83}{RF-83} \\
\hyperref[uc:fn_sostituisci]{UC45} & Sostituzione di una fonte & \hyperlink{req:RF-84}{RF-84}, \hyperlink{req:RF-85}{RF-85} \\
\hyperref[uc:fn_ritira]{UC46} & Ritiro di una fonte & \hyperlink{req:RF-86}{RF-86} \\
\hyperref[uc:fn_riprova]{UC47} & Nuova elaborazione di una fonte & \hyperlink{req:RF-87}{RF-87} \\
\hyperref[uc:tr_elenco]{UC48} & Consultazione del registro delle interazioni & \hyperlink{req:RF-89}{RF-89} \\
\hyperref[uc:tr_voce]{UC48.1} & Visualizzazione di una voce del registro & \hyperlink{req:RF-90}{RF-90} \\
\hyperref[uc:tr_filtro]{UC49} & Filtraggio del registro & \hyperlink{req:RF-91}{RF-91} \\
\hyperref[uc:tr_f_persona]{UC49.1} & Filtraggio per persona & \hyperlink{req:RF-92}{RF-92} \\
\hyperref[uc:tr_f_periodo]{UC49.2} & Filtraggio per periodo & \hyperlink{req:RF-93}{RF-93} \\
\hyperref[uc:tr_f_capacita]{UC49.3} & Filtraggio per capacità ed esito & \hyperlink{req:RF-94}{RF-94} \\
\hyperref[uc:tr_dettaglio]{UC50} & Visualizzazione del dettaglio di un'interazione & \hyperlink{req:RF-95}{RF-95}, \hyperlink{req:RF-99}{RF-99} \\
\end{longtable}}

\subsection{Riepilogo}
Il documento individua 68 casi d'uso, di cui 50 principali e 18 sottocasi, e 129 requisiti, riepilogati per tipo e priorità nella Tabella~\ref{tab:riepilogo}.
\begin{table}[H]
\centering
\small
\caption{Riepilogo dei requisiti per tipo e priorità}\label{tab:riepilogo}
\begin{tabular}{@{}lrrrr@{}}
\toprule
\intestazione{Tipo} & \intestazione{Obbligatori} & \intestazione{Desiderabili} & \intestazione{Facoltativi} & \intestazione{Totale} \\
\midrule
Funzionali & 73 & 15 & 4 & 92 \\
Di qualità & 14 & 1 & 0 & 15 \\
Prestazionali & 4 & 0 & 0 & 4 \\
Di vincolo & 14 & 4 & 0 & 18 \\
\midrule
\textbf{Totale} & 105 & 20 & 4 & \textbf{129} \\
\bottomrule
\end{tabular}
\end{table}

\end{document}
